% Gerado por tools/tei2latex.py a partir de revisado.xml. Regenerar, não editar.

%%>>> parte-1

%%>>> art-23

%%>>> t167

\pg{279} \tese{167}{A percepção sensitiva interna é mais elevada que a externa.}
Uma operação é mais elevada que outra, quando a faculdade, de que deriva a primeira, é mais elevada que a faculdade, de que deriva a segunda. Ora os sentidos \emph{internos} são mais elevados que os \emph{externos}; porque, embora todos os sentidos sejam compostos de um elemento \emph{material}, que é o \emph{orgão}, e de um elemento \emph{simples}, que é a \emph{energia sensitiva}, comtudo nos sentidos \emph{externos} predomina o elemento \emph{material}, ao passo que nos sentidos \emph{internos} predomina o elemento \emph{simples}. Logo a percepção sensitiva \emph{interna} é mais elevada que a \emph{externa} \footnote{Que o modo de operar dos sentidos \emph{internos} seja \emph{menos material} que o dos \emph{externos}, vê-se d'isto que os sentidos \emph{externos} não produzem o acto e não o proseguem, senão no tempo em que o objecto está presente --- ou \emph{por si mesmo}, como no tacto e no gosto, --- ou pela \emph{sua acção}, como nos outros sentidos; ao passo que os sentidos \emph{internos} operam, continuam ou tornam a operar, embora já não esteja presente o objecto externo. --- Sobretudo admiravel é a \emph{imaginação}, que, desprendendo-se da acção dos objectos, forma especies de coisas nunca vistas (embora com elementos anteriormente percebidos), e d'este modo imita a intelligencia, a qual, depois de ter recebido dos sentidos as especies dos obiectos externos, conserva essas especies (transformadas) e, para as reproduzir, não precisa do estimulo do objecto externo. --- Mas a imaginação encontra o seu perigo onde tem a sua gloria. Emquanto o sentido \emph{externo} não se engana na percepção do seu objecto (se não houver obstaculo), porque o proprio objecto está presente, a imaginação pode enganar-se, apprehendendo como presente um objecto ausente. Diz S. Thomaz: « Est quaedam vis apprehensiva, quae apprehendit speciem sensibilem, sensibili re praesente, sicut \emph{sensus proprius;} quaedam vero apprehendit eam, re absente, sicut \emph{imaginatio}. Et ideo \emph{sensus semper apprehendit rem, ut est,} nisi sit impedimentum in organo vel in medio; sed \emph{imaginatio apprehendit ut plurimum rem, uti non est,} quia appreheudit eam ut praesentem cum sit absens, et ideo dicit Philosophus (in IV Metaph.) quod sensus non est dominus falsitatis, sed phantasia » (\emph{De Verit.}, q. I, a. XI).\par *\par Aqui devemos accrescentar algumas coisas ácerca das operações da \emph{imaginação} durante o \emph{somno}.\par \emph{Somno} é a suspensão periodica das operações sensitivas ou da actividade dos centros nervosos superiores, devida a cançaço cerebral.\par Esse cançaço, de qué deriva o somno, é uma consequencia do trabalho, que o eixo cerebro-espinhal faz durante a vigilia, e que, accumulando neste orgão productos excrementicios, tira-lhe toda a excitabilidade. --- Durante o \emph{somno}, o sangue afflue menos abundantemente ao cerebro, e o systema nervoso repara as proprias perdas (\emph{Errera}, na \emph{Revue Scientifique}, 23 de julho de 1877).\par A suspensão das operações sensitivas, durante o \emph{somno}, póde não ser total, e então temos o \emph{sonho}.\par O \emph{sonho}, pois, é a reproducção parcial e irregular de certas imagens, mais\pgn{280}profundamente gravadas na phantasia, estando o exercicio dos sentidos externos quasi completamente suspenso. --- Por isso no \emph{sonho} ordinariamente se reproduzem as imagens, que fizeram mais impressão ou que mais preoccuparam durante a vigilia, e estas imagens ás vezes causam uma emoção tão forte, que excita e acorda o systema nervoso da vida sensitiva.\par Os phenomenos mais notaveis que podem dar-se no \emph{sonho}, são cinco: as \emph{allucinações}, a \emph{desordem nos juizos e nos raciocinios}, o \emph{somniloquio}, o \emph{somnambulismo}, e a \emph{lucidez natural}.\par 1.º) As \emph{allucinações}, que consistem em perceber o que não é, são devidás á \emph{vivacidade} das \emph{imagens} e á suspensão do exercicio dos sentidos externos; pois que, quando as imagens são tão vivas como as proprias percepções externas, e os sentidos externos são incapazes de comparar estas com aquellas, é natural que se confunda o mundo phantastico com o real.\par 2.º) A \emph{desordem nos juizos e nos raciocinios} dá-se, emquanto, no \emph{sonho}, as imagens se reproduzem incompleta e desordenadamente. --- Nem se objecte que os juizos e os raciocinios dependem da intelligencia, e que esta não dorme. Por quanto, a intelligencia, embora não durma, não póde, no seu exercicio, prescindir do concurso das imagens, e por isso a desordem na imaginação reflecte-se necessariamente na intelligencia.\par 3.º) O \emph{somniloquio} dá-se, quando se desenvolve inteiramente a excitabilidade em qualquer parte do cerebro, e especialmente na terceira circumvolução frontal esquerda, de que partem os impulsos para os movimentos da palavra. --- \emph{Maury}, no seu livro ácerca dos \emph{sonhos}, diz que os circumstantes podem entrar em conversa com o \emph{somniloquo}, se o \emph{ouvido} d'este tiver a excitabilidade necessaria, e podem \emph{suggerir-lhes} qualquer idea ou acção.\par 4.º) O \emph{somnambulismo} resulta da viva reproducção e realização de imagens, relativas a uma certa ordem de operações habituaes a uma pessoa. Nestes casos a excitabilidade dos pontos cerebraes, onde as imagens residem, communica-se a todas as partes do systema nervoso, de que dependem as sensações e os movimentos necessarios para a realização d'aquella ordem de operações. O resto do eixo cerebro-espinhal está adormecido, e por isso as operações, as sensações e os movimentos do somnambulo são todos exclusivamente relativos ao \emph{sonho}.\par 5.º) A \emph{lucidez natural} dá-se, quando o homem, durante o somno, exerce algumas operações, que exigem o exercicio da intelligencia ou da razão. Nem isto admira. Sendo a excitabilidade cerebral, durante o somno, menos extensa e por isso mais energica, a alma concentra toda a sua attenção sobre certas imagens e ideas; e não repugna que, nesse estado, ás vezes realize algumas operações com maior perfeição do que na vigilia.\par Além d'este somno, que é \emph{natural}, porque deriva de causas \emph{naturaes}, ha um outro somno, que é \emph{artificial}, porque é provocado por meios \emph{artificiaes}, e que é conhecido geralmente com o nome de \emph{hypnotismo}.\par O \emph{hypnotismo}, que antigamente se chamava \emph{magnetismo animal}, é definido por \emph{Morselli} « un somno artificial, mais ou menos profundo, em que algumas regiões do cerebro ficam paralysadas, ao passo que outras soffrem uma excitação extraordinaria ».\par Este phenomeno, que não se realiza em todo e qualquer individuo, produz-se\pgn{281} por varios modos. Mas todos estes modos têm um só fim, que é \emph{fatigar} os \emph{nervos sensitivos} (sobretudo o nervo \emph{optico}) e por meio d'elles o \emph{sentido commum}. Para isso basta que o hypnotizado tenha os olhos fixos --- ou nos olhos do hypnotizador, --- ou nalgum ponto luminoso, --- ou ouça um som monotono, --- ou imagine que vai entrar no somno hypnotico.\par Durante o \emph{somno hypnotico} dão-se muitos e diversos phenomenos.\par Os principaes phenomenos, que podem dar-se no estado hypnotico, são o \emph{exclusivismo}, o \emph{automatismo imitativo} e a \emph{suggestão}.\par \emph{a}) O \emph{exclusivismo} consiste em que, durante o hypnotismo, ha uma coisa, que está sempre presente ao hypnotizado. Essa coisa é o rosto e a palavra do hypnotizador. Este \emph{exclusivismo} dá-se, porque o hypnotizado adormece com a persuasão de que está sob o influxo do hypnotizador, e a imagem d'este é a unica que fica vivamente gravada na sua phantasia.\par \emph{b}) O \emph{automatismo imitativo} consiste em que o hypnotizado repete \emph{machinalmente} todos os movimentos que o hypnotizador faz na sua presença. --- A explicação d'este phenomeno não é difficil. A experiencia attesta que a imaginação de um certo e determinado movimento nos leva por um instincto natural a executal-o; visto que os pontos do cerebro, onde existem as imagens do mesmo movimento, communicam com os respectivos nervos motores. Por isso, quando a imagem de um movimento é muito viva, d'aquelle ponto do cerebro parte uma corrente, que excita os nervos motores. --- Esta \emph{automatismo imitativo} dá-se não só no hypnotizado, mas tambem no homem acordado; só ha uma differença, que neste o instincto da \emph{imitação} obedece á razão e á vontade, e naquelle é cego e involuntario.\par \emph{c}) A \emph{suggestão} dá-se, quando o hypnotizador \emph{suggere} ao hypnotizado a imagem de uma coisa, e este opèra como se aquella coisa fosse real. Assim, se ao hypnotizado for \emph{suggerida} a imagem de um incendio ou de uma alluvião, elle fará no somno tudo o que nesses casos faria na vigilia. --- O hyponizado, não tendo deante de si outra coisa senão a pessoa e a palavra do hypnotizador, só tem as imagens, que provêm da vontade e da acção d'este. Taes imagens, sendo tanto mais vivas quanto mais são concentradas, suscitam na alma emoções correspondentes, e levam o sujeito a executar as operações, que as mesmas imagens representam; pois o hypnotizado, não podendo corrigir a allucinação, toma a imagem pelo objecto real.\par O Dr. Venturoli explica d'este modo a \emph{razão} d'esses phenomenos.\par « A verdadeira causa organica dos phenomenos hypnoticos reside num estado particular do systema nervoso; e todos os processos, empregados para provocar o \emph{hypnotismo}, consistem geralmente na prolongada repetição da mesma sensação e no cançaço de algum orgão sensitivo. A attenção fixa e continua sobre o mesmo objecto acaba por cançar o orgão \emph{sensitivo} e o \emph{sentido commum}, e aquelle estado de cançaço predispõe para o somno...\par « Além d'isso, a imaginação excitada no individuo, que sabe dever soffrer um estado produzido por actos extranhos e extraordinarios, têm uma parte consideravel na producção d'esses phenomenos; e, nesta perturbação extraordinaria do systema nervoso e das faculdades sensitivas, é facil passar para a catalepsia e para o somnambulismo.\par « A comparação dos phenomenos \emph{espontaneos} da catalepsia e do somnambulismo\pgn{282} com os \emph{artificiaes} confirma a identidade da origem e da natureza d'estes e d'aquelles: a unica differença está em que uns são \emph{espontaneos}, e outros \emph{artificiaes}. E, embora estes se manifestem de um modo mais surprehendente, isso é devido aos artificios empregados sobre certos individuos, de modo que neste caso póde dizer-se que a arte vence a natureza » (\emph{La scienza italiana}, 1883, vol. II, pag. 451).\par Alem d'estes phenomenos que, embora extraordinarios, não excedem comtudo as forças \emph{naturaes} do homem, e não differem essencialmente dos que se dão em certos individuos que têm alterado o systema nervoso, o \emph{hypnotismo}, segundo alguns, apresenta outros phenomenos, que excedem as forças \emph{naturaes} do homem e constituem o \emph{hypnotismo preternatural}.\par Estes phenomenos são a \emph{visão de coisas secretas e remotas}, a \emph{transposição dos sentidos}, o \emph{ler nos livros fechados}, o \emph{fallar uma lingua desconhecida}, etc.\par A existencia dos phenomenos \emph{superiores} do \emph{hypnotismo} é attestada por \emph{Cervello}, \emph{Semmola}, \emph{Lombroso}, \emph{Morselli}.\par Qual será a \emph{causa} d'esses phenomenos?\par Os phenomenos \emph{superiores} do \emph{hypnotismo}, \emph{se existem}, hão de ter necessariamente uma causa proporcionada. Essa causa não póde ser \emph{natural}; pois o homem não póde \emph{naturalmente} ver as coisas secretas e remotas, ouvir com as mãos, ler em livros fechados, fallar uma lingua que nunca aprendeu. Logo deve ser \emph{preternatural}. --- Essa causa \emph{preternatural} deve ser um ente \emph{immaterial}, dotado de \emph{intelligencia} e de \emph{vontade}; aliás não seria proporcionada. Ora o ente \emph{immaterial}, dotado de \emph{intelligencia} e de \emph{vontade}, é um \emph{espirito separado da materia}. Logo a causa dos phenomenos \emph{superiores} do \emph{hypnotismo} é um \emph{espirito separado da materia}. --- Esta conclusão parecerá exaggerada, mas é logica e racional, \emph{supposta} a existencia dos \emph{phenomenos superiores}.}.

%%>>> t168

\tese{168}{Séde da percepção sensitiva interna.}
Sendo varios e especificamente diversos os sentidos \emph{internos}, a percepção sensitiva \pg{280} \pg{281}  \emph{interna} não pode derivar de um só principio \emph{proximo} e \emph{immediato}. Haverá, portanto, \emph{quatro} principios \emph{proximos} e \emph{immediatos} \pg{282}  da percepção sensitiva \emph{interna}, porque \emph{quatro} são os sentidos \emph{internos}, especificamente diversos. --- Como, porém, os quatro sentidos \emph{internos} residem no \emph{cerebro}, podemos dizer que o \emph{cerebro} é a séde da percepção sensitiva interna \footnote{Da \emph{séde} dos sentidos \emph{internos} tratámos, com o devido desenvolvimento, no cap. II da secção II d'esta III parte da Metaphysica.}.

%%>>> art-24

\section*{Artigo IV.\\ Appetite sensitivo, sua divisão e paixões.}
\addcontentsline{toc}{section}{Artigo IV. --- Appetite sensitivo, sua divisão e paixões.}

169. \textbf{Appetite} sensitivo, --- \emph{Appetite sensitivo} é a operação, pela qual o homem tende para um bem, que é conveniente á sua natureza \emph{sensitiva} e que foi percebido pelos sentidos.

Esta operação existe no homem. Por quanto, existindo no homem o \emph{appetite sensitivo}, como \emph{faculdade}, deve tambem existir nelle, como \emph{operação}; visto que a cada faculdade corresponde uma operação proporcionada. --- A séde ou o orgão d'esta operação é, como vimos, o \emph{cerebro}.

{\small\noindent\textsc{Summario}: --- Conhecimento intellectual, seu objecto e origem. --- Modo por que se produz o conhecimento. --- Modo por que se desenvolve o conhecimento. --- Modo por que se exprime o conhecimento. --- Volição, seu objecto e divisão.\par}

%%>>> t170

\pg{283} \tese{170}{Divisão do appetite sensitivo.}
O \emph{appetite sensitivo}, considerado como \emph{operação}, tem a mesma divisão, a que está sujeito, considerado como \emph{faculdade}; pois de faculdades distinctas não podem deixar de derivar operações distinctas. Ora o \emph{appetite sensitivo}, considerado como \emph{faculdade}, divide-se em \emph{concupiscivel} e \emph{irascivel}, segundo tende para um bem \emph{agradavel} ou para um bem \emph{arduo}. Logo o \emph{appetite sensitivo}, considerado como \emph{operação}, divide-se em \emph{concupiscivel} e \emph{irascivel}.

%%>>> t171

\tese{171}{Paixões.}
As operações da faculdade \emph{appetitivo-sensitiva}, se forem elevadas a um intenso grau de força e de energia, tomam o nome especial de \emph{paixões}. --- \emph{Paixões}, portanto, são fortes movimentos da faculdade \emph{appetitivo-sensitiva}, determinados pela percepção do bem ou do mal \footnote{\emph{Descartes} diz que as \emph{paixões} são « percepções, sentimentos e emoções da \emph{alma}, causadas, conservadas e fortalecidas por quaesquer movimentos dos espiritos ». --- A definição de \emph{paixões}, dada por Descartes, é falsa. Por quanto, as \emph{paixões} são movimentos, affeições do \emph{composto organico}, e não da \emph{alma} só; pois são actos do \emph{appetite sensitivo}, e este é faculdade organica. Se o sujeito das paixões fosse a alma só, os animaes, cujas operações dependem da materia, não teriam paixões verdadeiras, mas apenas apparentes; o que é evidentemente falso.\par \emph{Mallebranche} diz que as \emph{paixões} são emoções proprias da alma só, suscitadas por Deus na \emph{occasião} de movimentos extraordinarios dos espiritos animaes e do sangue. --- Tambem esta opinião é falsa. Como dissemos contra \emph{Descartes}, a \emph{alma}, por si só, não póde ser sujeito de paixões. Alem d'isso, vimos que o \emph{occasionalismo}, em que \emph{Mallebranche} funda a sua theoria das \emph{paixões}, é absurdo.\par Pergunta-se: As \emph{paixões} são boas, ou más?\par Ha duas opiniões contrarias a este respeito. --- Uma, seguida por \emph{Epicuro}, \emph{Fourier}, \emph{Saint-Simon}, admitte que as paixões são boas ou virtuosas. --- Outra, defendida por \emph{Platão} e pelos \emph{Stoicos}, ensina que as paixões são más e viciosas, são doenças da alma.\par Ambas estas opiniões são falsas. Por quanto, as paixões, sendo operações do \emph{appetite sensitivo}, que é destituido de intelligencia e de vontade, não são, \emph{por si mesmas}, nem boas nem más, nem virtuosas nem viciosas; pois o principio e o sujeito da bondade e da malicia é a vontade, esclarecida pela intelligencia. --- Se as paixões, ás vezes, se chamam boas ou más, taes qualidades derivam para ellas --- ou das qualidades dos \emph{objectos}, para que tendem, --- ou da vontade, que se serve d'ellas, como de instrumentos, para a consecução dos seus fins. Por isso a \emph{virtude} e o \emph{vicio} são coisas \emph{accidentaes} para as \emph{paixões}.\par Alguns philosophos dizem que deveriamos extirpar as \emph{paixões}, porque impedem a perfeição da nossa alma.\par A opinião d'estes escriptores é inacceitavel. As paixões não são por si o resultado de uma corrupção, mas são a consequencia legitima e natural da natureza sensitiva. E', pois, impossivel extirpal-as, como é impossivel extirpar a propria natureza. --- Além d'isso, as paixões, quando bem dirigidas (e sem\pgn{284}pre o devem ser), não só não impedem a perfeição da alma, mas promovem-na e completam-na; pois que os esforços da vontade são mais energicos, efficazes e generosos, quando favorecidos e acompanhados pela \emph{paixão}.}.

%%>>> t172

\pg{284} \tese{172}{Numero das paixões.}
As \emph{paixões} são \emph{onze}, a saber: \emph{amor} e \emph{odio}, \emph{desejo} e \emph{abominação}, \emph{alegria} e \emph{tristeza}, \emph{esperança} e \emph{desespero}, \emph{audacia} e \emph{medo}, e, finalmente, \emph{ira}.

D'estas \emph{onze} paixões as \emph{seis} primeiras derivam do appetite \emph{concupiscivel}; as \emph{cinco} ultimas emanam do appetite \emph{irascivel}.

Na verdade, o \emph{bem} e o \emph{mal}, para que tende ou de que foge o appetite \emph{concupiscivel}, podem considerar-se --- ou simplesmente como \emph{taes}, e produzem o \emph{amor} e o \emph{odio}, --- ou como \emph{ausentes}, e dão origem ao \emph{desejo} e á \emph{abominação}, --- ou como \emph{presentes}, e produzem \emph{alegria} ou \emph{prazer}, \emph{tristeza} ou \emph{dor}. --- Do mesmo modo, a \emph{difficuldade} do bem, que é objecto do appetite \emph{irascivel}, é \emph{vencivel} ou \emph{invencivel}, e temos a \emph{esperança} ou o \emph{desespero}; ou o mal \emph{imminente} póde \emph{evitar-se}, ou não, e temos a \emph{audacia} ou o \emph{medo}; ou, finalmente, queremos vingar-nos contra a causa do nosso mal, e temos a \emph{ira} \footnote{Entre todas as paixões a principal é o \emph{amor}; pois todas têm o principio e o termo no amor. E com razão; visto que o bem é naturalmente anterior ao mal (pois o mal é a privação do bem), e não fugimos do mal, senão emquanto nos priva do bem. Por isso \emph{Bossuet} conclue a sua analyse (que reproduzimos na \emph{Cosmologia}) com as seguintes palavras: \emph{Otez l'amour, il n'y a plus de passion: et posez l'amour, vous les faites naître toutes} (Traité de la Conn. de Dieu et de soi-même, c. 1, art. 6).\par O \emph{amor}, o \emph{desejo} e a \emph{alegria} são paixões \emph{attractivas}; --- o \emph{odio}, a \emph{abominação} e a \emph{tristeza} são paixões \emph{repulsivas}; --- o \emph{desespero} e o \emph{medo} são paixões \emph{deprimentes}; a \emph{esperança}, a \emph{audacia} e a \emph{ira} são paixões \emph{excitantes}.\par O P. \emph{Manuel Bernardes}, no livro --- \emph{Luz e calor}, com a costumada elegancia de estylo e precisão de ideas, trata das \emph{paixões}. Eis as suas palavras:\par « Move-se o \emph{appetite} por onze differentes modos de actos seus vitaes e immanentes, isto é, que procedem da alma, como vida do corpo, e se recebem e ficam no mesmo appetite, que. se moveu. O nome mais commum (d'estes actos) é o de \emph{paixões}, em razão de que a elles se segue, ou acompanha alguma alteração e padecimento do corpo. E por isso, ainda que na vontade espiritual ha tambem actos espirituaes semelhantes a estes materiaes do \emph{appetite}, comtudo se não chamam propriamente \emph{paixões}, se não operações, porque não commovem e fazem padecer proxima e immediatamente o corpo.\par « Os nomes d'estes onze actos, ou \emph{paixões}, são: --- \emph{amor} e \emph{odio}, --- \emph{desejo} e \emph{fuga} ou \emph{abominação}, --- \emph{deleitação} e \emph{tristeza}, --- \emph{esperança} e \emph{desesperação}, --- \emph{temor} e \emph{audacia}, --- e, finalmente, \emph{ira}.\par « As primeiras \emph{seis} pertencem á parte \emph{concupiscivel}. Porque o bem e o mal --- ou se consideram absoluta e precisamente, e então o bem causa \emph{amor}, e o mal \emph{odio}; --- ou se consideram emquanto ausentes, e então o bem causa \emph{desejo}, e o mal \emph{fugida} ou \emph{abominação}; --- ou se consideram emquanto já presentes, e então o bem causa \emph{deleitação} ou \emph{gozo}, e o mal \emph{tristeza} ou \emph{dor};\pgn{285}supposto que os nomes do \emph{gozo} e \emph{dor} mais se attribuem ás coisas espirituaes, e os de \emph{deleite} e \emph{tristeza} ás materiaes.\par « As outras cinco \emph{paixões} pertencem ao \emph{irascivel}. Porque o bem arduo e difficultoso, emquanto possivel, termina a \emph{esperança}, --- emquanto impossivel, a \emph{desesperação}. O mal futuro, emquanto difficil de se evitar, termina o \emph{temor}; mas, emquanto se apresenta vencivel e evitavel, termina a \emph{audacia}. A \emph{ira} não tem parelha ou contrario; porque o bem presente não causa \emph{paixão} no \emph{irascivel}, como a causa o mal presente, que é a dita paixão da \emph{ira}, que se levanta, não para escapar d'esse mal, pois já se suppõe presente, senão para se resarcir e vingar do damno que lhe causa. E assim á \emph{ira} somente se lhe oppõe o amansar-se, que propriamente não é seu contrario, senão uma privação, ou negação da mesma \emph{ira} cessando.\par « Outras \emph{paixões} se nomeam, que se podem reduzir a estas \emph{onze}: como debaixo da \textsc{tristeza}, a \emph{acedia}, \emph{anxiedade}, \emph{angustia}, \emph{compaixão} ou \emph{misericordia}, a \emph{inveja} e o \emph{zelo}; --- debaixo da \textsc{ira}, a \emph{mania} e o \emph{furor}; --- debaixo do \textsc{temor}, estas seis ou sete especies, \emph{covardia}, \emph{pejo}, \emph{vergonha}, \emph{admiração}, \emph{estupor}, \emph{agonia}, e \emph{pavor}. E todos estes actos se podem applicar espiritualmente á vontade racional, emquanto concorre com o \emph{appetite apaixonado}.\par « Nos brutos tambem ha paixões tirando as que suppõem uso da razão, como a \emph{inveja} e \emph{vergonha}, tomadas propriamente » (I parte, Dout. III).\par \emph{Descartes} enumera apenas seis \emph{paixões}, a saber: \emph{admiração}, \emph{amor}, \emph{odio}, \emph{desejo}, \emph{alegria} e \emph{tristeza}.\par A enumeração das \emph{paixões}, feita por \emph{Descartes}, \emph{é incompleta} e \emph{falsa}. --- E' \emph{incompleta}; porque não enumera varias paixões, que todavia existem no homem e no animal, e são irreductiveis ás enumeradas por elle. --- E' \emph{falsa}; pois que a \emph{admiração}, embora possa suscitar uma paixão, todavia não é paixão, mas é um estado do \emph{espirito}.\par A analyse das \emph{principaes} paixões poderia levar-nos a julgar que ellas têm por causa sufficiente um raciocinio da intelligencia. Mas não é assim. Os objectos estão naturalmente cercados de certos caracteres sensiveis e differentes, conforme aquelles são bons ou maus, faceis ou difficeis, etc.; e estes differentes caracteres produzem no ser sensitivo differentes impressões, e por isso differentes imagens e tendencias. D'ahi o encontrarem-se as \emph{paixões} nos animaes mais perfeitos.}.

%%>>> t173

\pg{285} \tese{173}{Paixões do appetite concupiscivel.}
As \emph{paixões} do appetite \emph{concupiscivel}, como acabamos de dizer, são as \emph{seis} seguintes: \emph{amor e odio, desejo e abominação, alegria e tristeza}. --- De cada uma daremos uma breve explicação.

\textbf{a) Amor e odio.} --- \emph{Amor} é uma paixão, pela qual o homem tende para um bem sensivel, percebido pelos sentidos. --- \emph{Odio} é uma paixão, pela qual o homem foge de um objecto sensivel, que é mal, ou que foi percebido como tal.

\textbf{b) Desejo e abominação.} --- \emph{Desejo} é uma paixão, pela qual o homem tende para um bem sensivel, quando \emph{ausente}. --- \emph{Abominação}, ou \emph{aversão}, é uma paixão, pela qual o homem evita o mal \emph{ausente}.

c) Alegria e tristeza. --- \emph{Alegria}, ou \emph{prazer}, é a paixão, pela qual o homem goza do bem sensivel, actualmente \emph{possuído}. --- \emph{Tristeza}, ou \emph{dor}, é a paixão, pela qual o homem se afflige pelo mal sensivel, que \emph{presentemente} o opprime \footnote{Não podemos deixar de tratar, ainda que brevemente, do \emph{prazer} e da \emph{dor}.\par Indicaremos a \emph{definição}, e a \emph{divisão} d'estas duas paixões, e os \emph{requisitos} para a producção do \emph{prazer}.\par a) DEFINIÇÃO DO PRAZER E DA DOR. --- O \emph{prazer} é a satisfação, que se prova, quando a faculdade tende para um objecto em plena conformidade com a sua inclinação. A \emph{dor} é o que se experimenta, quando a faculdade é contrariada na tendencia para o seu objecto. --- Porquanto, a faculdade, que tende naturalmente para o seu objecto, approximando-se ou afastando-se d'elle, alcançando-o ou perdendo-o, faz com que se experimentem emoções agradaveis ou desagradaveis, cuja natureza varia segundo a natureza e as disposições da faculdade, e segundo os graus de bondade do proprio objecto. --- Alem d'isso, os physiologistas \emph{Binet} e \emph{Féré}, no seu livro --- \emph{Sensation et mouvement}, constataram que a sensação agradavel ou desagradavel corresponde sempre a um exaggero ou a uma diminuição de energia. --- Por isso, o \emph{prazer} e a \emph{dor} não são propriamente actos das faculdades sensitivas, mas acompanham os mesmos actos.\par Muitos escriptores modernos pensam que o \emph{prazer} e a \emph{dor} são phenomenos meramente \emph{passivos}, irreductiveis ás percepções e ás sensações, e que por isso devem ser productos de uma faculdade especial, distincta das faculdades perceptivas e appetitivas. Esta faculdade seria a \emph{sensibilidade}. --- Esta opinião é falsa. O \emph{prazer} e a \emph{dor} são phenomenos que, além de não serem meramente \emph{passivos} (pois gozar, por ex., de um espectaculo não é ser meramente \emph{passivo}), são communs a \emph{todas} as faculdades.\par Todo o erro do certos escriptores modernos consiste em procurar sempre a classificação dos \emph{phenomenos} psychicos, sem lhes importar a classificação das faculdades; de modo que a cada phenomeno attribuem uma faculdade distincta. Ora este processo é falso. Para se fazer uma exacta classificação dos \emph{phenomenos}, é necessario ter sempre em vista o principio, de que derivam, e considerar se este póde ser a causa sufficiente d'aquelles; aliás --- ou separar-se-hão factos, que estão naturalmente unidos, e que são reductiveis ao mesmo principio, --- ou reunir-se-hão factos que, embora \emph{apparentemente} semelhantes, são comtudo \emph{essencialmente} diversos. --- Por terem sido esquecidos estes principios, creou-se sem necessidade uma faculdade \emph{especial} para o \emph{prazer} e para a \emph{dor}, e dividiram-se os phenomenos da alma em \emph{sensibilidade}, \emph{conhecimento} e \emph{actividade}, confundindo assim a percepção e a actividade \emph{sensitiva} com a percepção e a actividade \emph{espiritual}.\par b) Divisão DO PRAZER E DA DOR. --- O \emph{prazer}, como tambem a \emph{dor}, divide-se em \emph{sensivel} e \emph{espiritual}. \emph{O prazer sensivel} e a \emph{dor sensivel} affectam as faculdades \emph{sensitivas}. O \emph{prazer espiritual} e a \emph{dor espiritual} affectam as faculdades \emph{espirituaes}. Porquanto, não é só nas faculdades \emph{sensitivas} que sentimos o prazer e a dor, mas tambem na intelligencia e na vontade. Os\pgn{287}prazeres e as dôres nas faculdades intellectuaes são phenomenos mais nobres e elevados, e não menos intensos e vivos, que os prazeres e as dôres nas faculdades inferiores. Quem não provou como por vezes é \emph{laboriosa} a investigação da verdade, mas tambem como é \emph{doce} a sua contemplação?!\par c) REQUISITOS PARA A PRODUCÇÃO DO PRAZER. --- Sendo o prazer um effeito do acto, que a faculdade exerce em plena conformidade com a sua inclinação, é necessario que o mesmo \emph{acto} possua os seguintes requisitos:\par a) \emph{Não deve ser impedido}. Por quanto, a inacção, mais ou menos forçada, das nossas faculdades é causa de afflicção.\par \emph{b}) \emph{Deve ser conforme á inclinação da faculdade}. Com effeito, o exercicio violento de uma faculdade produz tristeza ou dôr.\par \emph{c}) \emph{O seu objecto deve ser proporcionado á natureza e ás actuaes disposições da faculdade}. Quanto maior fôr a proporção, tanto mais vivo será a prazer.\par \emph{d}) \emph{Deve tambem ser proporcionado á propria faculdade}. Este requisito exige que a medida da intensidade e da duração de um acto esteja em proporção com a respectiva faculdade. --- Semelhante medida é estreita para as faculdades \emph{organicas}, mas é bastante larga para as faculdades \emph{intellectuaes}. Um objecto sensivel, se for muito vivo e intenso, cança e perturba as faculdades sensitivas; assim uma luz demasiadamente viva perturba a vista. --- Não acontece o mesmo com o objecto das faculdades \emph{espirituaes}. A verdade, o bem, o bello, --- estes objectos nunca saciam a intelligencia e a vontade.\par Todavia o funccionamento das faculdades \emph{superiores} exige uma certa moderação, não por causa do seu \emph{objecto}, mas por causa das funcções organicas que o acompanham, e por causa da harmonia das faculdades, que poderia ser perturbada pelo demasiado desenvolvimento d'uma só faculdade.\par D'ahi se segue que o \emph{prazer} perfeito e illimitado só se encontra em \textsc{Deus}, \emph{acto} purissimo e perfeitissimo, e que a verdadeira felicidade do homem consiste na participação da mesma vida divina, para alcançar a qual, como diz Aristoteles (\emph{Moral a Nic.}, l. x. c. 7), devemos sacrificar os prazeres caducos e miseraveis da terra. (Cf. A. Farges --- \emph{Le Cerveau}, \emph{l'âme et ses facultés}).}.

%%>>> t174

\pg{286} \tese{174}{Paixões do appetite irascivel.}
As \emph{paixões} do appetite \emph{irascivel}, como dissemos, são \emph{cinco}, a saber: \emph{esperança} e \emph{desespero}, \emph{audacia} e \emph{medo}, e, finalmente, \emph{ira}.

a) Esperança e desespero. --- \emph{Esperança} é a paixão, pela qual o homem tende para um bem \emph{sensivel}, que, embora \emph{arduo}, é \emph{possivel} de se alcançar. --- \emph{Desespero} é a paixão, que nasce quando a consecução do bem sensivel é ou parece \emph{impossivel}.

b) Audacia e medo. --- Audacia é a paixão, pela qual o homem \emph{se esforça} para alcançar um bem \emph{arduo}, ou para fugir de um mal grave. --- \emph{Medo} é a paixão, pela qual o homem tenta afastar de si um mal imminente, difficil de se evitar.

c) Ira. --- \emph{Ira} é a paixão, pela qual o homem procura escapar ao mal \emph{presente}, que o opprime, --- ou, ao menos, vingar-se da causa que o infligiu \footnote{A \emph{ira} não tem contrario, a não ser que se colloque entre as paixões a inclinação de fazer bem a quem nos faz mal. Mas a collocação seria violenta e irracional; pois essa inclinação não traz comsigo a emoção ou a perturbação, que as paixões costumam trazer.}.

\textbf{OPERAÇÕES INTELLECTUAES} DO HOMEM

%%>>> art-25

\section*{Artigo I.\\ Conhecimento intellectual, seu objecto e origem.}
\addcontentsline{toc}{section}{Artigo I. --- Conhecimento intellectual, seu objecto e origem.}

%%>>> t175

\pg{287} \pg{288} \tese{175}{Operações intellectuaes do homem.}
As operações \emph{intellectuaes} do homem são as que emanam \emph{immediatamente} das suas faculdades \emph{intellectuaes}. --- E, como as faculdades \emph{intellectuaes} do homem são a \emph{intelligencia} e a \emph{vontade}, as suas operações \emph{intellectuaes} são o \emph{conhecimento intellectual} e a \emph{volição} \footnote{A \emph{volição} chama-se operação \emph{intellectual}, não porque deriva da intelligencia, mas porque a vontade é uma faculdade \emph{intellectual}, isto é, que deriva da intelligencia. Operação \emph{intellectual}, pois, é o \emph{genero}, a que estão subordinadas duas especies: o \emph{conhecimento intellectual} e a \emph{volição}.}.

Neste capitulo tratamos primeiramente do \emph{conhecimento intellectual}, da sua origem, do modo por que se exerce, desenvolve e exprime, e em seguida da \emph{volição} \footnote{Não é necessario que demonstremos aqui de proposito a \emph{espiritualidade} das operações \emph{intellectuaes}. Basta o que dissemos ácerca da \emph{espiritualidade} da alma e da intelligencia. Apresentamos apenas o argumento seguinte: Uma operação diz-se \emph{espiritual}, quando o seu objecto é \emph{espiritual} ou \emph{immaterial}. Ora o objecto das operações \emph{intellectuaes}, que o homem produz, é \emph{espiritual} ou \emph{immaterial}; porque a intelligencia percebe e a vontade ama a \emph{justiça}, a \emph{virtude}, a \emph{sciencia}, \textsc{Deus}, que são objectos \emph{espirituaes} ou \emph{immateriaes}. Logo as operações \emph{intellectuaes} são \emph{espirituaes}.}.

%%>>> t176

\tese{176}{Conhecimento intellectual.}
\emph{Conhecimento intellectual} é uma operação, pela o homem percebe a \emph{essencia} das coisas. --- Percebe-se, como dissemos, a \emph{essencia} das coisas, não só quando se indicam \emph{todos} os seus elementos \emph{especificos}, mas tambem, embora menos perfeitamente, quando se aponta um ou outro dos elementos mais \emph{genericos} ou \emph{transcendentaes} \footnote{Dissemos que o \emph{conhecimento} é uma operação. Porquanto, a operação póde ser considerada, --- ou em relação ao principio de que deriva, --- ou em relação ao termo, para que tende. Aqui a operação \emph{intellectual} considera-se em relação ao seu termo, que é o \emph{conhecimento}.}.

%%>>> t177

\tese{177}{Objecto do conhecimento intellectual.}
Objecto do \emph{conhecimento intellectual} é a \emph{essencia} das coisas. --- Divide-se em \pg{289}  \emph{adequado} e \emph{proporcionado}. O \emph{adequado} é a \emph{essencia} de todas as coisas, considerada no seu mais fundamental e commum elemento, que é o \emph{ente}. O \emph{proporcionado} é a \emph{essencia} das coisas \emph{materiaes}, \emph{abstrahida das condições materiaes} \footnote{No art. I do cap. III da secção II d'este tratado (n. 114) explicamos, com muito desenvolvimento, este ponto importantissimo.\par E' sabido que, para os \emph{materialistas}, \emph{sensualistas} e \emph{positivistas}, o objecto das operações intellectuaes não é o \emph{universal} e o \emph{immaterial}, mas é o \emph{concreto} e o \emph{material}. --- E' verdade que os \emph{sensualistas} affirmam que a \emph{imagem}, formada pela \emph{phantasia}, \emph{é universal}. Mas a affirmação é falsa; pois aquella \emph{imagem} representa uma coisa concreta, e por isso não é \emph{universal}.}.

%%>>> t178

\tese{178}{Origem do conhecimento intellectual.}
E' certo que a nossa intelligencia possue muitos e diversos conhecimentos. --- Mas d'onde provêm elles? Qual é a causa sufficiente e necessaria da sua presença no nosso espirito? --- E' a celebre questão ácerca da \emph{origem das ideas}.

A esta questão os philosophos deram varias soluções \footnote{Esta questão ácerca da origem das \emph{ideas} refere-se propriamente á origem das \emph{especies intelligiveis}, que são o principio do conhecimento intellectual; pois, recebidas as \emph{especies} na intelligencia, esta immediatamente produz o seu acto, que é o conhecimento. --- Esta parte da \emph{Anthropologia} costuma chamar-se propriamente \emph{ideologia}.}.

%%>>> t179

\tese{179}{Systemas ácerca da origem do conhecimento intellectual.}
Os systemas ácerca da origem do conhecimento intellectual são \emph{seis}; o \emph{innatismo}, o \emph{tradicionalismo}, o \emph{ontologismo}, o \emph{sensualismo}, o \emph{idealismo} e o \emph{psychologismo racional}.

Na verdade, o conhecimento ou é \emph{innato} ou é \emph{adquirido}.

Se admittimos que é \emph{innato}, temos o \emph{innatismo}.

Se concedemos que é \emph{adquirido}, devemos dizer que deriva ou de um principio \emph{extrinseco}, ou de um principio \emph{intrinseco}.

Se deriva de um principio \emph{extrinseco}, este principio --- ou é a \emph{auctoridade}, e temos o \emph{tradicionalismo}, --- ou é a \emph{visão immediata e intuitiva de} \textsc{Deus}, e temos o \emph{ontologismo}.

Se deriva de um principio \emph{intrinseco}, este é constituido --- ou \emph{exclusivamente} pelos \emph{sentidos}, e temos o \emph{sensualismo}, --- ou \emph{exclusivamente} pela \emph{intelligencia}, e temos o \emph{idealismo}, --- ou \emph{conjunctamente} pela \emph{intelligencia} e pelos \emph{sentidos}, isto é, pela intelligencia auxiliada pelos sentidos, e temos então o \emph{psychologismo racional}.

De cada um dos mencionados systemas daremos a definição e faremos uma breve critica \footnote{O primeiro systema, que os philosophos inventaram para explicar a origem dos conhecimentos \emph{intellectuaes} é o \emph{sensualismo}. Como reacção\pgn{290}contra o \emph{sensualismo}, appareceram o \emph{innatismo} e o \emph{idealismo}. Para destruir as consequencias do \emph{idealismo}, surgiram o \emph{ontologismo} e o \emph{tradicionalismo}. O \emph{psychologismo racional} evitou os dois extremos, dizendo que a origem dos nossos conhecimentos depende de uma causa \emph{intrinseca}, que é a propria \emph{intelligencia}, auxiliada pelos \emph{sentidos}.}.

%%>>> t180

\pg{290} \tese{180}{Innatismo.}
O \emph{innatismo} ensina que os nossos conhecimentos intellectuaes não são adquiridos, mas \emph{innatos} \footnote{O \emph{innatismo} foi defendido --- na antiguidade por \emph{Platão}, --- e nos tempos modernos por \emph{Descartes} e \emph{Leibnitz}. --- Estes escriptores concordam em que as ideas da nossa intelligencia são \emph{innatas}, mas divergem no \emph{modo} por que explicam este \emph{innatismo}.\par a) \emph{Platão} propoz a seguinte theoria: A alma humana, antes de ser unida ao corpo, viveu, \emph{num mundo ideal}, mais pura e mais unida a \textsc{Deus}, contemplando as ideas eternas. Quando, em castigo das suas culpas, foi encarcerada no corpo, perdeu completamente a memoria dos antigos conhecimentos. Todavia, os entes visiveis, quando actuam nas faculdades sensitivas, por serem formados á imagem dos \emph{archetypos} eternos, determinam a intelligencia á recordação dos conhecimentos adquiridos na primitiva visão divina, e latentes na memoria num estado de sopor. D'este modo a sciencia da alma humana seria apenas uma \emph{reminiscencia}.\par b) \emph{Descartes} divide todos os nossos conhecimentos intellectuaes em tres classss: \emph{adventicios}, \emph{facticios} e \emph{innatos}. Os conhecimentos \emph{adventicios} são os que adquirimos por intermedio dos sentidos, excitados pela impressão dos seres corporeos: como, por exemplo, a idea do \emph{sol}. Os \emph{facticios} são os que o nosso espirito fórma, e aos quaes nenhuma realidade corresponde; tal é a idea do \emph{centauro}. \emph{Innatos} são os conhecimentos, que foram impressos, pela acção divina, no nosso espirito, quando este foi creado, e que, pela acção dos sentidos, se recordam e tornam \emph{reflexos}; tal é a idea do \emph{infinito}, que só \textsc{Deus} (diz elle) podia escrever no nosso espirito. --- Todavia \emph{Descartes}, pouco coherente nos seus principios, affirma explicitamente que em nós não são \emph{innatos} os \emph{conhecimentos}, mas que é \emph{innata} a \emph{faculdade} de os adquirir, e que só neste sentido deve interpretar-se o \emph{innatismo} da idea do \emph{infinito}. (Cf. \emph{Lettres Philosophiques}, lettre 33).\par c) \emph{Leibnitz}, seguido por \emph{Wolf}, diz que a nossa alma, no momento da sua creação, recebe de \textsc{Deus} um schema ou uma representação \emph{confusa} de todo o mundo, e que, em seguida, esse schema se desenvolve em ideas particulares, claras e distinctas, á medida que os objectos sensiveis \emph{parecem} operar no corpo, e, por meio d'este, na alma. --- D'este modo a sciencia consistiria em tornar claros e distinctos os conhecimentos, que já teriam existido no nosso espirito num estado obscuro e confuso.}.

%%>>> t181

\tese{181}{O innatismo é falso.}

a) Todo o systema, que é contrario á experiencia, é falso. Ora o \emph{innatismo} é contrario á experiencia, a qual attesta que os nossos conhecimentos não existiam no nosso espirito, quando este começou a existir, mas que foram e vão sendo adquiridos de dia para dia. Logo o \emph{innatismo} é falso.

b) Todo o systema, que não está em harmonia com a con- \pg{291}  dição da nossa intelligencia, é falso; pois que a operação é um effeito, e todo o effeito deve ser proporcionado á causa. Ora o \emph{innatismo} não está em harmonia com a condição da nossa intelligencia; pois que esta, por estar unida ás faculdades sensitivas, deve, no seu exercicio e desenvolvimento, experimentar o influxo e o auxilio d'aquellas faculdades. Logo o \emph{innatismo} é falso.

c) Todo o systema, que leva ao \emph{idealismo}, é falso; porque, como veremos, os nossos conhecimentos são \emph{objectivos} ou \emph{reaes}. Ora o \emph{innatismo} leva ao \emph{idealismo}; pois os nossos conhecimentos, se fossem \emph{innatos}, seriam meras modificações da intelligencia, ás quaes nenhuma realidade corresponderia, visto que a nossa alma, por não ser a causa dos entes creados, não os póde representar. Logo o \emph{innatismo} é falso \footnote{Os adversarios respondem que os conhecimentos \emph{innatos}, pelo facto de terem sido infundidos no espirito por \textsc{Deus}, são \emph{reaes} ou \emph{objectivos}. --- Este recurso é vão. A \emph{realidade} ou \emph{objectividade} dos nossos conhecimentos é um facto natural, e um facto natural deve derivar de uma causa natural, e não da acção immediata de \textsc{Deus}.\par Nenhuma idea pode ser \emph{innata}, nem a de \emph{ente}. --- Porquanto, como dissemos muitas vezes, a nossa intelligencia, sendo faculdade de uma alma, que é forma substancial do corpo, deve operar de um modo proporcionado á sua natureza e ás condições da sua existencia. Por isso nenhuma idea pode despontar ou formar-se sem a cooperação dos sentidos. A nossa intelligencia é, como diz Aristoteles (\emph{De anima}, l. III; t. 14), \emph{sicut tabula in qua nihil est scriptum}. Com o andar do tempo, adquire ella muitas e varias ideas, mas sempre relativas a objectos, que \emph{de algum modo} actuaram nos sentidos; e por isso, se um sentido faltar, faltarão tambem as ideas dos objectos, perceptiveis por esse sentido (assim um cego de nascimento não tem a idea das côres). --- Nem a idéa de \emph{ente} é innata. Tambem ella, como a de \emph{unidade}, \emph{bondade}, ecc., é abstrahida dos coisas sensiveis ou dos phantasmas pela luz da intelligencia agente.\par A razão é sempre a mesma. Diz S. Thomaz: « Primae conceptiones intellectus (possibilis) lumine intellectus agentis cognoscuntur per species a sensibilibus abstractas, sive sint complexae (os \emph{axiomas}), sive incomplexae (\emph{as primeiras ideas}), ut ratio \emph{entis} et \emph{unius} et hujusmodi » (\emph{De mag.}, a. 1; --- cf. \emph{Summ, Theol.}, p. I, q. 84, aa. 3 et 4). --- E' necessario nunca perder de vista a grande verdade que S. Thomaz exprime nas seguintes palavras: « Naturalis nostra cognitio a sensu principium sumit; unde tantum se nostra naturalis cognitio extendere potest, in quantum manuduci potest per sensibilia » (\emph{Sum. Theol.}, p. I, q. 12, a. 12).\par Tendo no texto refutado o \emph{innatismo} em geral, devemos dizer alguma coisa ácerca das suas aspecies.\par a) O systema de \emph{Platão} é falso pelas seguintes razões:\par 1.º) Suppõe que as almas preexistem aos corpos, e que a união d'aquellas com estes é violenta e accidental. Ora, como demonstraremos, é falso\pgn{292}que as almas preexistem aos corpos, e, como vimos, a união da alma com o corpo não é violenta nem accidental, mas é natural e substancial.\par 2.º) O \emph{mundo ideal}, em que as almas teriam haurido os seus conhecimentos, e o sopor da alma, proveniente da sua união com o corpo, são invenções gratuitas e repugnantes.\par 3.º) A consciencia attesta que, quando aprendemos alguma coisa, não \emph{recordamos} um conhecimento passado e esquecido, mas alcançamos um novo conhecimento, que nunca tivemos.\par b) O systema de \emph{Descartes} é inacceitavel pelas seguintes razões:\par 1.º) Suppõe que a essencia da alma consiste no \emph{pensamento}; pois que, segundo este escriptor, a alma, assim como recebe de \textsc{Deus} a essencia, assim tambem recebe o pensamento. Ora essa supposição é falsa.\par 2.º) Não ha razão alguma que justifique a distincção entre ideas \emph{adventicias} e \emph{innatas}. Segundo Descartes, algumas ideas são \emph{innatas} pelo facto de serem \emph{universaes} e \emph{necessarias}; mas ha muitas ideas, que são \emph{adventicias} e que todavia são \emph{universaes} e \emph{necessarias}.\par c) O systema de \emph{Leibnitz} é falso pelas seguintes razões:\par 1.º) A nossa alma, sendo finita e especificamente distincta das outras creaturas, não póde ter a representação, embora confusa, do mundo inteiro. Porquanto, só \textsc{Deus}, por ser a causa de todas as coisas, contêm na sua essencia infinita a semelhança d'ellas.\par 2.º) Admitte uma \emph{harmonia preestabelecida} entre a alma e o corpo, --- \emph{harmonia} que, como vimos na \emph{Secção primeira}, repugna.}.

%%>>> t189

\pg{292} \tese{182}{Tradicionalismo.}
O \emph{tradicionalismo} ensina que os nossos conhecimentos \emph{intellectuaes} derivam unica e exclusivamente da \emph{tradição}, ou da \emph{auctoridade}, divina ou humana \footnote{O \emph{tradicionalismo} divide-se em \emph{rigoroso} e \emph{mitigado}.\par a) O tradicionalismo \emph{rigoroso}, ensinado por \emph{Bonald} e \emph{Lamennais}, ensina que o homem, sem a tradição oral e por isso sem a revelação, não póde conhecer nenhuma verdade; porque o homem não póde pensar sem o auxilio da \emph{linguagem}, e esta não póde ser inventada pelo homem.\par b) O tradicionalismo \emph{mitigado}, que teve. por principaes auctores e defensores M. \emph{Bonetty} e o P. Ventura, sustenta que a tradição \emph{oral} e por isso a \emph{linguagem} não são necessarias para descobrir \emph{todas} as verdades, mas só para adquirir as verdades da ordem \emph{metaphysica} e da \emph{moral}.}.

\textbf{183. O tradicionalismo é falso.}

a) E' falso no seu \emph{fundamento}. Este systema admitte que a nossa intelligencia não póde formar ideas universaes, nem conhecer os mais geraes principios da metaphysica e da moral. Ora tal opinião é falsa. Com effeito, a nossa intelligencia, sendo dotada de força comparativa e abstractiva, forma ideas universaes, e, apoiando-se no principio de \emph{causalidade} e na idea do bem, eleva-se dos effeitos ás causas, e descobre muitas verdades da ordem especulativa e moral. Logo o \emph{tradicionalismo} é falso no seu \emph{fundamento}.

b) E' falso na sua \emph{essencia}. O \emph{tradicionalismo} ensina que \pg{293}  todos os nossos conhecimentos intellectuaes derivam exclusivamente da \emph{revelação}, do \emph{magisterio externo} e da \emph{palavra}. Ora isto é falso. Por quanto, a \emph{revelação}, o \emph{magisterio externo} e a \emph{palavra}, embora sejam \emph{auxilios} para a acquisição de \emph{muitos} conhecimentos, não podem ser a \emph{unica} e \emph{exclusiva causa} de \emph{todos}; pois que o acto de fé suppõe um acto de razão, visto que não acreditariamos na auctoridade dos outros, se não conhecessemos que é necessario acreditar. Logo o \emph{tradicionalismo} é falso na sua \emph{essencia}.

c) Leva a \emph{consequencias} absurdas. Com effeito, suppondo o acto de fé um acto de razão, esta faculdade, se fosse por si mesma incapaz de alcançar toda e qualquer verdade, não poderia acceitar nenhuma das verdades ensinadas pela tradição; porque não poderia conhecer com certeza a auctoridade do testemunho. D'ahi o \emph{scepticismo universal} \footnote{Tendo refutado, no texto, o \emph{tradicionalismo} em geral, refutemos agora a sua duplice especie.\par a) O tradicionalismo \emph{rigoroso} é falso.\par a) Porquanto, este systema suppõe que o homem não póde pensar nem reflectir, sem o auxilio da \emph{linguagem}. Ora esta supposição é falsa, --- não só porque a idea procede a \emph{linguagem} (primeiramente pensamos e depois manifestamos o pensamento), --- mas tambem porque a \emph{linguagem} não póde entender-se, sem primeiro se entender a relação \emph{convencional} entre as palavras e as ideas, que ellas exprimem (aliás a palavra não poderia suggerir idea alguma). --- Para um homem destituido de ideas a linguagem é incomprehensivel (Cf. S. Thom. \emph{opusc}. 13). --- Por isso o principio da evolução da nossa intelligencia não depende do magisterio externo (aliás a intelligencia nunca chegaria a desenvolver-se); mas esse magisterio limita-se a ensinar o modo de comparar as ideas, que o discipulo já possue, e o modo de as deduzir umas de outras.\par \emph{b}) De mais, é falso que o homem não possa inventar a \emph{linguagem}. Por quanto, o homem tem uma tendencia natural --- não só para pensar e para manifestar os seus pensamentos, --- senão tambem para a producção organica da palavra. Ora isto basta para elle poder indicar uma coisa material com um som \emph{articulado}, e depois servir-se das palavras, que indicam coisas materiaes, para indicar as immateriaes.\par Advertimos que, embora o homem \emph{possa} inventar a \emph{linguagem}, comtudo é certo que de \emph{facto} não a inventou, mas lhe foi infundida pelo Creador. --- E aqui não podemos deixar de referir uma theoria dos \emph{racionalistas}, seguida ultimamente por \emph{E. Renan} (\emph{De l'origine du langage}, c. XII), segundo a qual a \emph{linguagem} seria um producto \emph{espontaneo} do homem. Esta theoria é falsa. Se assim fosse, todos os homens fallariam a mesma lingua, e, para a aprenderem, não precisariam de ensino. --- Voltaremos ao assumpto.\par b) O tradicionalismo \emph{mitigado} é falso. Porquanto, a nossa intelligencia, para se elevar ao descobrimento das principaes verdades da ordem metaphysica e moral, basta que applique os \emph{principios} aos casos concretos e\pgn{294}particulares. Está averiguado que homens, que nunca ouviram o som da palavra humana, só pela contemplação das creaturas se elevaram á contemplação do Creador, e descobriram os principios da moralidade e do bem.\par A nossa conclusão é certa tambem \emph{theologicamente}. --- O Concilio Vaticano definiu que a razão pode demonstrar e demonstra os fundamentos da fé, que são a existencia de Deus, a sua veracidade, o facto da revelação sobrenatural, etc. (Const. \emph{De fide cath}.). --- A S. Congregação do \emph{Index} obrigou M. \emph{Bonnetty} a subscrever as seguintes proposições: --- 1) Ratiocinatio Dei existentiam... cum \emph{certitudine} probare potest. --- 2) Rationis \emph{usus fidem praecedit}... et ad eam conducit. (Cf. \emph{Denzinger}, n. 1506).\par Todavia é necessario evitar os extremos. Se a razão humana não precisa da tradição nem da revelação para se elevar aos primeiros principios da ordem metaphysica e moral, não póde concluir-se que a divina revelação é inutil. Por quanto a revelação divina abrange as verdades da ordem \emph{sobrenatural}, e as verdades da ordem \emph{natural}. Ora a revelação das verdades \emph{sobrenaturaes}, supposta a elevação do homem ao estado sobrenatural, não só é util, mas é \emph{necessaria}. A revelação das verdades \emph{naturaes}, se não é necessaria, é \emph{util}; porque, sem ella, poucos homens, depois de muito tempo, teriam aprendido poucas verdades, e essas misturadas com erros, quando \emph{todos} os homens, desde o principio da sua existencia e sem mistura d'erro, devem saber as verdades, de que depende a sua felicidade, presente e futura. Diz S. Thomaz: « Salubriter divina providit clementia ut ea etiam, quae ratio investigare potest, fide tenenda praeciperet, ut sic omnes de facili possent divinae cognitionis participes esse et absque dubitatione et errore » (C. \emph{Gentes}, l. I, c. 4).}.

184. \textbf{Ontologismo}. --- O \emph{ontologismo} ensina que a nossa intelligencia tem uma \emph{intuição directa} e \emph{constante} da essencia de \textsc{Deus}, e que, existindo em \textsc{Deus} as ideas archetypas de todas as coisas creadas, nelle contemplamos o universo \footnote{O \emph{ontologismo} teve \emph{Nicolau Mallebranche} por fundador, e V. \emph{Gioberti} e A. \emph{Rosmini} por principaes defensores.\par a) \emph{Mallebranche} dizia que \textsc{Deus}, não podendo ser percebido por uma idea distincta d'elle, é objecto de uma \emph{intuição} ou visão \emph{directa} e \emph{immediata}, e que nessa visão percebemos todas as outras coisas, á excepção do sujeito pensante, que é attingido pela consciencia.\par b) V. \emph{Gioberti} resuscitou, neste seculo, o systema \emph{ontologico}. A theoria d'este escriptor, embora na essencia concorde com a de \emph{Mallebranche}, todavia nas circumstancias differe sensivelmente. Diz assim: Deve distinguir-se um duplice pensamento: um \emph{immediato} e \emph{directo}, que é a \emph{intuição}; outro \emph{mediato} e \emph{indirecto}, que é a \emph{reflexão}, pela qual consideramos e explicamos o objecto percebido pela \emph{intuição}. --- A \emph{reflexão} não póde exercer-se senão por meio de signaes; d'ahi a necessidade da linguagem, communicada pela revelação. A \emph{intuição}, porém, é natural ao homem desde o principio da vida intellectual, e o seu objecto immediato é a \emph{Idea}, isto é, a Verdade absoluta e eterna, que é \textsc{Deus}. --- Identificando-se a primeira Idea do nosso espirito com o primeiro Ente, segue-se que o primeiro \emph{psychologico} se confunde com o \emph{primeiro ontologico}, e ambos formam um \emph{primeiro}, que é o \emph{primeiro philosophico}. --- Ora este Ente, que está sempre presente ao nosso pensamento, percebe-se como é em si mesmo, isto é, actualmente creador de todas as\pgn{295}coisas finitas. Nessa percepção, pois, o nosso espirito descobre tres realidades: o Ente supremo, a creação livre, e as coisas finitas existentes pela creação. D'ahi a formula: o \emph{Ente cria o existente}, isto é: \emph{Deus cria o mundo}.\par c) \emph{Rosmini} ensinou o \emph{ontologismo} da maneira seguinte. Disse que \textsc{Deus}, quando produz a nossa alma, irradia nella a idea universal do ser, ou o \emph{ser ideal}, em que a alma vê tudo. Este \emph{ser ideal}, segundo \emph{Rosmini}, \emph{é o proprio Deus, o Verbo divino} (\emph{Rinnovamento della Filosofica}, lib. 3, c. 52). --- D'este modo « \emph{na esphera do creado manifesta-se á intelligencia humana alguma coisa de divino em si mesmo, isto é, que pertence á natureza divina} » (\emph{Teosofia}, v. 4, n.º 2). --- Nas suas ultimas obras, este escriptor dá ao ser \emph{ideal} o nome de ser \emph{virtual} e \emph{inicial}. « O ser \emph{virtual} e \emph{inicial} (diz \emph{Rosmini}) ou o \emph{ser visto por intuição natural} é necessario. Mas elle não é ente: logo é alguma coisa de um ente. Mas esse ente, de que aquelle é alguma coisa, não póde ser contingente, pois que o contingente é opposto ao necessario. Logo o ser, que o homem contempla por \emph{intuição}, é alguma coisa de um \emph{ente necessario} e \emph{eterno}, que cria, determina e limita todos os entes contingentes, e este ser é \textsc{Deus} » (\emph{Teosofia}, v. 1, l. 2, sez. 2, c. 2 a 5, n.º 298). --- Depois d'esta manifestação do ente \emph{ideal} ou \emph{virtual}, a nossa alma applica ás sensações d'este ou d'aquelle objecto a \emph{idea} do ser \emph{virtual} ou \emph{inicial}, e pronuncia o juizo: o \emph{que sinto é}. Estes juizos formam os nossos conhecimentos.}.

185. O ontologismo é \textbf{falso}.

a) Este systema admitte que \textsc{Deus} é o objecto, \emph{continua} e \emph{intuitivamente} presente á nossa intelligencia; de modo que esse conhecimento de \textsc{Deus} é a causa de todos os outros conhecimentos. Ora isto é falso. Nenhum homem tem consciencia de que é dotado da visão \emph{intuitiva} de \textsc{Deus}, --- visão que, se existisse, não escaparia á reflexão, tratando-se do principal objecto do nosso pensamento. Logo o \emph{ontologismo} é falso \footnote{Nem se diga que, nesta visão de Deus, não contemplamos propriamente a \emph{Essencia divina}, mas só as \emph{ideas archetypas} das coisas creadas. Em Deus as ideas archetypas identificam-se com a propria Essencia divina, porque essas ideas são a propria Essencia emquanto se refere ás creaturas. Além d'isso, o conhecimento \emph{relativo} de um ente é \emph{posterior} ao seu conhecimento \emph{absoluto}, isto é, um ente percebe-se primeiramente em si e depois com relação aos outros entes. Diz o Angelico Doutor: « Non est possibile quod aliquis videat rationes creaturarum in ipsa divina essentia, ita quod eam non videat; tum quia ipsa divina Essentia est ratio omnium eorum quae fiunt, ratio autem idealis non addit supra divinam essentiam, nisi respectum ad creaturas; tum etiam quia prius est cognoscere aliquid in se, quod est cognoscere Deum, ut est objectum beatitudinis, quam illud per comparationem ad alterum, quod est cognoscere Deum secundum rationes rerum in ipso existentes » (\emph{Sum. Theol.}, 2-2, q. 173, a. 1; cf. \emph{De ver.}, q. 12, a. 6).}.

b) A visão \emph{directa} ou \emph{intuitiva} de \textsc{Deus} não só não existe, mas \emph{agora} é absolutamente impossivel. Na verdade, para isso seria necessario que \textsc{Deus} se unisse á nossa intelligencia, --- ou por uma \emph{especie intelligivel}, --- ou pela sua \emph{essencia divina}. Ora \textsc{Deus} não póde unir-se á nossa intelligencia por uma \emph{especie} \pg{296}  \emph{intelligivel}; pois toda a \emph{especie intelligivel}, sendo creada e finita, não póde representar o Ente increado e infinito. Nem póde unir-se pela sua \emph{essencia divina}; pois que a \emph{essencia divina}, sendo purissima e infinita, não pode ser o objecto proporcionado de nenhuma intelligencia creada, muito menos da nossa intelligencia, que, embora espiritual, está comtudo unida á materia e \emph{directamente} só percebe o \emph{immaterial} existente na \emph{materia}. Logo o \emph{ontologismo} é falso \footnote{O argumento é convincente. Dá-se o conhecimento \emph{immediato} ou a visão de um objecto, quando a forma ou o principio, pelo qual o objecto é conhecido, é \emph{proprio} do mesmo objecto, de modo que elle, na ordem intencional e de cognoscibilidade, é perfeitamente semelhante ao objecto. Ora, quando se trata do conhecimento \emph{immediato} ou da \emph{visão} de Deus, tal principio não pode ser alguma coisa creada e distincta de Deus, porque a creatura tem apenas uma analogica semelhança com Deus. Logo deve ser a propria Essencia divina. Mas a Essencia divina não pode ser acto ou forma intelligivel da intelligencia creada, porque a potencia e o acto devem pertencer ao mesmo genero, e a Essencia divina nem \emph{genericamente} convem com a faculdade cognoscitiva creada. Ouçamos o S. Doutor: « Essentia divina non est \emph{naturalis} forma intelligibilis intellectus creati; quod sic patet. Actus enim et potentia semper sunt unius generis; unde forma naturalis intellectus non potest esse, nisi sit illius generis, in quo est potentia creati intellectus; unde forma sensibilis, quae est alterius generis, non potest esse forma ipsius (intellectus creati), sed forma immaterialis tantum, quae est generis sui. Sicut autem forma sensibilis est infra genus intellectivae potentiae creatae, ita Essentia divina est supra ipsum; unde Essentia divina non est forma, ad quam se extendit naturalis facultas intellectus creati » (\emph{De ver}., q. 8, a. III). --- Por isso a intelligencia creada, para que possa unir-se immediatamente á Essencia divina, deve ser revestida e confortada pela luz sobrenatural da gloria. --- Os \emph{ontologistas} destroem a differença essencial entre a ordem natural e a sobrenatural. Por quanto, a differença entre a visão beatifica e a visão natural de Deus, apregoada por esses escriptores, não seria essencial, mas apenas \emph{accidental}, porque a beatifica seria apenas mais viva e intensa que a natural, mas o mais e o menos não mudam a especie. E neste caso a nossa alma não precisaria da luz da gloria para ver a Essencia divina: o que é contrario á doutrina da Egreja, que, no Concilio de Vienna, condemnou o seguinte erro dos Beguardos e Beguinas: \emph{Anima non indiget lumine gloriae}, \emph{ipsam elevante}, \emph{ad videndum Deum}.}.

c) O \emph{ontologismo} admitte ou deve admittir tres pontos: --- que o nosso espirito, desde a sua creação, está \emph{continuamente} pensando, pois que a visão \emph{intuitiva} de \textsc{Deus} seria \emph{primordial} e \emph{continua}; --- que o objecto \emph{proporcionado} da nossa intelligencia não é o \emph{immaterial na materia}, mas o \emph{immaterial puro}, pois \textsc{Deus} é purissimo espirito; --- que a mesma intelligencia, no exercicio da sua actividade, não depende, de nenhum modo, dos sentidos, pois que \textsc{Deus} não está ao alcance dos sentidos.

Ora estes tres pontos, como demonstrámos, são falsos. \textbf{Logo} o \emph{ontologismo} é falso \footnote{Alem d'isso, se o homem tivesse a intuição \emph{directa} e \emph{immediata} da Essencia divina, não só seria perfeitamente feliz, mas nem poderia commetter nenhum erro ou ter alguma duvida no conhecimento de Deus e dos seus attributos. Ora estes consequentes são claramente falsos. Diz S. Thomaz: « Quidam dixerunt, quod primum, quod a mente humana cognoscitur etiam in hac vita, est Deus, qui est veritas prima et per hunc omnia alia cognoscuntur. Sed hoc aperte est falsum; quia cognoscere Deum per essentiam est hominis beatitudo; unde sequeretur omnem hominem beatum esse. Et praeterea cum in divina essentia omnia quae dicuntur de ipsa sint unum, nullus erraret circa ea, quae dicuntur de Deo; quod experimento patet esse falsum. Et iterum, ea quae sunt prima in cognitione intellectus, oportet esse certissima, unde intellectus certus est se ea intelligere, quod patet in proposito non esse » (In Boet. \emph{De Trin.}, q. I, a. 3).\par Não admira, pois, que a Egreja, guarda e mestra infallivel da verdade, condemnasse, pela Congregação do S. Officio (12 Set. 1861), sete proposições, que encerram o \emph{ontologismo}, das quaes a primeira é a seguinte: « Immediata Dei cognitio, habitualis saltem, intellectui humano essentialis est, ita ut sine ea nihil cognoscere possit: siquidem est ipsum lumen intellectuale ».\par São muitas e diversas as difficuldades, oppostas pelos \emph{ontologistas}. Apresentamos, para a solução das principaes, alguns principios. --- 1.º) Todas as coisas creadas recebem de Deus o ser e a intelligibilidade. Mas, assim como o ser das coisas, embora derive de Deus, não é constituido \emph{formalmente} pelo proprio ser de Deus, assim tambem a intelligibilidade das coisas, embora seja uma participação da intelligibilidade divina, todavia não é constituida formal e intrinsecamente por ella. Por isso Deus não é causa \emph{formal} do nosso conhecimento, mas é somente causa \emph{efficiente}, porque faz com que as coisas sejam intelligiveis, e com que a alma tenha a faculdade de entender e actualmente entenda, sendo Elle a causa primeira e universal de todo o ser. Por isso a luz intellectual, de que é revestida a nossa alma, não é Deus, mas deriva de Deus (Cf. \emph{Sum}. \emph{Th}., p. I, q. 12, a. 2). --- 2.º) Quando se diz que todas as coisas são por nós conhecidas na primeira Verdade, e que é por meio d'esta que de tudo se julga, não se deve entender que a Verdade increada é o principio pelo qual entendemos e julgamos, mas que pela luz, que Deus nos infundiu e que é uma semelhança da sua luz infinita, conhecemos e julgamos de todas as coisas. --- 3.º) Deus está presente á nossa intelligencia --- não como objecto ou forma intelligivel, --- mas como principio efficiente e causa conservadora. --- 4.º) Deus é o primeiro intelligivel em si mesmo e em relação á intelligencia divina, mas não em relação á intelligencia creada. O ser divino está fóra do genero de toda a intelligencia finita, e por isso é proporcionado unicamente para a intelligencia divina. --- 5.º) O ente infinito não pode ser conhecido pelas especies, abstrahidas dos entes finitos, de um modo \emph{adequado} e por um conhecimento \emph{proprio}, mas pode ser conhecido de um modo \emph{inadequado} e por um conhecimento \emph{analogico}. --- 6.º) A ordem logica ou do conhecimento deve corresponder á ordem ontologica ou da realidade --- emquanto a intelligencia deve conhecer as coisas como são em si mesmas, de modo que o que se affirma das coisas se encontre verdadeiramente nellas, --- e não emquanto a ordem dos conhecimentos\pgn{298}deva seguir a ordem das existencias, de modo que seja conhecido primeiramente o que primeiramente existe (se assim fosse, nunca poderiamos subir dos effeitos á causa, pois o effeito é naturalmente posterior á causa).\par Façamos agora uma breve critica do \emph{ontologismo} de \emph{Mallebranche}, de \emph{Gioberti} e de \emph{Rosmini}.\par a) O \emph{ontologismo} de \emph{Mallebranche} é falso, não só pelas razões apresentadas no texto, mas tambem pelas seguintes:\par 1.ª) \emph{Baséa-se num falso fundamento}. Por quanto suppõe que se não possa perceber o Ente supremo por uma idea, que seja distincta d'Elle. Ora isto é falso; pois, se não podemos ter uma idea, que seja distincta de \textsc{Deus} e ao mesmo tempo \emph{adequada}, podemos ter e effectivamente temos uma idea, que é distincta de \textsc{Deus}, embora \emph{inadequada} e infinitamente distante do objecto, que ella representa.\par 2.ª) \emph{Mallebranche}, como todos os \emph{Ontologistas}, confundem o Ente supremo e infinito com o ente commum e abstracto. Ora esta confusão leva ao \emph{pantheismo}. Com effeito, o ente commum e abstracto é um attributo essencial das creaturas, pois que dizemos: o homem é ente, o sol é ente, etc. Logo, se o ente commum e abstracto é o proprio \textsc{Deus}, \textsc{Deus} será o attributo essencial das creaturas, e teremos inevitavelmente o \emph{pantheismo}.\par b) O \emph{ontologismo} de \emph{Gioberti} é mais absurdo que o de \emph{Mallebranche}.\par 1.º) \emph{Gioberti} diz que nós, pela \emph{reflexão}, percebemos clara e distinctamente o que obscura e confusamente tinhamos percebido pela \emph{intuição}. Ora, pela \emph{intuição}, não só teriamos percebido o proprio \textsc{Deus}, mas tambem o mundo, e o acto, com que \textsc{Deus} o cria. Logo, pela \emph{reflexão}, deveriamos perceber conjunctamente todos estes tres objectos. E todavia não é assim; pois a consciencia attesta que, se percebemos o mundo, não percebemos o proprio \textsc{Deus} actualmente creador.\par 2.º) Une o \emph{tradicionalismo} ao \emph{ontologismo}, porque admitte a necessidade da linguagem para o primeiro acto da reflexão; quando é certo que a intelligencia humana, embora dependa das faculdades sensitivas, todavia póde encetar e executar as suas operações, sem precisar de algum auxilio exterior. --- Nem se diga que a palavra serve para determinar ou limitar o objecto da reflexão. Por quanto esse objecto ou é infinito, ou finito. Se é infinito, não póde nem deve ser determinado ou limitado. Se é finito, já não precisa ser determinado ou limitado.\par 3.º) Este systema leva ao \emph{pantheismo}. Na verdade, \emph{Gioberti} diz que as substancias creadas são destituidas de toda a intelligibilidade. Se as substancias creadas são destituidas de toda a intelligibilidade, são tambem destituidas de toda a realidade; pois aquella é uma consequencia necessaria d'esta. Se as substancias creadas são destituidas de toda a realidade, haverá uma unica substancia, que é a increada. Ora isto é \emph{pantheismo}. --- A intelligibilidade das coisas, e por isso a sua realidade, segundo \emph{Gioberti}, é a Idea, é \textsc{Deus}. Diz assim: « \emph{A Idea}, sendo a propria \emph{intelligibilidade das coisas, irradia a si mesma e o mundo com o proprio esplendor} ». (\emph{Introd}., pag. 8).\par c) O systema de \emph{Rosmini} é inacceitavel, não só porque contém o \emph{ontologismo}, mas tambem porque leva ao \emph{pantheismo} e ao \emph{idealismo}.\par 1.º) Contém o \emph{ontologismo}. Por quanto \emph{Rosmini} ensina que o ser \emph{ideal} ou\pgn{299}\emph{virtual} se manifesta immediatamente á nossa intelligencia. Ora este ser \emph{ideal} ou \emph{virtual}, segundo o mesmo escriptor, é o \emph{proprio Deus}, é o Verbo \emph{divino}. --- E' verdade que \emph{Rosmini} ensina que o ser \emph{ideal} ou \emph{virtual}, por não exprimir a subsistencia divina, não deve nem póde ser chamado \emph{Deus}. Mas esta emenda é inutil para o escopo. Por quanto, a razão, adduzida por \emph{Rosmini}, seria valida, se o ser \emph{ideal} ou \emph{virtual} fosse apenas uma noção abstracta, isto é, formada pela contemplação das creaturas; porque então poderiamos considerar uma propriedade divina, sem considerar as outras. Mas não é assim. Aquelle ser \emph{divino}, que intuitiva e continuamente se manifestaria á nossa intelligencia, não é, na opinião de \emph{Rosmini}, uma noção abstracta, formada pela contemplação das creaturas, mas é uma coisa \emph{concreta}, isto é, percebe-se como está \emph{em si} (\emph{Teosofia}, V. II, pag. 165). Ora tudo o que é \emph{divino}, como está \emph{em si}, é o proprio \textsc{Deus}. Logo é o proprio \textsc{Deus} que se manifesta á intelligencia, quando ella contempla o \emph{ser inicial} ou \emph{virtual}.\par 2.º) \emph{Leva ao pantheismo}. Para \emph{Rosmini}, o \emph{ser inicial}, que é o proprio \emph{ser} do Verbo divino (exceptuada a existencia), é o \emph{ser} de que toda e cada uma das creaturas é essencialmente constituida; e por isso o ente finito é um termo ou forma impropria do proprio ser; o ser das creaturas não é \emph{feito}, mas \emph{posto}; e, quando a creatura perece, o ser não \emph{perece}, \emph{retira-se} da propria creatura (\emph{Teosofia}, vol. I, pag. 269 e 350). Se o ser das creaturas não é \emph{feito} ou \emph{creado}, o \emph{pantheismo} é inevitavel.\par 3.º) \emph{Leva ao idealismo}. \emph{Rosmini} diz que percebemos os objectos externos, applicando o \emph{ser inicial} ás sensações. Mas, se reflectirmos que, segundo \emph{Rosmini}, as nossas sensações não são objectivas, mas inteiramente subjectivas, e que o \emph{ser ideal} ou \emph{inicial}, que se manifesta á intelligencia, não deriva das coisas creadas mas do proprio \textsc{Deus}, deveremos concluir que o nosso conhecimento não é \emph{objectivo} e \emph{real}, mas é um phantasma vão e esteril.\par Neste ponto \emph{Rosmini} concorda com \emph{Kant}, com a differença que este admitte a necessidade de muitas formas innatas, produzidas pelo nosso proprio espirito, ao passo que o escriptor italiano admitte uma só idea innata, proveniente da \emph{immediata} intuição de \textsc{Deus}.\par O systema de Rosmini foi condemnado pela Egreja, com o decreto da Suprema Congregação do S. Officio do dia 14 de Dezembro de 1887. --- E' a proposição 37ª que contém o \emph{ontologismo}.}.

186. \textbf{Sensualismo}. --- O \emph{sensualismo} ensina que \emph{todos} os nossos conhecimentos, não só os sensitivos mas tambem os \pg{299}  intellectuaes, derivam unica e exclusivamente das nossas faculdades \emph{sensitivas} \footnote{O \emph{sensualismo} foi seguido --- na antiguidade por \emph{Leucippo}, \emph{Democrito}, \emph{Epicuro}, --- e nos tempos modernos por \emph{Locke}, \emph{Condillac}, \emph{Laromiguière}, \emph{Comte}, \emph{Mill}, \emph{Taine}, \emph{Bain}, \emph{H. Spencer}, e por todos os \emph{materialistas} e \emph{positivistas}.\par a) \emph{Leucippo}, \emph{Democrito}, \emph{Epicuro}, querendo explicar a origem dos nossos conhecimentos intellectuaes, disseram que dos corpos se destacam umas imagens ou superficies subtilissimas, as quaes, atravessando o ar, penetram no cerebro e na alma.\par b) \emph{Locke}, escriptor inglez (1632-1704), seguindo as pizadas do francez \emph{P. Gassendi} (1592-1655), e do inglez \emph{T. Hobbes} (1588-1679), ensinou a seguinte theoria: A causa unica de todas as nossas ideas é constituida pela \emph{sensação} e pela \emph{reflexão}. Pela \emph{sensação} percebemos os \emph{objectos externos}, que chegam á nossa alma pelos orgãos dos sentidos externos, e que são as \emph{côres}, os \emph{sabores}, etc. Pela \emph{reflexão} percebemos os objectos da sensibilidade interna, que\pgn{300}são as \emph{operações}, com que a alma percebe os objectos externos. Emquanto reflecte sobre as suas operações, a alma forma novas ideas, diversas das ideas dos objectos externos, como são as ideas de \emph{percepção}, de \emph{duvida}, de \emph{pensamento}, de \emph{volição}, etc. E', pois, pela \emph{sensação} e pela \emph{reflexão} que adquirimos todos os conhecimentos.\par c) \emph{Condillac}, francez (1715-1780), acceitou os principios de \emph{Locke}, mas levou o \emph{sensualismo} ao exaggero. Principiou por dizer que o homem é mera machina; em seguida, gradual e successivamente, attribuiu-lhe este ou aquelle sentido, esta ou aquella sensação; finalmente concluiu que todas as operações da alma são apenas diversos aspectos ou diversas \emph{transformações} da sensação. --- Descendo ao particular, \emph{Condillac} enumerou seis operações ou faculdades da alma, a saber: a \emph{attenção}, a \emph{comparação}, o \emph{juizo}, a \emph{reflexão}, a \emph{imaginação} e o \emph{raciocinio}. A \emph{attenção} é a mais viva e intensa das sensações determinadas pelos objectos externos. A \emph{comparação} é uma duplice attenção a dois objectos. O \emph{juizo} é a juxtaposição dos objectos da comparação, para se conhecer a sua conveniencia ou desconveniencia. A \emph{reflexão} é uma attenção mais intensa. A \emph{imaginação} é a synthese das qualidades percebidas pelos sentidos. O \emph{raciocinio} é a deducção de um juizo de outro. --- Na opinião, pois, de \emph{Condillac}, o conhecimento é uma \emph{sensação transformada}.\par d) \emph{Laromiguière} (1756-1837) procurou mitigar a doutrina de \emph{Condillac}, dizendo o seguinte: Além da \emph{sensação}, deve admittir-se uma força activa, que opéra nas proprias sensações. Esta actividade divide-se em tres faculdades, que são a \emph{attenção}, a \emph{comparação} e o \emph{raciocinio}. A intelligencia humana, pois, comprehende tres faculdades. Pela \emph{attenção}, que concentra a actividade, --- pela \emph{comparação}, que divide a mesma actividade e que é uma duplice attenção, --- pelo \emph{raciocinio}, que ainda a divide e que é uma duplice comparação, o espirito torna-se uma faculdade (\emph{Leçons de philos.}). --- D'onde se vê que, na opinião d'este escriptor, a \emph{attenção}, operando nas sensações, é unica causa da origem das ideas. E' uma reproducção do systema de \emph{Locke}.\par e) \emph{A. Comte} (1798-1857) renovou o \emph{sensualismo}, inventando uma theoria, a que deu o nome de \emph{positivismo}. --- Digamos alguma coisa ácerca da \emph{natureza} e das \emph{bases} d'este systema, que tantos espiritos tem illudido.\par a) \emph{Positivismo}. --- O \emph{positivismo} admitte que o espirito humano é capaz de attingir as verdades \emph{positivas} ou da ordem experimental, mas não é capaz de resolver as questões \emph{metaphysicas}, que não podem ser verificadas pela observação ou pela experiencia \emph{sensivel}.\par O \emph{positivismo}, pois, admittindo, como certo, só o que está sujeito aos methodos da observação e da experiencia \emph{sensivel}, não conhece senão os \emph{factos} e as \emph{leis}, que os regulam e que não passam de \emph{factos} geraes. A \emph{origem} e a \emph{essencia} dos seres, como tambem as \emph{causas supremas} e \emph{finaes}, são coisas desconhecidas e impenetraveis. --- Por isso póde dizer-se que o \emph{positivismo} é\pgn{301}um \emph{dogmatismo physico} e um \emph{scepticismo metaphysico}. E' um \emph{dogmatismo physico}; pois que affirma o objectividade do mundo physico. E' um \emph{scepticismo metaphysico}; porque não quer pronunciar-se ácerca da existencia e da natureza dos objectos metaphysicos.\par b) \emph{Bases do positivismo}. --- \emph{A. Comte} funda o seu systema sobre duas \emph{bases}, que são a \emph{lei dos tres estados} e a \emph{classificação das sciencias}.\par c) \emph{Lei dos tres estados}. --- A \emph{lei dos tres estados}, inventada para expulsar a \emph{metaphysica} fora dos limites do saber humano, representa, segundo \emph{Comte}, a historia da evolução scientifica.\par Diz \emph{Comte} que o espirito humano, no seu desenvolvimento, passou successivamente por tres estados, que são: 1.º) o estado \emph{theologico} ou \emph{ficticio}, que tem differentes phases (\emph{fetichismo}, \emph{polytheismo} e \emph{monotheismo}), e em que o espirito humano explica os phenomenos por meio de vontades transcendentes ou de agentes sobrenaturaes; --- 2.º) o estado \emph{metaphysico} ou \emph{abstracto}, em que o espirito humano explica os phenomenos por meio de forças ou entidades occultas e abstractas, como o \emph{principio vital}, etc.; --- 3.º) o estado \emph{positivo} ou \emph{scientifico}, em que o espirito humano explica os phenomenos, subordinando-os a leis experimentalmente demonstradas. Todas as sciencias passaram pelos dois primeiros estados, e só se constituiram, quando chegaram ao terceiro. O estado \emph{positivo}, pois, é o termo fixo e definitivo, em que o espirito humano descança e encontra a \emph{sciencia}.\par d) \emph{Classificação das sciencias}. --- Partindo do principio de que o objecto da sciencia é só o \emph{positivo}, isto é, só o que póde estar sujeito ao methodo da observação e da experiencia sensivel, \emph{A. Comte} não reconhece senão as sciencias \emph{experimentaes} ou positivas, que tratam dos \emph{factos} e das suas \emph{leis}, e que, \emph{dispostas na sua ordem hierarchica}, são as seguintes: a \emph{mathematica}, a \emph{astronomia}, a \emph{physica}, a \emph{chimica}, a \emph{biologia}, e a \emph{sociologia}.\par Esta theoria, cujos vestigios se encontram principalmente em \emph{Turgot}, \emph{Saint-Simon}, acha-se exposta e largamente desenvolvida no livro de \emph{Comte}, intitulado: \emph{Cours de Philosophie Positive}.\par f) \emph{Stuart Mill} (\emph{Philosophie de Hamilton}), \emph{A. Bain} (\emph{Les Sens et l'Intelligence}), e \emph{H. Taine} (\emph{De l'Intelligence}), deram origem a um novo systema, chamado \emph{associacionismo}. --- Eis o resumo d'esta theoria:\par As nossas ideas não passam de sensações \emph{associadas}.\par Basta que dois factos ou duas ideas sejam uma vez percebidas simultaneamente, para que os dois phenomenos adquiram a \emph{tendencia} para se reproduzir ao mesmo tempo. Se, por qualquer circumstancia, se reproduzir bastantes vezes a mesma associação, aquella tendencia transforma-se em \emph{habito}; e este, fortalecido pela experiencia, é \emph{irresistivel}.\par Ora os dois factos ou as duas ideas, assim \emph{associadas}, dão origem a certos \emph{juizos}, que são \emph{universaes} e \emph{necessarios}; --- \emph{necessarios}, porque é impossivel\pgn{302}separar os dois elementos, que os constituem; --- \emph{universaes}, porque as associações, em que se fundam, são communs a todos os homens.\par Tal é a origem de todos os \emph{principios da razão}, nos quaes se baseam todas as sciencias, tanto \emph{experimentaes} como \emph{racionaes}.\par As sciencias \emph{experimentaes} ou \emph{positivas} baseam-se sobre a \emph{inducção}, que é uma \emph{invencivel expectativa dos mesmos consequentes depois dos mesmos antecedentes}. Mas esta \emph{expectativa} é um habito resultante da experiencia, a qual attesta que \emph{ao mesmo antecedente succede sempre o mesmo consequente}. E' o pretendido \emph{principio de causalidade}.\par As sciencias \emph{mathematicas} ou \emph{racionaes} derivam tambem da experiencia; pois todos os \emph{axiomas} não passam de verdades \emph{experimentaes}. Assim o \emph{principio da contradicção} é um habito contrahido pela experiencia, que muitas vezes attestou que a \emph{affirmação e a negação são dois actos do espirito, que se exluem naturalmente}.\par g) \emph{Herbert Spencer} (\emph{Principes de Psychologie}), seguido por \emph{Lewes}, \emph{Murphy} e pela eschola \emph{evolucionista}, procurou remediar aos inconvenientes do \emph{associacionismo} por meio da theoria da \emph{hereditariedade}, dizendo que os \emph{principios da razão} não passam de \emph{habitos hereditarios}. --- D'ahi o nome de \emph{hereditarismo}, dado ao systema de \emph{Herbert Spencer}.\par O escriptor inglez expõe assim a sua theoria: O \emph{habito} é o principio que explica tudo. Mas o \emph{habito} não fica \emph{individual}. Toda a percepção, toda a \emph{associação}, fortalecida pelo \emph{habito}, tende para se tornar \emph{organica}, isto é, transmissivel pela \emph{hereditariedade}. D'este modo o \emph{habito} do pae torna-se natureza do filho. Cada geração accrescenta alguma coisa aos habitos recebidos das gerações anteriores. E, como as gerações successivas obedecem á \emph{lei suprema da evolução}, os \emph{habitos} tornam-se constantes e immutaveis, e constituem as \emph{formas do pensamento}, os \emph{principios da nossa razão}.\par Como se vê, o \emph{hereditarismo} e o \emph{associacionismo} em parte concordam, e em parte differem entre si. --- Concordam em que ambos reconhecem a experiencia, como fonte \emph{unica} dos nossos conhecimentos. --- Differem em que --- no \emph{associacionismo} a experiencia, que se diz a fonte dos nossos conhecimentos, é \emph{individual}, --- ao passo que no \emph{hereditarismo} a experiencia é \emph{collectiva}, pois que é de \emph{todas as gerações passadas}.}.

187. \textbf{O sensualismo é falso.}

a) Se a causa necessaria e sufficiente de \emph{todos} os nossos conhecimentos fosse constituida pelos \emph{sentidos}, não poderiamos ter noções de coisas \emph{abstractas} e \emph{immateriaes}; pois que pelos sentidos só percebemos o \emph{concreto} e o \emph{singular}. Ora possuimos noções de coisas \emph{abstractas} e \emph{immateriaes}, como são as noções de \emph{virtude}, de \emph{justiça}, de \textsc{Deus}, etc. Logo a causa necessaria \pg{301}  e sufficiente de \emph{todos} os nossos conhecimentos não é constituida pelos \emph{sentidos}.

b) O \emph{sensualismo} admitte que \emph{todas} as nossas faculdades são sensitivas e organicas. Ora esta opinião é absurda; visto que, como demonstrámos, o homem é dotado de intelligencia, e esta é uma faculdade \emph{espiritual}; pois que os objectos, que ella percebe, são espirituaes. Logo o \emph{sensualismo} é falso.

c) O \emph{sensualismo} destroe o fundamento de todas as sciencias. Na verdade, todas as sciencias basêam-se em principios universaes, absolutos e immutaveis, de que deduzem conclusões necessarias e independentes das coisas contingentes e mutaveis. Ora, se todos os nossos conhecimentos derivassem unica e exclusivamente dos \emph{sentidos}, as sciencias fundar-se-hiam sobre bases mutaveis e contingentes; porque mutaveis e contingentes são os conhecimentos \emph{sensitivos}. Logo o \emph{sensualismo} é falso \footnote{Refutemos agora as diversas especies do \emph{sensualismo}.\par a) O \emph{sensualismo} de \emph{Leucippo}, de \emph{Democrito} e de \emph{Epicuro} é falso.\par 1.º) Este systema confunde a intelligencia com os sentidos; quando é certo que aquella, por ser espiritual e immaterial, é distincta e diversa d'estes, que são organicos e materiaes.\par 2.º) Repugna que dos corpos se destaquem certas imagens ou superficies subtilissimas, para, atravez do ar, penetrarem no cerebro e na alma. Na verdade\pgn{303}, essas imagens ou superficies são materiaes e accidentaes. Ora uma coisa material não póde penetrar num outro objecto material (pois a materia é naturalmente impenetravel), nem uma coisa accidental póde separar-se do sujeito que a sustenta.\par 3.º) Se fosse verdadeiro tal systema, todos os nossos conhecimentos deveriam estar encerrados na esphera do sensivel e do concreto. Mas não é assim; porque a nossa intelligencia forma ideas espirituaes e abstractas.\par b) O systema de \emph{Locke} é absurdo pelas seguintes razões.\par 1.º) Confunde a intelligencia com os sentidos, e por consequencia nega a espiritualidade da nossa alma, --- confusão e negação que são absurdas.\par 2.º) A nossa intelligencia percebe objectos abstractos e immateriaes. Mas não os poderia perceber, se fosse verdadeiro o systema de \emph{Locke}. Porquanto, segundo esse escriptor, a \emph{sensação} e a \emph{reflexão} são a causa dos nossos conhecimentos. Ora, pela \emph{sensação} e \emph{reflexão} não podemos perceber objectos abstractos e immateriaes; pois pela \emph{sensação} só percebemos objectos concretos e materiaes, e pela \emph{reflexão} percebemos de um modo mais distincto os mesmos objectos, que foram percebidos de um modo confuso pela \emph{sensação}, e que são concretos e materiaes. Logo o systema de \emph{Locke} é absurdo.\par c) Não menos absurdo é o \emph{sensualismo} de \emph{Condillac}. Na verdade:\par 1.º) Este escriptor, não reconhecendo em nós outras faculdades senão as sensitivas, rebaixa o homem á condição dos animaes irracionaes.\par 2.º) \emph{Condillac} reduz todas as forças á \emph{sensibilidade}, e depois diz que a sensibilidade se transforma em faculdades. Ora isto é contradictorio. Aquella transformação deve necessariamente suppor uma faculdade ou uma força, que opére na sensibilidade para a transformar. Essa força não póde ser a propria sensibilidade, porque o artista distingue-se da materia em que opéra. Logo devemos dizer --- ou que a sensibilidade não se transforma em faculdades, --- ou que nem todas as forças se reduzem á sensibilidade.\par 3.º) As faculdades, que \emph{Condillac} enumera, não podem explicar-se unicamente pela sensação. Na verdade, esse escriptor diz que a \emph{attenção} é a sensação de um objecto, que nos impressiona. Ora isto é falso; não só porque muitas vezes a nossa \emph{attenção} se fixa em objectos meramente espirituaes, que não impressionam as nossas faculdades sensitivas, --- mas tambem porque muitas vezes, deante de uma vivissima sensação, produzida por um objecto, distrahimos a nossa attenção para um objecto diverso e até contrario. --- O \emph{juizo} tambem não póde ser explicado pela sensação. Por quanto, a intelligencia não póde emittir um juizo, sem conhecer a relação entre o sujeito e o predicado. Ora esta relação, por ser uma coisa abstracta, não póde ser percebida pelas faculdades sensitivas, que só percebem objectos concretos. --- Se a \emph{attenção} e o \emph{juízo} não podem ser explicados pela sensação, muito menos o poderá ser o \emph{raciocinio}. Com effeito, quando \emph{raciocinamos}, estabelecemos, como principio e base, uma verdade universal, de que deduzimos uma conclusão particular. Ora a verdade universal não póde ser percebida pelos sentidos, mas só por uma faculdade inorganica ou espiritual. --- Diga-se o mesmo com relação ás outras faculdades, enumeradas por \emph{Condillac}.\par d) O systema de \emph{Laromiguiére} é falso.\par 1.º) Se todas as nossas faculdades se reduzissem á \emph{attenção}, como anterior\pgn{304}a esta é só a \emph{sensação}, todos os nossos conhecimentos seriam sensitivos; pois que pela \emph{attenção} não percebemos novos objectos, mas reflectimos sobre os objectos anteriormente percebidos.\par 2.º) Sendo o systema de \emph{Laromiguiére} uma reproducção do de \emph{Locke} (com a unica differença que o escriptor francez attribue á \emph{attenção} tudo o que o escriptor inglez attribue á \emph{reflexão}), aquelle refuta-se com os mesmos argumentos com que foi refutado este.\par e) O \emph{positivismo} de \emph{Comte} é falso na sua \emph{natureza} e nas suas \emph{bases}.\par \emph{a}) E' falso na sua natureza pelas seguintes razões:\par 1.ª) Este systema, só reconhecendo como certo o que está ao alcance dos sentidos, admitte o \emph{sensualismo} com todos os antecedentes e consequentes. Ora o \emph{sensualismo}, como demonstramos no texto, é falso.\par 2.ª) Suppõe que a experiencia attesta simplesmente os \emph{factos}, e não leva até ao descobrimento da \emph{causa}. Ora esta supposição é falsa; porque, se a experiencia \emph{vulgar} só se limita a attestar os \emph{factos}, a experiencia \emph{scientifica} leva até ao descobrimento da \emph{causa} d'esses factos.\par 3.ª) Encerra affirmações \emph{gratuitas}; pois não adduz razões para provar --- que o nosso espirito só póde alcançar o \emph{relativo}, --- que os \emph{sentidos} são os \emph{unicos} meios para a acquisição da certeza, --- que o testemunho da razão e da \emph{consciencia} é falso, --- que as \emph{causas} são inaccessiveis ao nosso espirito, --- e que a idea de \emph{causa} é absurda.\par 4.ª) Encerra muitas contradicções; por quanto --- affirma que não possuimos principios racionaes e abstractos, e depois emprega a \emph{inducção}, que, em ultima analyse, se basêa naquelles principios; --- rejeita as noções transcendentes, e depois admitte a \emph{certeza} dos factos e das leis, \emph{certeza} que não póde alcançar-se sem aquellas noções; --- não admitte o \emph{principio} de \emph{causalidade}, e todavia dos phenomenos sobe ás leis geraes, o que só póde fazer-se em virtude d'aquelle principio; --- nega o \emph{absoluto}, e depois sustenta a proposição de que só é \emph{certo} o que foi verificado pela experiencia sensivel, proposição que encerra uma coisa \emph{absoluta}; --- não admitte questões ácerca da \emph{origem} dos seres, e depois quer explicar a \emph{origem} do mundo, da vida, do homem; --- protesta a sua \emph{neutralidade} entre o \emph{espiritualismo} e o \emph{materialismo}, e depois admitte que a materia é \emph{eterna}, que o pensamento é uma funcção do cerebro, que a \emph{psychologia} é reductivel á \emph{biologia}; --- affirma que não nega a existencia de \textsc{Deus} e da nossa alma espiritual, e depois ensina que \textsc{Deus} e o espirito são entidades \emph{ficticias} e \emph{imaginarias}. O proprio \emph{Huxley} diz: « A philosophia positivista é um tecido de contradicções, um cumulo de absurdos ».\par \emph{b}) E' falso na \emph{lei dos tres estados}. --- A. \emph{Comte} diz que os tres estados --- \emph{theologico}, \emph{metaphysico} e \emph{positivo} --- por que passa o espirito humano, são \emph{successivos}, e o posterior exclue o anterior. Ora, se esta affirmação fosse verdadeira, deveriamos encontrar, em toda a humanidade em geral, e em cada povo e em cada individuo em particular, a applicação d'aquella lei; de modo que, tanto na humanidade inteira, como em cada povo e em cada individuo, primeiramente deveria dominar o espirito \emph{theologico}, depois o \emph{metaphysico}, e finalmente o \emph{positivo}. Mas a propria historia, para a qual A. \emph{Comte} appella, demonstra que não foi assim.\par Se consideramos a evolução do espirito da humanidade inteira, desde a\pgn{305}mais remota antiguidade até aos nossos dias, veremos que os \emph{tres estados} não dominam \emph{successivamente}, segundo a doutrina de \emph{Comte}, mas que na \emph{antiguidade} e \emph{nos tempos modernos} dominam o espirito \emph{metaphysico} e o \emph{positivo}, e na \emph{edade media} domina o espirito \emph{theologico}.\par Estudando os \emph{tres} principaes periodos da philosophia grega, concluimos que os \emph{tres estados} se succedem numa ordem diametralmente opposta á que foi indicada por \emph{Comte}. Na verdade, os primeiros philosophos, como \emph{Thales}, \emph{Anaximenes}, \emph{Heraclito}, \emph{Democrito}, eram principalmente \emph{positivistas}, porque procuravam a explicação dos phenomenos do mundo no proprio mundo; --- os philosophos, que succederam a estes e que se chamaram \emph{Platão} e \emph{Aristoteles}, foram sobretudo \emph{metaphysicos}, porque dos phenomenos passageiros remontavam ás essencias immutaveis; --- e finalmente, os philosophos da eschola de \emph{Alexandria}, os \emph{Gnosticos} e os Padres da Egreja foram principalmente \emph{theologos}, visto que reconheciam a vontade divina como causa efficiente ou permissiva de todos os acontecimentos humanos.\par A historia dos seculos XVII e XVIII ensina que naquelle dominavam o espirito \emph{metaphysico} e o \emph{theologico}, sem excluir o \emph{positivo}, ao passo que neste dominava o espirito \emph{positivo}, mas sem excluir o \emph{metaphysico}, que se desenvolveu principalmente na Allemanha.\par Se, finalmente, considerarmos a evolução intellectual do individuo, veremos que o espirito \emph{theologico}, \emph{metaphysico} e \emph{positivo} se impõem e dominam ao mesmo tempo. Na verdade, a nossa inteligencia não se satisfaz com a descoberta da \emph{lei} de um phenomeno; quer. saber tambem a \emph{origem} e o \emph{fim} d'esse phenomeno. Se ás questões ácerca da \emph{lei} de um Phenomeno responde a sciencia \emph{positiva}, ás questões ácerca da \emph{origem} e do \emph{fim} do mesmo phenomeno respondem a sciencia \emph{theologica} e a \emph{metaphysica}. E' por isso que os mais celebrados talentos, como \emph{Leibnitz}, \emph{Newton}, Ampére, \emph{Cauchy}, depois de se terem elevado ao mais alto grau do espirito \emph{scientifico} ou \emph{positivo}, se abandonavam ás mais sublimes especulações \emph{theologicas} e \emph{metaphysicas}, e, pelo facto de serem bons \emph{christãos} e profundos \emph{metaphysicos}, não deixaram de ser \emph{sabios naturalistas}.\par Aqui advertimos que, censurando a \emph{lei dos tres estados}, não queremos dizer que os tres modos de pensar --- \emph{theologico}, \emph{metaphysico} e \emph{positivo} --- não são \emph{naturaes}; só queremos affirmar que aquelles tres modos de pensar, por isso mesmo que são \emph{naturaes}, não são \emph{successivos}, mas \emph{simultaneos}, pois o homem possue sempre toda a sua essencia. A \emph{religião}, a \emph{metaphysica} e a \emph{sciencia experimental} existiram sempre e simultaneamente, embora mais ou menos desenvolvidas, e sempre e simultaneamente hão de existir para o bem e o progresso dos individuos e da sociedade.\par Devemos advertir tambem que na refutação da \emph{lei dos tres estados}, que são o \emph{theologico}, o \emph{metaphysico} e o \emph{positivo}, demos a esses adjectivos --- \emph{theologico}, \emph{metaphysico} e \emph{positivo} --- o sentido em que os homens de criterio costumam tomal-os, dizendo que --- o estado \emph{theologico} é o que refere os phenomenos do mundo á vontade de \emph{Deus}, --- o estado \emph{metaphysico} é o que explica as operações intellectuaes por um principio \emph{espiritual}, e que investiga o \emph{fim} e a \emph{origem} das coisas, --- e o estado \emph{positivo} é o que explica os phenomenos\pgn{306}sensiveis pelas causas \emph{proximas}. Tomados neste sentido, os tres \emph{estados}, como dissemos, existem não \emph{successiva} mas \emph{simultaneamente}.\par Mas se os adjectivos --- \emph{theologico}, \emph{metaphysico} e \emph{positivo} --- se tomarem no sentido, em que os toma A. \emph{Comte}, dizendo que o estado \emph{theologico} é o que attribue a uma vontade \emph{transcendente} e \emph{ficticia} os conhecimentos humanos, --- que o estado \emph{metaphysico} é o que explica os pensamentos e as volições por entidades \emph{abstractas} e \emph{imaginarias}, --- e que o estado \emph{positivo} é o que se limita a observar os factos sem lhe importar a causa d'elles, então diremos que os \emph{tres estados} não existem nem nunca existiram. Por quanto, é falso --- que o espirito humano tenha começado e comece por attribuir os acontecimentos a uma vontade \emph{transcendente} e \emph{ficticia}, --- que em seguida tenha referido e refira os pensamentos e as volições a umas entidades \emph{abstractas} e \emph{imaginarias}, --- e que finalmente se tenha limitado e limite a observar os phenomenos, \emph{sem lhe importar a causa d'elles}.\par Nem se diga que o espirito humano passa por vezes de um erro para outro. O erro é anormal ao espirito, e tudo o que é anormal não constitue um \emph{estado}, muito menos uma \emph{lei}.\par \emph{c}) E' falso na \emph{classificação das sciencias}. Com effeito, a \emph{classificação}, apresentada por A. \emph{Comte}, além de ser falsa pelas razões expostas na \emph{Logica}, está baseada num falso supposto. Na verdade, excluindo do quadro das sciencias a \emph{metaphysica} e a \emph{theologia}, o fundador do \emph{positivismo} suppoz que a \emph{theologia} explica os acontecimentos por uma vontade \emph{transcendente} e \emph{ficticia}, e que a \emph{metaphysica} refere os pensamentos e as volições a um principio \emph{abstracto} e \emph{imaginario}. Ora este supposto é falso. A \emph{theologia} explica as verdades, de que trata, pelo infallivel testemunho de \textsc{Deus}, e \textsc{Deus} não é um ente \emph{ficticio}, mas eminentemente \emph{real}, que, embora incomprehensivel na sua essencia infinita, comtudo é cognoscivel e conhece-se nas suas obras; e a \emph{metaphysica} não attribue os pensamentos e as volições a um principio \emph{abstracto} ou \emph{imaginario}, mas a um principio \emph{concreto} e \emph{real}, que é a alma.\par f) O \emph{associacionismo} é falso pelas seguintes razões:\par 1.ª) Sendo este systema uma reproducção dos systemas de \emph{Locke} e de \emph{Condillac}, a sua falsidade demonstra-se --- não só pelos argumentos apresentados no texto contra o \emph{sensualismo} em geral, --- mas tambem pelos argumentos com que reprovámos os systemas d'aquelles dois escriptores.\par 2.ª) Segundo o \emph{associacionismo}, os \emph{principios da razão} percebem-se pelo \emph{habito}. Ora isto é falso. Por quanto, os \emph{principios da razão}, se se percebessem pelo \emph{habito}, deveriam adquirir \emph{gradual} e \emph{successivamente} a sua certeza invencivel, pois o \emph{habito} forma-se pela repetição dos actos. Mas não é assim. Os \emph{principios da razão}, logo que se manifestam á intelligencia, apparecem immediatamente dotados de certeza invencivel.\par 3.ª) Este systema não explica a \emph{universalidade}, nem a \emph{necessidade} dos \emph{principios da razão}. --- Não explica a \emph{universalidade}. Por quanto, as \emph{associações} devem variar segundo os \emph{habitos} de cada um; e se algumas associações são \emph{identicas} em todos os homens, essa \emph{identidade} não póde derivar de um principio \emph{subjectivo} e \emph{logico}, mas de um principio \emph{objectivo} e \emph{real}. --- Não explica a \emph{necessidade}. Com effeito, a \emph{necessidade} dos \emph{principios da razão} é \emph{absoluta} e \emph{real}, ao passo que o \emph{habito} de \emph{associar indissoluvelmente} duas\pgn{307}ideas é um facto meramente \emph{subjectivo}. Não é o \emph{habito} do nosso espirito que impõe a lei ás coisas; mas são as coisas que impõem e mudam os \emph{habitos} e por isso as associações das ideas.\par Não nos demoramos em refutar as opiniões dos \emph{associacionistas} acerca dos \emph{principios da razão} em geral e dos principios de \emph{contradicção} e de \emph{causalidade} em particular, pois já foram refutadas na \emph{Logica} e na \emph{Ontologia}.\par g) O \emph{hereditarismo} de \emph{H}. \emph{Spencer} é falso pelas seguintes razões.\par 1.ª) Basea-se num falso principio. \emph{Spencer} admitte a formação de novos typos, causada pela gradual accumulação de certos caracteres, transmittidos pela \emph{hereditariedade}. Ora este principio é falso; pois que, como provámos na \emph{Cosmologia}, as especies são \emph{immutaveis}.\par 2.ª) E' contraria aos factos. Com effeito, a historia attesta que os homens descobriram e applicaram, desde o principio da sua existencia e com a mesma exactidão, os principios que applicam hoje e que \emph{Spencer} diz terem sido formados, depois de muitas gerações, pelos \emph{habitos hereditarios}.\par 3.ª) Não explica a \emph{universalidade} nem a \emph{necessidade} dos \emph{principios da razão}. Porquanto, a experiencia, embora repetida innumeraveis vezes e durante muita gerações, só attesta \emph{factos mutaveis} e \emph{contingentes}, e por isso não póde produzir os \emph{principios da razão}, cuja verdade é \emph{immutavel} e \emph{necessaria}.\par 4.ª) Destroe a \emph{liberdade} do homem; pois que subordina a nossa vontade ao poder irresistivel dos habitos intellectuaes e moraes, recebidos dos antepassados. Ora é certo que os habitos, se inclinam a nossa vontade para um certo acto, não a determinam fatalmente.}.

188. \textbf{Idealismo}. --- O \emph{idealismo} ensina que todos os nossos conhecimentos derivam unica e exclusivamente da \emph{intelligencia}, \pg{304}  sem nenhuma dependencia das faculdades sensitivas e dos objectos externos; de modo que a intelligencia não conhece \pg{305}  os objectos, porque estes existem, mas os objectos existem, porque a intelligencia os conhece. --- Este systema chama-se \pg{306}  tambem \emph{transcendentalismo}; pois, querendo explicar a origem dos nossos conhecimentos, independentemente do exercicio das \pg{307}  faculdades sensitivas e dos objectos externos, \emph{transcende} ou excede toda a experiencia, interna e externa \footnote{O \emph{idealismo}, chamado tambem \emph{racionalismo transcendental} ou simplesmente \emph{transcendentalismo}, foi inventando por \emph{E}. \emph{Kant} e defendido por \emph{Fichte}, \emph{Schelling}, \emph{Kegel} e \emph{Cousin}.\par a) \emph{Kant} desenvolve a theoria \emph{idealista} no seu livro --- \emph{Critica da razão pura} --- de que damos o resumo.\par Principia por dizer que o nosso conhecimento se compõe de dois elementos, --- um \emph{extrinseco} que deriva do objecto externo, --- outro \emph{intrinseco} que resulta do sujeito pensante. O primeiro é a \emph{materia}, o segundo é a \emph{forma}.\par Não podemos explicar a origem e a natureza do conhecimento sem investigar a sua parte \emph{formal}, que em nós existe \emph{a-priori}. As notas, que caracterisam a parte \emph{formal} do conhecimento, são a \emph{universalidade} e a \emph{necessidade}.\par Em seguida, applicando estes principios á \emph{sensibilidade}, \emph{Kant} observa --- que os \emph{objectos} dos sentidos \emph{externos} se nos representam como situados no \emph{espaço}, e tão necessariamente que, se todo o mundo fosse reduzido ao nada,\pgn{308}ainda em nós ficaria a representação do \emph{espaço} incircumscripto ou illimitado, --- e que as \emph{sensações internas} se nos revelam numa successão tão continua, que, embora abstrahissemos d'ellas, todavia em nós ficaria a representação de successão ou de \emph{tempo}. --- O \emph{espaço}, pois, e o \emph{tempo} pertencem á parte \emph{formal} do conhecimento, e existem em nós \emph{a-priori}. Estas \emph{formas}, não dependendo da \emph{experiencia}, podem chamar-se \emph{intuições puras} do \emph{espaço} e do \emph{tempo}. \emph{A intuição pura} de \emph{espaço} é a forma da sensibilidade \emph{externa}; a de \emph{tempo} é a forma da sensibilidade \emph{interna}.\par Mais elevada que a da \emph{sensibilidade} é a actividade da \emph{intelligencia}; pois aquella limita-se a formar visões \emph{empiricas}, ao passo que esta produz \emph{pensamentos}. As \emph{formas} do \emph{pensamento} são tantas, quantas são as \emph{formas} do \emph{juizo}; pois pensar é julgar. Ora as formas do \emph{juizo} são \emph{doze}. Por quanto, em todos os juizos, seja qual for o seu objecto, encontramos sempre a \emph{quantidade}, a \emph{qualidade}, a \emph{relação} e o \emph{modo}. Ha, pois, quatro \emph{formas} existentes \emph{a-priori} na intelligencia: a \emph{quantidade}, a \emph{qualidade}, a \emph{relação}, e a \emph{modalidade}. Cada uma d'estas quatro \emph{formas} divide-se em tres classes. Por causa da \emph{quantidade}, o juizo é \emph{singular}, \emph{particular} ou \emph{universal}; --- por causa da \emph{qualidade}, \emph{é affirmativo}, \emph{negativo}, ou \emph{infinito}; --- por causa da \emph{relação}, \emph{é categorico}, \emph{hypothetico} ou \emph{disjunctivo}; --- e, finalmente, por causa da \emph{modalidade}, \emph{é problematico}, \emph{assertorio} ou \emph{necessario}. D'ahi as doze \emph{formas} ou \emph{categorias} da intelligencia, a saber: \emph{unidade}, \emph{multidão}, \emph{totalidade} (\emph{QUANTIDADE}); \emph{affirmação}, \emph{negação}, \emph{limitação} (\emph{QUALIDADE}); \emph{substancia} e \emph{modo}, \emph{causa} e \emph{effeito}, \emph{acção} e \emph{paixão} (\emph{RELAÇÃO}); \emph{possibilidade} e \emph{impossibilidade}, \emph{existencia} e \emph{inexistencia}, \emph{necessidade} e \emph{contingencia} (\emph{MODALIDADE}).\par A intelligencia, applicando estas \emph{formas} aos elementos sensiveis, produz os objectos do nosso conhecimento. Por isso os juizos, pelos quaes conhecemos as coisas, são \emph{syntheticos a-priori}; --- \emph{syntheticos}, porque o predicado não está incluido na essencia da sujeito, --- \emph{a-priori}, porque não se fundam na experiencia, mas derivam \emph{espontaneamente} da intelligencia.\par Mais elevada que a \emph{intelligencia} é a \emph{razão}; pois que todo o conhecimento, passando dos sentidos para a intelligencia, completa-se na \emph{razão}. A razão reduz á unidade os conceitos, e, por meio d'estes, os phenomenos, em virtude de tres ideas: --- da idea do \emph{sujeito pensante}, isto é, da unidade absoluta das condições internas, --- da idea do \emph{mundo}, isto é, da unidade absoluta das condições externas, --- da idea de \emph{Deus}, isto é, da unidade absoluta das condições de todas as coisas, que se alcançam pelo pensamento.\par D'esta theoria \emph{Kant} deduz os dois seguintes corollarios: --- 1.º) Não po\pgn{309}demos ter nenhum conhecimento das coisas, que não estão ao alcance da experiencia sensivel; como são a \emph{espiritualidade} e a \emph{liberdade} da nossa alma, a \emph{existencia} de \textsc{Deus}, etc. --- 2.º) Pela experiencia sensivel só percebemos os phenomenos; pois que os objectos, como são em si mesmos, ou os \emph{noumenos}, são incognosciveis.\par b) \emph{Fichte}, \emph{Schelling} e \emph{Hegel} levaram ao exaggero o \emph{transcendentalismo} de \emph{Kant}. O systema d'estes tres escriptores foi exposto na \emph{Cosmologia}.\par c) V. \emph{Cousin}, seguindo os vestigios dos philosophos \emph{transcendentaes}, explicou a \emph{origem das ideas} pela revelação de uma razão impessoal, que nos se manifestaria espontaneamente. Disse que o principio e o germen de todas as nossas ideas se encontra num acto primitivo, que a intelligencia produz sem reflexão livre, e pelo qual percebemos tudo o que a razão impessoal confusamente nos manifesta. Esse acto compõe-se de tres objectos ou elementos, a saber: o sujeito pensante, o mundo exterior, e as relações de ambos com a sua causa, que é \textsc{Deus}.}.

189. \textbf{O idealismo é falso.}

a) O \emph{idealismo} não quer attribuir, na origem dos nossos conhecimentos, nenhuma parte ás faculdades sensitivas nem aos objectos exteriores; porque parte do principio de que a percepção \emph{sensitiva} é apenas affectiva e subjectiva. Ora este principio é falso; pois que a percepção \emph{sensitiva} não é \textbf{apenas} affectiva e subjectiva, mas é tambem representativa e objectiva. Logo o \emph{idealismo} é falso.

b) O \emph{idealismo} ensina que os nossos conhecimentos derivam unica e exclusivamente da intelligencia, sem dependencia dos sentidos e dos objectos exteriores; de modo que é a intelligencia que cria os objectos. Ora esta doutrina é falsa. Por quanto, a experiencia attesta que a nossa intelligencia não cria os objectos, pelo facto de os perceber, mas percebe os objectos, pelo facto d'estes existirem; pois é levada á percepção por uma determinação proveniente dos objectos exteriores, que existem em si, e independentemente do nosso pensamento. Logo o \emph{idealismo} é falso.

c) O \emph{idealismo} não só desvirtua a origem dos nossos conhecimentos, mas leva até ao mais completo \emph{scepticismo}. Com effeito, se a origem dos conhecimentos dependesse unica e exclusivamente da \emph{intelligencia}, sem dependencia alguma dos objectos exteriores, não poderia existir a verdade \emph{absoluta}, tudo seria \emph{relativo}, uma e a mesma coisa seria verdadeira para uns e falsa para outros, e os nossos conhecimentos seriam sonhos da intelligencia. Ora estas consequencias são absurdas. Logo o \emph{idealismo} é falso \footnote{Refutemos agora as diversas especies de \emph{idealismo}.\par a) O systema de \emph{Kant} é absurdo no \emph{principio}, no \emph{processo}, e no \emph{termo}.\par \emph{a}) E' absurdo no \emph{principio}. Porquanto, \emph{Kant} propõe-se seguir o methodo \emph{a-priori}, prescindindo inteiramente da experiencia, e depois começa por admittir a existencia do sujeito e do objecto, --- existencia que só póde ser conhecida pela experiencia. --- Nem se diga que a existencia do sujeito e do objecto é admittida por \emph{Kant} como hypothese necessaria para raciocinar. Porquanto, uma simples hypothese, e \emph{que no fim se rejeita como falsa}, não póde ser o fundamento do raciocinio.\par \emph{b}) E' absurdo no \emph{processo} pelas seguintes razões:\par 1.º) \emph{Kant} attribue as representações \emph{abstractas} do \emph{espaço} e do \emph{tempo} á \emph{sensibilidade}. Ora isto é falso: porque a \emph{sensibilidade}, por ser uma faculdade \emph{organica}, só póde perceber o \emph{concreto} e o \emph{singular}.\par 2.ª) Segundo \emph{Kant}, as doze \emph{categorias} são noções \emph{abstractas}, ás quaes nenhuma realidade corresponde. Embora. Mas como explicar o apparecimento das \emph{categorias}? --- São \emph{innatas}? Mas então temos o \emph{innatismo} com todos os seus absurdos. --- Foram infundidas por \textsc{Deus}? Não; porque \textsc{DEus} não engana, e nos enganaria, se tivesse posto na nossa intelligencia ideas, que nada mostram, e representações, que nada representam. --- Derivaram dos objectos? Não; porque \emph{Kant} diz que nenhum objecto lhes corresponde. --- Foram produzidas pela intelligencia, independentemente dos objectos? Mas, neste caso,\pgn{310}a questão não se resolve; porque poderiamos sempre perguntar \emph{porque} e como a intelligencia produziu essas formas.\par 3.ª) De mais, a divisão e a disposição das \emph{doze} categorias, como demonstrámos na \emph{Logica}, são falsas.\par 4.ª) Finalmente, os juizos \emph{syntheticos a-priori}, como já provámos, são contrarios á natureza da nossa intelligencia; porque, se não são formados pela analyse do sujeito, nem pela experiencia, devem necessariamente derivar de um instincto cego e irracional; o que é falso.\par c) E' absurdo no \emph{termo}. \emph{Kant}, dizendo, que não podemos conhecer os objectos que escapam ao alcance dos sentidos, tira ás percepções sensitivas toda a objectividade, e faz do proprio sujeito pensante um mero phenomeno. Ora isto destroe toda a certeza do conhecimento humano, e introduz um \emph{scepticismo} mais pernicioso que o de \emph{Pyrro} e de \emph{Sexto Empirico}.\par b) Os systemas de \emph{Fichte}, \emph{Schelling} e \emph{Hegel} são falsos, porque negam a \emph{objectividade} das ideas, confundem a ordem ideal com a real, defendem o \emph{pantheismo}, e contêm outros absurdos, que apontámos na \emph{Cosmologia}.\par c) O systema de \emph{Cousin} é absurdo pelas seguintes razões:\par 1.ª) O escriptor francez, quando affirma que ha uma razão impessoal, que é distincta de nós, apparece em nós e costitue a intelligencia de cada um, além de professar o mais pernicioso \emph{pantheismo ideal}, contradiz a consciencia, a qual attesta que cada um de nós possue uma intelligencia pessoal, propria, contingente e limitada.\par 2.ª) Este systema destroe toda a certeza do conhecimento humano. Porquanto, aquelle acto primitivo, de cuja existencia \emph{Cousin} não dá uma razão plausivel, não é evidente, acredita-se cega e inconscientemente, e por isso todos os outros actos, que derivam d'elle, são destituidos de evidencia; pois que o principiado não póde exceder em perfeição o seu principio.\par 3.ª) Repugna que a primeira apprehensão da intelligencia seja composta d'aquelles tres objectos ou elementos; porque o Ente infinito, que é \textsc{Deus}, não se conhece pela \emph{intuição}, mas pelo \emph{raciocinio}. --- O proprio \emph{Cousin} concorda com a nossa doutrina. Com effeito, elle ensina que o Infinito se nos manifesta, emquanto a intelligencia, conhecendo que o sujeito e o objecto têm limites e por isso não encerram em si a razão sufficiente da sua existencia, remonta á Causa suprema, que deve ser infinita e absoluta. Ora isto é raciocinar, é deduzir de um juizo um outro, baseando-se no principio de que o finito depende do infinito, o relativo do absoluto.}.

%%>>> t190

\tese{190}{Psychologismo racional.}
O \emph{psychologismo racional} admitte que os nossos conhecimentos \emph{intellectuaes} derivam da \pg{294} \pg{295} \pg{297} \pg{298} \pg{300} \pg{302} \pg{303} \pg{308} \pg{309} \pg{310}  intelligencia, auxiliada pelos sentidos; emquanto os sentidos apresentam o objecto, que, abstrahido das propriedades materiaes e concretas, se une á intelligencia \footnote{O \emph{psychologismo racional} foi ensinado por \emph{Aristoteles} e seguido por \emph{Santo Agostinho}, \emph{Santo Thomaz}, \emph{S}. \emph{Bonaventura}, e por todos os melhores escholasticos da edade media e moderna.}.

%%>>> t191

\tese{191}{O psychologismo racional é verdadeiro.}

a) Um systema, que explica a origem dos nossos conhecimentos por um modo conforme com a indole da intelligencia, é verdadeiro; pois que a operação, derivando da natureza do ente, deve exprimir a sua indole. Ora o \emph{psychologismo racional}, \pg{311}  ensinando que os nossos conhecimentos derivam da intelligencia, auxiliada pelos sentidos, explica a origem d'aquelles por um modo conforme com a indole da propria intelligencia; porque esta faculdade, no seu exercicio, depende \emph{objectivamente} dos sentidos (n.os 116 e 117). Logo o \emph{psychologismo racional} é

verdadeiro \footnote{E' certo que o homem pensa, emquanto tal. Ora o homem é um composto de alma e de corpo. Logo é necessario que a alma e o corpo concorram, na medida da sua natureza, para o acto do pensamento. O celebre \emph{Lacordaire} dizia: « O homem não pensa senão com a totalidade do seu ser, assim como não vive senão com a totalidade do seu ser. Todos os systemas \emph{idealistas} e \emph{materialistas} são falsos, porque fazem d' elle mera intelligencia ou mera materia » (Conf. 48). E' o que antes tinha ensinado S. Thomaz nas seguintes palavras, já citadas: « Operatio proportionatur virtuti et essentiae; intellectivum autem hominis est in sensitivo; et ideo propria operatio ejus est intelligere intelligibilia in phantasmatibus » (\emph{De mem}. et \emph{rem}., l. I).}.

b) O \emph{psychologismo racional} é constituído por estes dois princípios: --- 1.º) a \emph{essencia} das coisas, que forma o objecto da nossa intelligencia, encontra-se nas proprias coisas, embora por um modo \emph{singular} e \emph{material}; --- 2.º) a mesma \emph{essencia}, para de \emph{singular} e \emph{material} se tornar \emph{universal} e \emph{immaterial}, exige uma força \emph{abstractiva}. Ora estes dois principios são verdadeiros; porque todas as coisas possuem uma essencia real, pela qual são o que são e pertencem a uma determinada especie; e porque a nossa alma possue effectivamente uma \emph{força abstractiva}, que é a \emph{intelligencia agente} (n.º 120). Logo o \emph{psychologismo racional} é verdadeiro.

c) A experiencia attesta que os nossos conhecimentos \emph{intellectuaes} começam, quando subimos do particular para o universal, do phenomeno para a lei, do effeito para a causa. Ora não poderiamos produzir essas operações, se não empregassemos, ao mesmo tempo, a \emph{intelligencia} e os \emph{sentidos}; pois que o particular, o phenomeno, o effeito só póde ser percebido pelos \emph{sentidos}, --- e o universal, a lei, a causa só póde ser conhecida pela \emph{intelligencia}. Logo os nossos conhecimentos \emph{intellectuaes} derivam da \emph{intelligencia}, auxiliada pelos \emph{sentidos}. Logo o \emph{psychologismo racional} é verdadeiro \footnote{O \emph{psychologismo racional} admitte tudo o que de verdadeiro se encontra nos outros systemas, rejeitando tudo o que nelles se contém de falso.\par Concede ao \emph{innatismo} que os conhecimentos intellectuaes exigem uma virtude superior á dos sentidos; mas contesta que a alma, desde o momento da sua creação e infusão no corpo, traga impressos esses conhecimentos.\par Reconhece com o \emph{tradicionalismo} que, sem a revelação, o homem não póde conhecer \emph{todas} as verdades da ordem metaphysica e moral; mas nega\pgn{312}que o homem, sem a revelação ou sem um auxilio exterior, não possa conhecer \emph{algumas} d'aquellas verdades.\par Admitte com o \emph{ontologismo} que só a intelligencia é capaz de conhecer os objectos abstractos e immateriaes; mas nega que esses objectos sejam contemplados na visão \emph{intuitiva} de \textsc{Deus}.\par Concorda com o \emph{sensualismo} em que os \emph{sentidos} concorrem para a producção dos conhecimentos \emph{intellectuaes}; mas nega que os mesmos sentidos concorram como causa unica e principal.\par Reconhece com o \emph{idealismo} que, na acquisição das ideas, a intelligencia opéra; mas não admitte que a intelligencia opére independentemente dos sentidos e dos objectos externos.\par \emph{Paulo Janet} diz: « Sem tomar partido decisivo por nenhuma d'estas theorias (acerca da origem dos conhecimentos), a que nos parece mais simples, menos conjectural e que mais se approxima dos factos, é a theoria \emph{aristotelica} acerca da \emph{intelligencia agente} » (\emph{Traité de Phil}., Psyc., ch. IX).\par Advirta-se a differença entre o \emph{psychologismo racional} e o \emph{sensualismo}. Concordam ambos estes systemas, quando dizem que os \emph{sentidos} concorrem para a producção dos nossos conhecimentos intellectuaes. Divergem, quando tratam de explicar a natureza d'aquelle concurso. Na verdade, o \emph{sensualismo} ensina que a nossa intelligencia é completamente \emph{passiva}, e precisa dos sentidos, como de mestres, que lhe communiquem as ideas; de modo que os sentidos são a causa principal e efficiente dos conhecimentos \emph{intellectuaes}. Pelo contrario, o \emph{psychologismo racional} sustenta que a intelligencia é tambem \emph{activa}, e precisa dos sentidos, como de servos, para estes lhe apresentarem o objecto, que só ella percebe: de modo que os sentidos são apenas causa instrumental dos nossos conhecimentos.}.

%%>>> art-26

\section*{ARTIGO II.}
\addcontentsline{toc}{section}{ARTIGO II.}

\pg{312} \textbf{Modo por que se produz o conhecimento}

\textbf{intellectual}.

%%>>> t199

\tese{192}{Estado da questão.}
No artigo antecedente dissemos que os nossos conhecimentos \emph{intellectuaes} derivam da \emph{intelligencia}, auxiliada pelos \emph{sentidos}. --- Agora devemos indicar o modo por que esses conhecimentos se produzem, enumerando as \emph{funcções}, que, para isso, a \emph{intelligencia} e os \emph{sentidos} exercem \footnote{As theorias, relativas ao modo por que se produz o conhecimento intellectual, foram expostas, na maxima parte, no art. I do cap. III da secção II d'este tratado. Para ahi remettemos o leitor. --- Como dissemos, ninguem deve admirar-se de que, tendo destinado esta secção III para o estudo das \emph{operações humanas}, fallassemos d'ellas na secção II, que era destinada para as \emph{faculdades}. As \emph{operações} fazem conhecer e distinguir as \emph{faculdades}, e por isso, tratando d'estas, não podiamos deixar de fallar naquellas.}.

193. \textbf{Os sentidos externos} apresentam o objecto do conhecimento \textbf{intellectual}.

a) O objecto do conhecimento \emph{intellectual} --- ou é um ente \pg{313}  \emph{material}, --- ou é um ente \emph{immaterial}. Se é um ente \emph{material}, é apresentado pelos \emph{sentidos externos}; pois o \emph{material} só póde ser percebido por faculdades \emph{materiaes}, como são os \emph{sentidos}. Se é um ente \emph{immaterial}, tambem é apresentado, de algum modo, pelos \emph{sentidos externos}; pois que só podemos elevar-nos ao immaterial por meio do material, e este é objecto dos \emph{sentidos}.

b) A experiencia attesta que os homens, destituidos de algum sentido \emph{externo}, não possuem a sciencia dos objectos \emph{próprios} d'esse sentido. Ora isto mostra que os \emph{sentidos externos} apresentam o objecto do conhecimento \emph{intellectual} \footnote{Cf. os n.os 116 e 117 com as respectivas \emph{notas}.}.

194. \textbf{A imaginação forma a especie sensivel do objecto.} --- O objecto do conhecimento \emph{intellectual} deve ser apresentado pelos \emph{sentidos externos}. Sendo assim, é necessario que o mesmo objecto se una, por meio de uma sua \emph{imagem}, aos mesmos sentidos; porque não póde unir-se com a sua realidade physica. A \emph{imagem} do objecto, passando pelo \emph{sentido commum}, vae unir-se com a imaginação. Esta faculdade, unida com a imagem, reprodul-a e exprime-a em si mesma. Essa reproducção ou expressão é a \emph{especie sensivel} do objecto.

A \emph{especie sensivel}, que representa a \emph{essencia}, cercada das condições materiaes, é indispensavel para a percepção \emph{intellectual}. Por quanto, se a \emph{imaginação}, por qualquer desarranjo, não poder formar a mesma \emph{especie}, ou não a formar semelhante ao objecto, a intelligencia será tambem perturbada no seu exercicio, não percebendo o objecto, ou não o percebendo como elle é em si mesmo \footnote{Nas \emph{notas} ao art. IV do cap. II da secção I d'este tratado explicámos, com o sufficiente desenvolvimento, as funcções da \emph{imaginação} no conhecimento \emph{intellectual}. --- Repetimos que o \emph{phantasma} ou a \emph{imagem sensivel} representa tambem a \emph{essencia} das coisas; porque, pelos sentidos \emph{externos}, percebemos as qualidades sensiveis, não \emph{abstractas} e \emph{separadas} da substancia, mas \emph{concretas} e \emph{inherentes} á propria substancia.}.

195. \textbf{A intelligencia agente torna a especie sensivel capaz de operar na intelligencia possivel.} --- A \emph{especie sensivel}, formada pela \emph{imaginação}, deve operar na \emph{intelligencia possivel}, para determinar esta faculdade á percepção do objecto. Mas a \emph{especie sensivel}, sendo material e concreta, não póde, só por si, operar na intelligencia, que é immaterial. E' necessaria, pois, uma força que, operando na \emph{especie sensivel}, abstraia da essencia as propriedades materiaes, e assim a torne capaz de operar na

intelligencia, e de produzir a \emph{especie intelligivel}, necessaria para a percepção intellectual. Essa força \emph{abstractiva} é a \emph{intelligencia agente}. Logo a \emph{intelligencia agente} torna a especie \emph{sensivel} capaz de operar na \emph{intelligencia possivel} \footnote{Nos n.º 119 e 120 (com as respectivas \emph{notas}) provámos a necessidade da \emph{especie intelligivel}, e indicámos o modo por que a \emph{especie sensivel} ou o \emph{phantasma} e a \emph{intelligencia agente} concorrem para a producção d'ella.\par A intelligencia \emph{possivel} é a que chamamos simplesmente \emph{intelligencia}. Diz-se \emph{possivel}, porque, como demonstrámos, ella \emph{póde} produzir sempre o seu acto, mas não o está produzindo sempre.}.

\textbf{196. A intelligencia possivel, determinada pela especie intelligivel, produz o verbo mental.} --- A intelligencia \emph{possivel} não póde perceber um certo objecto, se este, por meio de uma sua imagem immaterial, que é a \emph{especie intelligivel}, a não determine ao acto. Determinada pela \emph{especie intelligivel}, a mesma intelligencia reproduz e exprime em si mesma a semelhança espiritual do objecto percebido. Esta reproducção ou expressão chama-se \emph{verbo mental}; porque é por esta semelhança que a intelligencia \emph{falla} comsigo mesma, indicando ou exprimindo o objecto percebido. --- A intelligencia humana, sempre que entende, produz o \emph{verbo mental} \footnote{Tres são os elementos, que concorrem para o acto da percepção intellectual: a \emph{especie intelligivel}, que informa a intelligencia e a determina ao acto, --- o \emph{verbo}, que a intelligencia, fecundada pela especie intelligivel, forma em si mesma e que é a expressão ideal do objecto conhecido, --- e o \emph{objecto}, que é o termo \emph{extrinseco} do proprio acto perceptivo (Cf. S. Thom., \emph{De pot.}, q. 8, a. 1). Analysemos estes elementos e indiquemos o modo por que cada um concorre para o conhecimento. Neste numero trataremos da \emph{especie intelligivel} e do \emph{verbo}; no numero seguinte, do \emph{objecto}.\par a) \textsc{EsPECIE} INTELLIGIVEL. --- A \emph{especie intelligivel} é, como dissemos, a \emph{semelhança}, ou a \emph{representação}, ou o \emph{imagem} de um objecto, \emph{abstrahida das condições materiaes}. --- Esta especie é \emph{necessaria} para o conhecimento intellectual. Faz-se o conhecimento pela união do objecto conhecido com o sujeito cognoscente. Ora o objecto, que não pode unir-se com o cognoscente pela sua realidade physica, deve unir-se com elle por uma sua imagem, ou representação. Diz S. Thomaz: «\emph{Intelligens et intellectum, prout ex eis est effectum unum quid, quod est intellectus in actu, sunt unum} principium \emph{hujus actus, qui est intelligere. Et ideo ex eis efficitur unum quid inquantum intellectum coniugitur intelligenti sive per essentiam suam, sive per similitudinem} » (\emph{De cognit. ang.}, a. 1). D'este modo, o objecto conhecido tem o ser no cognoscente; embora o ser que tem no cognoscente seja diverso d'aquelle, que tem em si mesmo: pois o ser que tem em si é \emph{real}, e o ser que tem no cognoscente é \emph{intencional}. Este modo de existir num sujeito cognoscente revela a perfeição dos seres intellectuaes. --- E', portanto, evidente a \emph{natureza} e o \emph{officio} da especie intelligivel. Esta é uma forma \emph{objectiva}, isto é, representativa do objecto, que reduz ao acto a potencia intellectual, e com esta\pgn{315}constitue o principio completo e adequado do conhecimento, sendo a potencia principio \emph{passivo} e a especie principio \emph{activo}. --- Mas uma coisa é o \emph{principio} de um acto, outra é o seu \emph{termo}. Por isso nem a intelligencia nem a especie podem ser o termo \emph{directo} da acto perceptivo; mas são apenas o \emph{principio} ou o \emph{meio} para a percepção dos objectos. Diz S. Thomaz: « \emph{Similitudo rei intellectae, quae est species intelligibilis, est forma secundum quam intellectus intelligit} » (\emph{Summ. Theol.}, p. I. q. 85, a. 2). Por isso, devendo a percepção \emph{directa} ter um objecto, este não pode ser nem a intelligencia, nem o seu acto, nem a especie, mas deve ser alguma coisa distincta do sujeito pensante e das suas modificações. Diz S. Thomaz: « \emph{Id quod intelligitur primo est res, cujus species intelligibilis est similitudo} » (Ib.). E a especie não é o que se conhece, mas é o \emph{meio} pelo qual uma coisa é conhecida: « \emph{Species intelligibilis non est id quod intelligitur, sed id quo intellectus intelligit} » (Ib.) E noutro logar: « \emph{Objectum intellectus est ipsa rei essentia; quamvis essentiam rei cognoscat per ejus similitudinem, sicut per medium cognoscendi, non sicut per objectum, in quod primo fertur ejus visio} » (\emph{De Mente}, a. 4, ad 4). --- Note-se bem que a intelligencia percebe \emph{primeiramente}, pelo acto \emph{directo}, o objecto, de que a especie é semelhança; embora, \emph{secundariamente}, ou pelo acto \emph{reflexo}, perceba tambem a especie, que a informa: « \emph{Quia intellectus supra ipsum reflectitur, secundum eamdem reflexionem intelligit et suum intelligere, et speciem qua intelligit; et sic species intellecta secundario est id quod intelligitur} » (\emph{Sum. Th.}, p. I, q. 85, a. 2).\par Portanto a especie intelligivel é o \emph{meio}, pelo qual uma coisa é conhecida. Mas, como indicámos no tratado das operações \emph{sensitivas}, ha um triplice meio de percepção: o meio \emph{sub quo}, o meio \emph{quo} e o meio \emph{in quo}. O meio \emph{sub quo} é uma perfeição ou energia, que torna a intelligencia capaz de perceber as coisas; e, na ordem \emph{natural}, é a luz da \emph{intelligencia agente}. O meio \emph{quo} é o que determina a faculdade, e com esta constitue um principio completo e adequado do acto. O meio \emph{in quo} é um objecto, que, percebido primeiramente em si mesmo, leva a faculdade á percepção de uma coisa distincta d'elle. Ora o meio \emph{sub quo} e o meio \emph{quo} são apenas \emph{meios} e não \emph{objectos} da percepção \emph{directa}; porque a percepção não tem por termo \emph{directo} e \emph{immediato} esses meios, mas o proprio objecto. Mas o meio \emph{in quo} é, ao mesmo tempo, \emph{meio} e \emph{objecto} da percepção \emph{directa}; porque esse meio primeira ou directamente é percebido em si mesmo e depois leva á percepção do objecto, que, por isso, é conhecido de um modo \emph{mediato}; tal é o \emph{effeito}, que, depois de ter sido percebido em si mesmo, leva á percepção da causa respectiva. Ora a especie intelligivel é um meio \emph{quo}; porque determina a faculdade intellectual, e por isso não é o termo \emph{directo} ou \emph{immediato} da \emph{percepção}, mas é o \emph{meio} pelo qual \emph{directa} ou \emph{immediatamente} percebemos o objecto, que ella representa (Cf. S. Th., \emph{De verit}, q. 10; \emph{De mente}, a. 8 ad 2).\par b) \textsc{VERBO} MENTAL. --- A intelligencia, informada pela especie intelligivel, produz uma imagem espiritual, que exprime e representa a essencia do objecto. Esta imagem diz-se \emph{verbo} ou \emph{palavra mental}, porque é por ella que o espirito falla comsigo e indica a si mesmo o objecto percebido. O \emph{verbo mental}, pois, é a semelhança ou imagem intencional do objecto, gerada pela\pgn{316}intelligencia. Diz S. Thomaz: « \emph{Id, quod intellectus in concipiendo format, est verbum} » (\emph{Sum. Th.}, p. 1, q. 34, a. 1 ad 2).\par O \emph{verbo mental} distingue-se do \emph{verbo oral}, que é um signal externo, manifestativo do proprio verbo mental. --- Pelo \emph{verbo mental} fallamos comnosco, e exprimimos directamente o objecto; pelo \emph{verbo oral} fallamos com os outros, e indicamos directamente o nosso verbo mental. --- O nome \emph{verbo} foi primeiramente empregado para significar a externa locução da bocca, e, em seguida, para significar a interna locução da intelligencia; pois costumamos indicar a coisa significada pelo nome do signal, e vice-versa. --- Todavia, se se considera o valor da significação, e não a origem da nomenclatura, vê-se que o nome \emph{verbo}, ou \emph{palavra}, convem primeiramente ao conceito da intelligencia e depois á expressão externa. Na verdade, emquanto o \emph{verbo oral} depende necessariamente, na sua existencia e na sua significação, do \emph{verbo mental}, como da sua causa efficiente e final, porque a palavra externa é pronunciada por quem tem um conceito e para exprimir um conceito, o \emph{verbo mental} é independente, na existencia e na significação, do \emph{verbo oral}. Diz S. Thomaz: « \emph{Nomina imponuntur secundum quod cognitionem de rebus accipimus... Et ideo, quia verbum exterius, quum sit sensibile, est magis notum nobis quam interius, secundum nominis impositionem per prius vocale verbum dicitur verbum, quam verbum interius; quamvis verbum interius naturaliter sit prius, utpote exterioris causa efficiens et finalis} » (\emph{De Ver.}, q. 4, a. 1; cf. \emph{Sum. Th.}, p. I, q. 34, a. 1).\par O \emph{verbo mental} é proprio de cada ser intelligente. Toda a percepção intellectual tem o seu verbo; porque toda a acção deve ter o seu termo, e o termo da acção intellectual é o \emph{verbo}. Diz S. Thomaz: « \emph{Semper, cum actu intelligitur aliquid, verbum formatur} » (\emph{De nat. verbi intellectus}). --- O verbo deriva de dois principios: --- do principio \emph{passivo}, que é a faculdade intellectual, --- e do principio \emph{activo}, que é a especie intelligivel, representativa do objecto. Sendo uma verdadeira conceição, o \emph{verbo} deve exprimir e revestir, e effectivamente exprime e reveste, a natureza d'esses dois principios; porque elle, na sua simplicidade e espiritualidade \emph{entitativa}, é a imagem viva e bella da intelligencia, e, na sua significação \emph{objectiva}, é a imagem fidelissima do objecto, que representa. Diz o S. Doutor: « \emph{Quicumque intelligit, ex hoc ipso quod intelligit, procedit aliquid intra ipsum, quod est conceptio rei intellectae, ex vi intellectiva proveniens, et ex ejus (rei) notitia procedens. Quam quidem conceptionem vox significat, et dicitur verbum cordis (mentis), significatum verbo oris} » (\emph{Sum. Th.}, p. I, q. 27, a. 1). --- E, como a operação de nossa intelligencia é duplice, isto é, a simples apprehensão (idéa) e a composição ou divisão (juizo affirmativo ou negativo), assim ha um duplice verbo --- um que exprime a idéa, outro que exprime o juizo. Diz S. Thomaz: « \emph{Verbum intellectus nostri... est id ad quod operatio intellectus nostri terminatur..., quod dicitur conceptio intellectus; sive sit conceptio significabilis per vocem incomplexam, ut accidit quando intellectus format quidditates rerum, sive per vocem complexam, ut accidit quando intellectus componit et dividit} » (\emph{De Verit.}, q. 4, a. 2).\par O \emph{verbo} diz-se \emph{termo} da percepção intellectual, emquanto esta se considera como \emph{acto}, e não como \emph{conhecimento}. Na verdade, o \emph{acto perceptivo},\pgn{317}se considerado como \emph{acto} tem por termo o \emph{verbo}, considerado como \emph{perceptivo} tem por termo o \emph{objecto}, representado no verbo. Por isso costuma dizer-se que o \emph{verbo} é o termo \emph{subjectivo} e \emph{intrinseco} da percepção intellectual, e o \emph{objecto} é o termo \emph{objectivo} e \emph{extrinseco}. --- D'ahi se segue que, embora no ordem do \emph{tempo}, a formação do verbo não se distinga da percepção do objecto, comtudo, na ordem da \emph{natureza}, a formação do verbo precede a percepção do objecto, expresso pelo verbo. Porquanto, primeiramente a intelligencia é informada pela specie intelligivel, que é o principio sufficiente do conhecimento, pois basta para determinar a intelligencia ao acto perceptivo. Não temos ainda a percepção formal, mas só em raiz, porque previa á formação do verbo. A intelligencia, informada pela especie, produz o verbo, que é o termo do acto intellectual; e, produzindo o verbo, percebe nelle, e immediatamente, o objecto. Mas assim como o objecto precede o acto intellectual, porque o objecto determina o acto, assim tambem a formação do verbo precede a percepção do objecto, porque é no verbo e pelo verbo que a intelligencia percebe o objecto. Não é uma prioridade de \emph{tempo}, mas de \emph{natureza}. Diz S. Thomaz: « \emph{Prius natura est intellectus informatus specie, quae est principium sufficiens intelligendi, quam gignatur verbum; et ideo intelligere in radice prius est verbo, et verbum est terminus actionis intellectus. Sed quia objectum non habetur nisi in verbo; objectum autem prius est, quam quaelibet actio ad ipsum terminata; ideo verbum prius est quam intelligere} » (\emph{De natura verbi intellectus}).\par Dissemos que, se o termo \emph{subjectivo} e \emph{intrinseco} do acto intellectual é o \emph{verbo}, o termo \emph{objectivo} e \emph{extrinseco} do mesmo acto é o \emph{objecto}. A intelligencia pronuncia o verbo, e, pronunciando-o, entende o objecto, que no verbo se manifesta. Diz S. Thomaz: « \emph{Dicere importat principaliter habitudinem ad verbum conceptum; nihil enim est aliud dicere quam proferre verbum. Sed mediante verbo importat habitudinem ad rem intellectam, quae in verbo prolato manifestatur intelligenti} » (\emph{Sum. Th.}, p. I, q. 34, a. 1 ad 3). Na verdade, o verbo é o conceito mental; ora o conceito mental é o conhecimento de um objecto. Portanto, a intelligencia não vê no verbo o representante de um objecto, nem o modo de representar, mas vê directamente o proprio objecto representado. Nem podia ser d'outro modo. Recebendo a \emph{especie intelligivel}, que é a forma intencional do objecto, a intelligencia fica como que transformada no proprio objecto, e produz o verbo, e no verbo conhece e diz o objecto. Diz o S. Doutor: « \emph{Anima enim quasi transformata est in rem per speciem, qua agit quidquid agit: unde cum intellectus ea informatus est actu, verbum producit, in quo rem illam dicit, cujus speciem habet} » (\emph{De nat. verbi int.}) --- E' verdade que o \emph{verbo} é o termo da percepção, emquanto esta é \emph{acto}, mas tambem é verdade que o \emph{objecto} é o termo da percepção, emquanto esta é \emph{conhecimento}; porque no verbo, que é o termo do acto, resplandece e percebe-se o proprio objecto. O angelico Doutor compara o verbo a um espelho, cujas dimensões não excedem as do objecto, que nelle se reflecte: « \emph{Verbum est tamquam speculum in quo res cernitur, sed non excedens id quod in eo cernitur} » (\emph{De nat. verbi int.}). Ora, assim como quem olhasse para esse espelho, perceberia primeira e directamente o objecto, que nelle se reflecte, assim tambem a intelligencia, pronunciando o verbo,\pgn{318}vê primeira e directamente o objecto, que elle representa. Esconde-se o espelho aos olhos do corpo para só manifestar o objecto, que nelle resplandece; esconde-se o verbo aos olhos da intelligencia para só pôr em evidencia o objecto, que elle representa. --- Mas quem nos assegura --- que o verbo é a expressão fiel do objecto representado, --- e que no verbo e pelo verbo percebemos directamente o proprio objecto? --- A estas perguntas responderemos nas notas, relativas ao numero seguinte.\par Do que deixamos dito vê-se o erro de \emph{Rosmini}, quando no livro --- \emph{nuovo saggio sull'origine delle idee} --- affirma que o \emph{verbo} consiste \emph{exclusivamente} no \emph{juizo} e se produz unicamente pela \emph{reflexão}. --- E' falso que o \emph{verbo} consiste e se dá \emph{exclusivamente} no \emph{juizo}. Por quanto, todo o acto intellectual deve ter o seu termo \emph{intrinseco}; porque é impossivel que se dê uma acção sem o effeito respectivo, e que o producto de uma acção \emph{immanente} exista fóra da causa. Ora a \emph{simples apprehensão} (á qual pertence tambem a \emph{definição}, que é a expressão da essencia de uma coisa) é um acto intellectual. Logo a \emph{simples apprehensão} deve ter por termo o \emph{verbo}. Logo é falso que o \emph{verbo} consiste e se dá \emph{exclusivamente} no \emph{juizo} (Cf. S. Th. \emph{De diff. div. verbi et hum.}). --- De mais, assim como o \emph{verbo oral}, para ser \emph{verbo}, não é necessario que seja \emph{proposição}, mas basta que seja \emph{termo}; assim tambem o \emph{verbo mental}, para ser \emph{verbo}, não é preciso que seja \emph{juizo}, mas basta que seja \emph{simples apprehensão}. --- E' falso tambem que o \emph{verbo} se produz unicamente pela \emph{reflexão}; pois o \emph{verbo} produz-se na \emph{simples apprehensão}, e esta não é uma reflexão. --- Se fosse verdadeira a opinião de \emph{Rosmini}, todo o conhecimento seria \emph{reflexo}, porque todo o conhecimento produz o seu \emph{verbo}. Ora é falso que todo o conhecimento é \emph{reflexo}; porque o acto \emph{reflexo} suppõe necessariamente o \emph{directo} (Cf. S. Th. \emph{De nat. verbi int.}).\par \textsc{c)} \textsc{ANALOGIAS} DO VERBO MENTAL. --- O \emph{verbo mental} tem analogia com a \emph{especie intelligivel}, com a \emph{especie expressa sensivel}, e com o Verbo \emph{divino}.\par \emph{a}) \emph{Analogia} do \emph{verbo mental com a especie intelligivel}. --- 1.º) O \emph{verbo mental} e a \emph{especie intelligivel} convêm em que um e outra são formas espirituaes e semelhanças do objecto externo. Esta semelhança deriva da especie para o verbo; pois a intelligencia não gera o verbo, isto é, a imagem expressa do objecto, senão emquanto é informada pela especie, isto é, pela imagem impressa do objecto, visto que todo o agente opera em virtude e em conformidade da sua forma. --- 2.º) O \emph{verbo mental} e a \emph{especie intelligivel} differem em que --- a \emph{especie} é o producto da intelligencia \emph{agente}, e o \emph{verbo} é o producto da intelligencia \emph{possivel}; --- a \emph{especie}, determinando a intelligencia á percepção actual, é o \emph{principio} do conhecimento, ao passo que o \emph{verbo}, sendo o conceito que a intelligencia forma, quando conhece o objecto, é o \emph{termo} do conhecimento; --- a \emph{especie} representa o objecto como \emph{cognoscivel}, e por isso chama-se \emph{especie impressa}, mas o \emph{verbo} exprime ou representa o objecto como \emph{conhecido}, e por isso denomina-se \emph{especie expressa}. --- S. Thomaz explica esta analogia com as seguintes palavras: « \emph{Intellectus per speciem (intelligibilem) rei formatus, intelligendo format in se ipso quamdam intentionem rei intellectae (verbum).... Haec autem intentio intellecta, cum sit quasi terminus intelligibilis operationis, est aliud a specie}\pgn{319}\emph{intelligibili, quae facit intellectum in actu, quam oportet considerari ut intelligibilis operationis principium, licet utrumque} (species et verbum) \emph{sit rei intellectae similitudo. Per hoc enim quod species intelligibilis, quae est forma intellectus et intelligendi principium, est similitudo rei exterioris, sequitur quod intellectus intentionem} (verbum) \emph{format illi rei similem; quia quale est unumquodque, talia operatur} » (\emph{C. Gentes}, l. I, c. 53).\par Pergunta-se: Quae são as relações entre o \emph{verbo} e a \emph{idéa}? --- Respondemos que a \emph{idéa} pode tomar-se em dois sentidos: --- ou emquanto significa a especie intencional do objecto, que determina a intelligencia ao conhecimento, e assim é \emph{principio} da operação, --- ou emquanto significa a imagem do objecto, que a intelligencia, informada pela especie, produz e exprime em si mesma, e assim é \emph{termo} da operação. Tomada no primeiro sentido, a \emph{idéa} identifica-se com a \emph{especie intelligivel}, e distingue-se do \emph{verbo}, como o \emph{principio} de um acto se distingue do seu \emph{termo}. Tomada no segundo sentido, a \emph{idéa} não se distingue \emph{realmente} do \emph{verbo}, embora se distinga \emph{logicamente}. E distingue-se \emph{logicamente}; porque a \emph{idéa}, sendo uma forma representativa do objecto, indica apenas a relação do conceito com o proprio objecto; ao passo que o \emph{verbo}, exprimindo um \emph{acto} e uma \emph{percepção}, indica a relação do conceito com o principio, de que deriva, e com o objecto, que representa. De facto, não podemos pensar no \emph{verbo}, ou \emph{palavra}, sem que a entendamos como procedente do orgão da voz por um interior impulso vital, e como representativa de uma coisa. --- Por isso, embora o nosso \emph{verbo mental} possa chamar-se tambem \emph{idéa}, não se pode chamar \emph{idéa} o Verbo \emph{divino}; porque, sendo o Verbo \emph{divino} subsistente e pessoal, e consistindo tal subsistencia e personalidade na relação de opposição com o seu Principio, essa relação deve ser expressa, e não é expressa pelo vocabulo de \emph{idéa}, mas pelo de Verbo.\par \emph{b}) \emph{Analogia do verbo com a especie sensivel expressa}. --- O \emph{verbo} e a \emph{especie sensivel expressa} convém em que --- um e outra são \emph{imagens intencionaes} de um objecto externo; --- um e outra são \emph{termos} do conhecimento, sendo o verbo termo do conhecimento intellectual e a especie termo do conhecimento sensitivo; --- um e outra exprimem o objectivo respectivo; --- um e outra são como um espelho, em que se reflecte e se conhece o objecto. --- Mas, se ao \emph{conceito} da \emph{intelligencia} e á sua \emph{expressão externa} convem rigorosamente o nome de \emph{verbo}, este nome não pode convir á \emph{especie sensivel expressa}, nem a um signal externo, que de algum modo manifeste essa especie. Na verdade, o \emph{verbo mental}, exprimindo um juizo ou uma idéa, não pode indicar essa especie sensivel, que não é juizo nem idéa. Como tambem o \emph{verbo oral}, sendo expressão de conceitos intellectuaes, não pode significar a especie sensivel expressa. Não ha palavra, que exprima o phantasma na sua \emph{singularidade}; de modo que o conhecimento sensitivo, mesmo no homem, é \emph{mudo}. « A palavra, diz Müller, é o pensamento \emph{encarnado} ».\par \emph{c}) \emph{Analogia do nosso verbo com o Verbo divino}. --- Deus conhece e comprehende a sua Essencia infinita, e exprime o seu conhecimento numa Imagem, num Verbo, que reproduz e contem todas as perfeições divinas. Ora o Verbo divino e o nosso verbo convêm em que --- um e outro são o termo de um conhecimento, semelhante ao objecto conhecido, pois o nosso\pgn{320}verbo é a imagem expressa da coisa percepida, e o Verbo divino é a perfeitissima Imagem do Pae, que a exprime na contemplação da sua Essencia; --- o nosso verbo é espiritual, como espiritual é a intelligencia que o gera, mas é distincto d'ella, e o Verbo divino é a propria Imagem do Pae, e é uma Pessoa distincta do Pae. --- Mas ha enormes differenças entre o nosso verbo e o Verbo divino. O nosso verbo não procede da subtancia da nossa alma, mas da nossa intelligencia, e por isso é uma coisa \emph{accidental}, porque \emph{accidental} é toda a faculdade humana; o Verbo divino, nascendo da intelligencia do Pae, nasce da propria Essencia divina; porque, em Deus, as faculdades se identificam com a Essencia, e por isso é subsistente. O nosso verbo não é coevo ao seu principio, não é sempre em acto, não é immutavel; o Verbo divino é coeterno ao Pae, immutavel e eternamente gerado. O nosso verbo é effeito da insufficiencia ou impotencia intellectual, e por isso é finito e multiplicavel; o Verbo divino nasce da infinita perfeição e fecundidade de Deus, e por isso é infinito e unico. (Cf. S. Th. \emph{De differentia verbi divini et humani}; --- \emph{Sum. Th.} p. I, q. 34).}.

\textbf{197. A intelligencia possivel conhece no verbo a essencia das coisas.} --- O \emph{verbo}, que é o fructo ou o termo da operação \emph{intellectual}, \pg{316}  é, por sua natureza, uma imagem ou representação espiritual. Ora toda a imagem ou representação leva naturalmente ao \pg{317}  conhecimento de um objecto; aliás será uma imagem, que não se refere a nenhum exemplar, e uma representação, que

nada representa: o que é absurdo. Logo a intelligencia \emph{possivel} conhece no \emph{verbo} o objecto respectivo. E como o objecto, representado pelo \emph{verbo}, é a \emph{essencia} das coisas, segue-se que \pg{319}  a intelligencia, no verbo e pelo verbo, percebe a mesma \emph{essencia}. --- A \emph{essencia} das coisas, que a intelligencia percebe no verbo e pelo verbo, existe realmente nas coisas, ainda que \pg{320}  de um modo diverso d'aquelle, por que existe na intelligencia; pois nas coisas existe de um modo \emph{concreto}, e na intelligencia de um modo \emph{abstracto} \footnote{O \emph{verbo}, embora seja \emph{subjectivo}, emquanto reside num sujeito, comtudo é \emph{objectivo}, emquanto, representando a \emph{essencia}, exprime uma coisa objectiva e real. Desenvolvamos este ponto, que é da maxima importancia. O verbo é \emph{subjectivo}, emquanto é forma da intelligencia, e fructo de uma operação \emph{imwunanente}, a qual tem o seu principio e o seu termo no sujeito. --- Mas esta \emph{subjectividade} do verbo não torna \emph{subjectivo} o \emph{objecto}, que elle representa. O officio do verbo não é o de \emph{mudar} o objecto, mas sim o de o \emph{representar}. As qualidades, que competem ao verbo por causa do \emph{sujeito}, ao qual adhere, não se confundem com as qualidades, que lhe competem por causa do \emph{objecto}, que representa. Residindo num sujeito, o verbo tem uma \emph{entidade} proporcionada e correspondente á natureza d'esse sujeito, e por isso é increado ou creado, espiritual ou material, conforme increado ou creado, espiritual ou material é o seu sujeito. Portanto, a natureza das qualidades do verbo é especificada pela natureza do \emph{sujeito}, a que adhere, e não do \emph{objecto}, que representa; aliás deveriamos dizer --- que o verbo é material ou corporeo, porque representa coisas materiaes e corporeas, --- e que Deus, o Anjo e a mente humana não podem ter a especie do mundo corporeo, porque incapazes de receber a materia. --- Mas as qualidades, que do \emph{sujeito} provêm para o verbo, são distinctas e diversas das que lhe provêm do \emph{objecto}; assim um verbo, \emph{entitativa} ou \emph{subjectivamente} finito, pode representar um objecto \emph{infinito}, \emph{e} um verbo \emph{espiritual} pode exprimir um objecto \emph{material}. E' simplesmente absurdo pensar que a qualidade do objecto representado se torne qualidade do verbo representante, e que por isso seja material o verbo, que representa um objecto material. Diz S. Thomaz: « A forma in- \&ellectual (o verbo) tem uma \emph{duplice} relação: uma para com o \emph{objecto}, de que é forma representante, outra para com o \emph{sujeito}, de que é forma inherente. Se se considera em relação ao \emph{objecto}, não se diz que é dotada d'esta ou d'aquella qualidade, mas que representa este ou aquelle objecto; porque não se pode chamar material uma forma que representa um objecto materiali nem sensivel a que exprime uma coisa sensivel. Se se considera em\pgn{321}relação ao \emph{sujeito}, diz-se dotada d'esta ou d'aquella qualidade, porque segue a condição do sujeito, em que subsiste » (\emph{De ideis}, a. 2 ad 5). --- O que dizemos do \emph{verbo} deve applicar-se á \emph{especie intelligivel}, e tambem á \emph{idéa}, emquanto esta se emprega para significar o \emph{verbo} ou a \emph{especie intelligivel}.\par O \emph{verbo}, pois, é a imagem de um \emph{objecto}. Ora á imagem não se representa \emph{directamente} a si mesma, mas representa o \emph{objecto}. Por isso o verbo, sendo, por sua natureza, o proprio objecto emquanto representado, leva \emph{directamente} a intelligencia a perceber o objecto. E, como o objecto, representado pelo \emph{verbo}, á a \emph{essencia} das coisas, segue-se que á intelligencia, no verbo e pelo verbo, percebe a' mesma \emph{essencia}. É o que ensina S. Thomaz: « Objectum intellectus est ipsa rei \emph{essentia}; quamvis essentiam rei cognoscat per eius \emph{similitudinem}, sicut per \emph{medium cognoscendi}, non sicut objectum in quod \emph{primo} fertur eius visio » (\emph{De mente}, a. 4 ad 4). --- \emph{A essencia} das coisas, que a intelligencia percebe no verbo e pelo verbo, existe realmente nas coisas, ainda que de um modo diverso d'aquelle, por que existe na intelligencia; pois nas coisas existe de um modo \emph{concreto}, isto é, cercada de propriedades individuaes, ao passo que na intelligencia existe de um modo \emph{abstracto}, isto é, livre d'essas propriedades. Assim a \emph{humanidade}, isto é, a natureza constituida pelos dois elementos, que são á \emph{animalidade} e a \emph{racionalidade}, existe \emph{realmente} no mundo, fóra da nossa intelligencia: porque no mundo existem individuos realmente dotados de vida \emph{animal} e \emph{racional}; embora a \emph{hwumanidade}, que se encontra nos homens, seja determinada e definida por caracteres, que são proprios de um individuo e não de outro, ao passo que á \emph{Auwmanidade}, que é o objecto da nossa intelligencia, prescinde d'esses caracteres individuaes (Cf. \emph{Sum}. \emph{Th}, p. I, q. 85, a. 2 ad 2).\par D'ahi se vé a futilidade da velha objecção, resuscitada ultimamente por \emph{Kant}; á qual diz: Como á possivel que o objecto, expresso pelo verbo mental, seja o que existe na realidade, quando esse objecto, considerado na intelligencia, é \emph{universal}, \emph{abstracto}, \emph{necessario}, \emph{immutavel}, ao passo que, considerado na realidade, a \emph{singular}, \emph{concreto}, \emph{contingente}, \emph{nmutavel} ?\par Este dificuldade funda-se num supposto --- que as coisas existam na mente do mesmo modo, por que existem na realidade. Ora este supposto é falso, como vamos provar. Às coisas são conhecidas emquanto existem no cognoscente; e existem no cognoscente de um modo conforme á natureza d'este; d'ahi o principio: \emph{receptwm}. \emph{est in recipiente per qodum recipientis}. Se, pois, o modo de ser do cognoscente é superior ao modo de ser da coisa conhecida, esta existe mais perfeitamente no cognoscente do que em si mesma; assim uma pedra existe mais perfeitamente na. nossa intelligencia do que em si\par 21 Vor. II.\pgn{322}mesma, porque em si existe de um modo material e na nossa mente existe de um modo espiritual. Como, pelo contrario, se \emph{o} modo de ser do cognoscente é inferior ao modo de ser da coisa conhecida, esta existe mais perfeitamente em si mesma do que no cognoscente; assim os attributos divinos existem mais perfeitamente em Deus do que na nossa intelligencia, porque em Deus existem de um modo simplicissimo e infinito, \emph{e} na nossa mente de um modo composto e finito. Se é, portanto, certo que o objecto conhecido deve de \emph{algum modo} encontrarse no cognoscente, é falso que elle deve encontrar-se \emph{do mesmo modo} por que existe em si mesmo, porque isto repugna á propria natureza do conhecimento (\emph{Sum}. \emph{Th}., p.l, q. 12, a. 4). Não admira, pois, que as coisas, que em si mesmas são \emph{singulares} e \emph{materiaes}, se encontrem na nossa mente num estado \emph{immaterial} e \emph{abstracto}. --- Mas essa diversidade refere-se simplesmente aos \emph{modos} de existencia, que são \emph{accidentaes} e não alteram a \emph{natureza} das coisas; porque, se considerarmos os \emph{predicados essenciaes}, os \emph{elementos constitutivos} das coisas, esses \emph{predicados}, esses \emph{elementos} não só se encontram na nossa mente, mas tambem fóra d'ella, na realidade, porque é por esses predicados ou elementos que as coisas são o que são e se acham constituidas na propria especie. Portanto não repugna que um e o mesmo objecto seja \emph{universal} e \emph{singular}, \emph{abstracto} e \emph{concreto}, conforme se considera existente na mente ou na realidade. --- Devemos, porém, advertir que um e o mesmo objecto é, tanto na nossa mente como fóra d'ella, \emph{contingente} e \emph{Twutavel} em ordem á \emph{existencia}, e é \emph{necessario} e \emph{immutavel} em ordem \& \emph{essencia} (\emph{Sum}. \emph{Th}, p. I, q. 85, a. 9 ad 2).\par Os \emph{sensualistas}, sequazes de \emph{Condilac}, e os \emph{idealistas}, discipulos de \emph{Kant}, negam a \emph{objectividade} dos nossos conceitos intellectuaes, e affirmam que conhecemos apenas a \emph{modificação} da nossa inteligencia, e não os \emph{objectos} representados pelas ideas. :\par Esta opinião é \emph{falsa} e \emph{absurda}. --- E \emph{falsa} não só pelo que deixamos dito, nas tambem porque a genesis do nosso conhecimento intellectual mostra claramente a sua \emph{objectividade}. De facto, o \emph{verbo}, que é o termo do acto intellectual e que exprime a \emph{essencia} do objecto, procede, como de principio, da \emph{especie intelligivel}, \& qual é semelhante. A \emph{especie intelligivel} é produzida pela \emph{intelligencia agente}, que a tira ou abstrahe do \emph{phantasma}. O \emph{phantasma} é formado pela \emph{imaginação} com os elementos subministrados pelos \emph{sentidos externos}. Os \emph{sentidos externos} estão em contacto \emph{immediato} com os objectos, e não se enganam no seu acto perceptivo. Logo, pelo facto de o processo intellectual comecar dos objectos, e, depois de ter passado por uma continua serie de causas, terminar na intelligencia, devemos concluir que as nossas ideas são \emph{objectivas}, e que por isso percebemos os proprios objectos (Cf. S. Th., c. \emph{Gentes}, l, c. 58) --- E' \emph{absurda}, porque repugna ao \emph{senso commum}. O valor \emph{objectivo} das nossas ideas é uma verdade fundamental, evidente por\pgn{323}si mesma, e é uma condição essencial de todo o progresso scientifico. Uma idea sem valor \emph{objectivo} seria um conhecimento sem objecto conhecido, uma representação que nada representa. Esse valor não \emph{pode ser} nem \emph{demonstrado}, nem \emph{negado}, nem \emph{posto em duvida}, sem que, por isso mesmo, seja \emph{affirmado}; porque toda a demonstração, directa ou indirecta, presuppõe necessariamente principios e ideas, e seria completamente illusoria, se se prescindisse do valor \emph{objectivo} das ideas nella empregadas. Todos os homens, á excepção dos loucos, quando fallam, attribuem um valor \emph{objectivo} ás suas ideas, e entendem exprimir uma coisa \emph{objectiva}, \emph{e} não um delirio da imaginação; e os proprios adversarios, quando affirmam que as ideas não têm valor \emph{objectivo}, implicitamente admittem que a sua affirmacüo tem um valor \emph{objectivo}; aliás seria inutil e ridiculo fallar e insistir na propria opinião. E' com a permissão \textbar{} d'esses disparates que Deus castiga o orgulho de certas intelligencias! \textbar{} (1) O \emph{desenvolvimento} do conhecimento \emph{intellectual} póde considerar-se \textbar{} sob dois aspectos; \emph{subjectivo} e \emph{objectivo}. --- Considerado sob o aspecto \emph{subjectiv}0, o desenvolvimento do conhecimento \emph{intellectual} refere-se ás diversas j \emph{funcções} ou \emph{operações}, que a intelligencia produz. --- Considerado sob o aspecto\par \textasciicircum{}, \emph{objectivo}, esse desenvolvimento refere-se aos diversos \emph{objectos}, que a. mesma inteligencia percebe ou póde perceber.\par Tratando, no texto, do desenvolvimento do conhecimento \emph{intellectual}, \textbar{} emquanto \emph{objectivo}, diremos, em nota e no fim do artigo, alguma coisa d'elle, \textbar{} j emquanto \emph{subjectivo}. \textbar{} 1 (2) Tratando do \emph{objecto} da intelligencia (n. 114), provámos, com o suf- \textbar{} ficiente desenvolvimento, que a \emph{essencia} das coisas \emph{corporaes}, \emph{abstrahida} das \textbar{}}.

%%>>> art-27

\section*{Artigo III.}
\addcontentsline{toc}{section}{Artigo III.}

\pg{314} \pg{315} \pg{318} \pg{321} \textbf{Modo por} que \textbf{se desenvolve} o \textbf{conhecimento}

\textbf{intellectual}.

198. Desenvolvimento objectivo do conhecimento. --- Depois de termos considerado as funcções, que a \emph{intelligencia} e os \emph{sentidos} exercem para a producção do conhecimento \emph{intellectual}, resta-nos \pg{322}  examinar os diversos \emph{objectos}, que conhecemos, e o \emph{modo} por que os conhecemos.

Os \emph{objectos}, que a intelligencia percebe, são as substancias, tanto \emph{materiaes} como \emph{immateriaes}. As substancias \emph{materiaes} - podem considerar-se --- ou na sua \emph{universalidade}, isto é, na sua \emph{essencia}, --- ou na sua \emph{singularidade}, como existem na realidade. As substancias \emph{emunateriaes} são a \emph{alma humana} e \textsc{Deus}.

\pg{323} As substancias \emph{materiaes} epnsideradas na sua \emph{universalidade}, isto é, á \emph{essencia} das substancias \emph{materiaes} conhece-se por um modo \emph{directo}. As substancias \emph{materiaes} consideradas na sua \emph{singularidade}, como tambem as substancias \emph{immateriaes}, como a \emph{alma humana} e \textsc{Deus}, são conhecidas por um modo \emph{indirecto}. \textbar{} Desenvolveremos esta doutrina no presente artigo \footnote{\pgn{324}condições materiaes, --- é o objecto \emph{natural} ou \emph{proporcionado} á condição actual da nossa intelligencia, --- é conhecida \emph{anteriormente} aos outros objectos e pela sua especie \emph{propria}, --- e é a razão \emph{pela qual} e \emph{segundo a qual} conhecemos todos os mais objectos. Não é, pois, necessario insistir nestes pontos. Só indiearemos uma ou outra consequencia, que deriva d'essa verdade.\par Lembramos --- que todo o conhecimento \emph{intellectual} suppõe necessariamente um conhecimento \emph{sensitivo}, --- que nada pode entender a intelligencia que não tenha: sido, de algum modo, percebido e apresentado pelo sentidos, --- e que, por isso, a inteligencia só pode chegar até onde os sentidos a podem guiar. Diz S. Thomaz: « Naturalis nostra cognitio a sensu principium sumit; unde tantum se nostra naturalis cognitio extendere potest, in quantum manuduei potest per sensibilia » (\emph{Sum}. \emph{Th}. p. T, q. 12, a. 12). Ora o primeiro objecto que as coisas sensiveis ou materiaes offerecem á contemplação da intelligencia é a \emph{quiddidade} ou \emph{essencia} d'ellas: « \emph{Primum}, quod intelligitur a nobis, secundum statum praesentis vitae, est \emph{quidditas} rei \emph{materialis}, quae est nostri intellectus objectum » (\emph{Sum}. \emph{TA}., p. I, q. 88, a. 3).\par Nem todas as \emph{essencias} das substancias \emph{materiaes} são conhecidas por nós de um modo \emph{directo} ou \emph{immediato}. Muitas vezes, para distinguirmos uma substancia de outra, devemos servir-nos do complexo das propriedades \emph{accidentaes}, --- complexo que suppre a falta da \emph{differenga especifica}. Diz o Angelico: « Quia \emph{essentiales} rerum differentiae sunt ignotae frequenter et innominatae, oportet interdum uti \emph{accidentalibus differentiis} ad substantiales differentias designandas » (\emph{De pot}., a. 2 ad 2).\par Na percepção \emph{directa} das substancias \emph{materiaes}, a intelligencia começa pelos conceitos, que têm mais extensão e menos comprehensao, e chega a um conceito, que tem menos extensão e mais comprehensão e que por isso é mais proprio do objecto representado. --- Tal processo é conforme com a indole da nossa intelligencia, e corresponde á experiencia. E' conforme com a indole da nossa intelligencia, que, sendo uma faculdade \emph{potencial}, passa naturalmente do imperfeito e incompleto para o perfeito e completo; ora o conceito universal é menos distincto e menos completo que o particular. (Cf. \emph{Sum}. \emph{Th}., p. I, q. 85, a. 3). Corresponde á experiencia, a qual aítesta que os povos e os individuos pouco civilisados tém pobreza de linguagem, empregando sempre palavras genericas e indeterminadas, ao passo que as nações 2 as pessoas civilisadas têm riqueza de linguagem, empregando palavras especificas e proprias de um certo e determinado objecto; ora a palavra é a expressão oral do conceito. --- D'ahi se segue que o primeiro conceito, formado pela abstracção, é o conceito de \emph{ente}, que tem a maxima extensão e a minima comprehensão. E' neste conceito que se resolvem e é a este que se addicionam todos os outro conceitos. --- Aqui devemos advertir que,\pgn{325}quando dizemos ser o conhecimento do \emph{universal} anterior ao do \emph{singular}, nos referimos ao conhecimento \emph{intellectual}; porque, se se trata do conhecimento \emph{em geral}, o conhecimento do \emph{singular} é anterior ao do \emph{universal}, sendo o acto da intelligencia precedido necessariamente pelo acto dos sentidos, e estes conhecem o \emph{singular}.\par Concebida a idéa de \emph{ente}, a intelligencia exprime o \emph{principio de contradicção}, que é o primeiro e a base de todos os outros principios racionaes. De facto, dizendo \emph{ente}, immediatamente percebemos a sua repugnancia com o \emph{nada} (pois o \emph{nada} é a negação do ente), e enunciamos este juizo: \emph{o que é não pode deivar de ser emquanto} é. E' o \emph{primeiro juizo}, porque se refere á \emph{primeira idéa}. --- \emph{O principio de contradicção} diz-se \emph{fundamento} de todos os outros principios --- não emquanto só elle é percebido por 8i, immediatamente, --- mas emquanto está implicitamente contido em todos os outros principios. --- Os primeiros principios conhecem-se pela luz da intelligencia agente, isto é, pela força abstractiva da nossa mente, e não pelo ' raciocinio, mas pelo simples conhecimento dos termos. --- Esta facilidade e promptidao natural com que a nossa intelligencia, depois de perceber a essencia das coisas, pronuncia os \emph{primeiros principios}, chama-se \emph{habito innato} dos primeiros principios. Este habito diz-se \emph{innato}, não emquanto são \emph{innalas} as idéas, de que se compõem os principios, --- mas emquanto ó \emph{innata} a luz intellectual e a facilidade ou propensão para os perceber, quando nos são propostos. As idéas, em que se baséam os principios, são adquiridas pela virtude abstractiva da nossa intelligencia.\par Percebidas as \emph{essencias e} formados os \emph{primeiros} principios, a intelligencia reflecte nos seus actos e nos dos sentidos e percebemos a nossa alma e as substancias materiaes, e depois, pelo principio de \emph{causalidade}, subimos ao Principio increado de todas as coisas.}.

%%>>> t199-2

\tese{199}{A nossa intelligencia conhece a essencia das substancias materiaes por um modo directo.}

2) Conhece a \emph{essencia} das substancias \emph{materiaes}, Porquanto, 1 a consciencia attesta que percebemos a \emph{essencia} do \emph{mineral}, da \emph{planta}, do \emph{animal}, do \emph{homem}, etc. Ora estas substancias são \emph{materiaes}. \textbar{}

b) Por um modo \emph{directo}. Um objecto é conhecido por um modo \emph{directo}, quando é conhecido \emph{por si}, isto é, por meio da sua especie ou imagem propria. Ora a \emph{essencia} das substancias \textbar{} \emph{materiaes} é conhecida \emph{por si}; por quanto as substancias \emph{mate}-

rá \emph{riaes} são as que \emph{primeiramente}, pelas suas qualidades concretas, impressionam os \emph{sentidos}, e, pela sua \emph{essencia abstrahida} \textbar{} das propriedades individuaes, actuam na \emph{intelligencia}. Logo a nossa intelligencia conhece a \emph{essencia} das substancias \emph{materiaes} por um modo \emph{directo} (2).

%%>>> t200

\pg{324} \tese{200}{A nossa intelligencia conhece as substancias materiaes na sua singularidade por um modo indirecto.}

a) Conhece as substancias \emph{materiaes} na sua \emph{singularidade}. Se a intelligencia não conhecesse as substancias \emph{materiaes} na sua \emph{singularidade}, não poderia instituir juizos nem raciocinios ácerca d'elles. Ora instituimos juizos e raciocinios ácerca de semelhantes objectos. Logo a nossa intelligencia conhece as substancias \emph{materiaes} na sua \emph{singularidade}.

b) Por um modo \emph{indirecto}. Um objecto conhece-se por um modo \emph{indirecto}, quando não é percebido pela especie \emph{propria}, mas por meio de outro objecto, com que elle está unido. Ora as substancias \emph{materiaes}, consideradas na sua \emph{singularidade}, são percebidas pela intelligencia, não pela especie \emph{propria}, mas por meio do \emph{phantasma}, que as representa. Logo a nossa intelligencia conhece as substancias \emph{materiaes} na sua \emph{singularidade} por um modo \emph{indirecto} \footnote{Pergunta-se: Como é que a intelligencia póde perceber as substancias \emph{materiaes} na sua \emph{singularidade}, quando o objecto d'aquella faculdade é o \emph{universal}, o \emph{immaterial}?\par Desenvolvamos a resposta a esta pergunta, dada resumidamente no texto. E' certo que a nossa intelligencia percebe as substancias \emph{materiaes} na sua \emph{singularidade}. A negação d'esse conhecimento repugna á excellencia d'esta faculdade, e aos officios que ella desempenha no homem. A nossa vida racional, quer na ordem especulativa quer na pratica, funda-se toda no conhecimento intellectual que temos do mundo corporeo, não só \emph{em abstracto}, mas tambem \emph{em concreto}. --- Além d'isso, a nossa intelligencia fórma proposições ou juizos relativos ás substancias \emph{singulares} e \emph{materiaes}, attribuindo-lhes predicados \emph{communs}, que exprimem uma natureza ou essencia \emph{commum}. Ora taes juizos seriam impossiveis, se uma e a mesma faculdade não\pgn{326}percebesse ambos os termos, isto é, o sujeito \emph{singular} ou \emph{material} e o predicado \emph{commum}, assim como não poderiamos perceber a differença entre o sabor doce e a côr branca, se não existisse uma faculdade sensitiva commum, capaz de perceber uma e outra qualidade. Diz S. Thomaz: « Sicut non possemus sentire differentiam dulcis et albi, nisi esset una potentia sensitiva communis, quae cognosceret utrumque, ita etiam non possemus cognoscere comparationem \emph{universalis} ad \emph{particulare}, nisi esset \emph{una} potentia, quae cognosceret utrumque. Intellectus igitur utrumque cognoscit, sed \emph{alio} et \emph{alio modo} ». (\emph{De anima}, lib. III, lect. 8). --- Vamos indicar os \emph{diversos} modos, por que a intelligencia percebe o \emph{universal} e o \emph{singular}.\par O conhecimento, que a intelligencia tem do \emph{universal}, \emph{é directo}. Porquanto, o conhecimento de um objecto é \emph{directo}, quando o objecto se conhece pela sua \emph{especie} ou \emph{imagem propria}. Ora o \emph{universal} conhece-se pela sua \emph{especie propria}, que é abstrahida do phantasma pela luz da intelligencia \emph{agente}. --- Mas o \emph{singular} ou \emph{material}, se pode ser conhecido na sua especie \emph{propria} pelas faculdades \emph{sensitivas}, não o pode ser pela intelligencia. Na verdade, semelhante especie ou imagem, representativa do objecto e necessaria para a determinação da intelligencia, --- não pode ser formada pela acção do proprio objecto na intelligencia, porque um objecto \emph{singular} e \emph{material}, não pode actuar numa faculdade espiritual, como é a nossa intelligencia, --- nem pode ser formada pelo mesmo objecto, \emph{elevado pela intelligencia agente}, como o é o phantasma, porque então produzir-se-hia uma \emph{especie intelligivel}, e esta não representaria o objecto \emph{singular} e \emph{material}, mas a \emph{essencia universal}, isto é, \emph{abstrahida} de todas as notas ou condições individuaes. Logo o conhecimento, que a intelligencia tem do objecto \emph{singular} e \emph{material}, não é \emph{directo}.\par Do exposto resulta que a razão proxima e adequada, pela qual não podemos perceber \emph{directamente} o \emph{singular corporeo}, não se deve procurar na \emph{singularidade} d'esse objecto, nem na \emph{immaterialidade} da especie intelligivel, mas unicamente no modo connatural por que a nossa intelligencia agora entende. Deus percebe o \emph{singular corporeo}, porque, pela sua intelligencia, produz a \emph{totalidade} de uma coisa, isto é, produz não só a \emph{forma}, que é principio de especificação, mas tambem a \emph{materia}, que é principio de individuação, e, por isso, nas suas idéas archetypas, conhece as coisas em \emph{todos} os seus elementos, não só especificos e formaes mas tambem individuaes e materiaes. As intelligencias separadas tambem percebem o \emph{singular corporeo}, porque recebem de Deus as especies intelligiveis, semelhantes ás idéas archetypas divinas, que, como dissemos, representam \emph{todos} os elementos das coisas. Mas a nossa intelligencia, no estado actual da união da alma com o corpo, não pode perceber senão por meio das \emph{especies intelligivis}, \emph{abstrahidas} de todas as condições materiaes ou individuantes; e, como estas especies representam apenas o \emph{universal}, isto é, a essencia \emph{abstracta}, e não as notas \emph{individuantes} e \emph{materiaes}, é claro que a intelligencia, por meio d'estas especies, só pode perceber o \emph{universal}, e não o \emph{singular}. Logo, se não podemos perceber \emph{directamente}\pgn{327} o \emph{singular}, isto não depende da \emph{singularidade} d'esse objecto, nem da \emph{immaterialidade} da especie intelligivel, aliás nem Deus nem os Anjós poderiam conhecer o \emph{singular}, mas depende unicamente do modo connatural, porque a nossa intelligencia entende no estado actual da união da alma com o corpo (Cf. S. Th., \emph{De anima}, a. 20).\par Se o conhecimento, que a intelligencia tem dos objectos \emph{singulares} e \emph{materiaes}, não é \emph{directo}, deve ser e é \emph{indirecto}. O conhecimento de uma coisa é \emph{indirecto}, quando a coisa não é conhecida \emph{por si}, isto é, pela \emph{especie} ou \emph{imagem propria}, mas por meio de uma \emph{outra} coisa, com que está unida e que é conhecida pela \emph{especie propria}. Ora o objecto \emph{singular} e \emph{material} não é conhecido \emph{por si}, mas por meio do \emph{universal}, com que elle está naturalmente unido. Na verdade, a nossa intelligencia percebe, de um modo \emph{directo}, o \emph{universal}, mas percebe-o pela \emph{abstracção}, emquanto no individuo singular e material, representado pelo phantasma, considera, pela virtude da intelligencia agente, a \emph{essencia}, sem considerar os principios \emph{individuantes}. Ora é impossivel que a intelligencia considere e apprehenda no individuo singular e material a \emph{essencia}, abstrahida dos principios individuantes e por isso \emph{universal}, sem que, ao mesmo tempo e de um modo \emph{indirecto} e \emph{accidental}, considere e apprehenda o mesmo individuo singular e material, em que está contida essa \emph{essencia}, ou esse \emph{universal}; assim como é impossivel que a vista perceba uma certa côr, sem que, de um modo indirecto e accidental, perceba o sujeito, em que a côr subsiste. Logo a intelligencia conhece o objecto \emph{singular} e \emph{material} por meio do \emph{universal}, e por isso de um modo \emph{indirecto}.\par Tal é a doutrina de S. Thomaz. O principio, diz elle, que torna \emph{singulares} as coisas materiaes, é a materia individual. Ora a nossa intelligencia entende emquanto abstrahe da materia individual a especie intelligivel. O que se abstrahe da materia individual é \emph{universal}. Por isso, a nossa intelligencia não conhece \emph{directamente} senão o \emph{universal}; mas \emph{indirectamente}, \emph{e por uma especie de reflexão}, pode conhecer o \emph{singular}. Por quanto, mesmo depois de abstrahir as especies intelligiveis, ella não pode por meio das ditas especies \emph{actualmente} perceber, sem que olhe para os phantasmas, nos quaes entende as especies intelligiveis. D'este modo, a intelligencia percebe \emph{directamente}, pela especie intelligivel, o \emph{universal}, e \emph{indirectamente} percebe o \emph{singular}, representado pelo phantasma. Eis as palavras do S. Doutor: « Principium \emph{singularitatis} in rebus materialibus est \emph{materia individualis}. Intellectus autem noster intelligit abstrahendo speciem intelligibilem ab huiusmodi materia. Quod autem a materia individuali abstrahitur, est \emph{universale}. Unde intellectus noster \emph{directe} non est cognoscitivus, nisi \emph{universalium}. \emph{Indirecte} autem, \emph{et quasi per quandam reflexionem}, potest cognoscere \emph{singulare}. Quia, sicut supra dictum est (p. I, q. 84, a. 7), etiam postquam species intelligibiles abstraxerit, non potest secundum eas \emph{actu} intelligere, \emph{nisi convertendo se ad phantasmata}, \emph{in quibus species intelligibiles intendit}, ut dicitur in III \emph{de anima}. Sic igitur ipsum \emph{universale} per speciem intelligibilem \emph{directe} intelligit; \emph{indirecte} autem \emph{singularia}, quorum sunt phantasmata. Et hoc modo format propositionem: \emph{Socrates est homo} » (\emph{Summ. Theol.}, p. I, q. 86, a. 1). Notemos que, segundo S. Thomaz, a intelligencia percebe o \emph{singular}, --- não pela \emph{reflexão} propriamente dita, porque a intelligencia, neste caso, não reflecte\pgn{328}sobre si mesma, --- mas por \emph{uma especie} de reflexão, porque a intelligencia reflecte \emph{sobre} o \emph{phantasma}, e por isso o conhecimento é \emph{indirecto}, \emph{per accidens}. --- Procuremos explicar o que é \emph{esta especie de reflexão}, que a intelligencia exerce sobre o phantasma e pela qual chega a perceber o \emph{singular}.\par Ha uma grande distancia entre a intelligencia e os sentidos; porque ha uma grande distancia entre o \emph{universal}, que é o objecto da intelligencia, e o \emph{singular}, que é o objecto dos sentidos. Todavia, uma distancia, embora grande, não é uma separação. Todas as faculdades humanas, não só as intellectuaes mas tambem as sensitivas, brotam da mesma essencia da alma e nella subsistem, como tambem o \emph{universal} está encerrado, embora em estado potencial, no \emph{singular}. Brotando da mesma essencia, as faculdades humanas devem estar subordinadas, de modo que as inferiores sirvam ás superiores; e, estando o \emph{universal} contido no \emph{singular}, o conhecimento de um deve levar ao conhecimento do outro. E nada mais natural. Entre os actos da intelligencia e os das faculdades sensitivas existe tal continuação e connexão, que o que se passa nos sentidos deve ter uma repercussão, um reflexo na intelligencia, e vice-versa. --- Mas ha uma differença. Sendo a energia sensitiva inferior á intellectual, o conhecimento que a intelligencia tem do que se passa na esphera dos sentidos excede, em extensão e perfeição, o conhecimento que os sentidos têm do que se passa na intelligencia. Diz S. Thomaz: « Virtus sensitiva est infra virtutem intellectivam; unde magis potest intellectus se extendere ad id quod est sensus, quam sensus ad id quod est intellectus » (IV \emph{Sent.}, Dist. 50, a. 3 ad 4).\par E', pois, por meio dos sentidos que a intelligencia deve adquirir e adquire o conhecimento do \emph{singular}. --- Mas os sentidos auxiliam a intelligencia no conhecimento do \emph{singular} --- não emquanto apresentam o phantasma, do qual se possa abstrahir a \emph{especie intelligivel}, porque esta especie representaria o \emph{universal}, e não o \emph{singular}, --- mas emquanto apresentam o phantasma, no qual a intelligencia, pela reflexão nos seus actos, percebe de algum modo o \emph{singular}. --- Na verdade, é certo que a nossa intelligencia não só entende, mas entende que entende, isto é, tem uma percepção intuitiva e immediata dos seus actos concretos e individuaes. Para esta percepção dos seus actos, não precisa a intelligencia de novas imagens ou de especies distinctas, representativas dos proprios actos (ainda que essas especies ou imagens possam formar-se em seguida, recordando aquella percepção consciente e considerando exclusivamente o elemento subjectivo); basta que ella não concentre toda a sua energia na percepção do objecto, mas a extenda até ao seu acto. D'este modo a intelligencia, sem producção de novas imagens e por uma simples elasticidade do seu acto, percebe o objecto e o proprio acto, isto é, percebe alguma coisa a mais que não está representada na \emph{especie intelligivel}. --- Percebendo o proprio acto, a intelligencia deve perceber o \emph{principio}, de que o acto deriva e que é a \emph{especie intelligivel}. Percebendo a \emph{especie intelligivel}, percebe a \emph{origem} d'ella, isto é, percebe que a especie não é innata, não foi infundida do alto, mas foi abstrahida do \emph{phantasma}, representativo do objecto entendido. Percebendo que a especie foi abstrahida do \emph{phantasma}, a intelligencia percebe o proprio \emph{phantasma}, que reside, como no seu sujeito proprio, na faculdade da \emph{imaginação}. E como o\pgn{329}phantasma representa o \emph{material} e o \emph{singular}, isto é, a essencia emquanto está envolvida nas propriedades individuaes, a intelligencia percebe o \emph{singular} e o \emph{material}. Logo a intelligencia, pelo facto de perceber o seu acto, percebe, \emph{de algum modo}, a \emph{singularidade} do objecto entendido. --- Dissemos \emph{de algum modo}. Por quanto, a intelligencia percebe o \emph{singular} --- não emquanto forma uma \emph{especie intelligivel}, representativa das differenças e notas individuaes ou singulares, e por isso a sua percepção não é \emph{directa}, --- mas emquanto, na consciencia do proprio acto, opera uma especie de \emph{reflexão} sobre o \emph{phantasma}, e por isso esta sua percepção é \emph{indirecta}; porque a percepção de uma faculdade é \emph{indirecta}, quando se refere a um objecto, percebido por uma outra faculdade, com a qual a primeira está unida.\par Esta explicação é muito natural. A \emph{especie intelligivel} e o \emph{phantasma} são duas coisas tão estreitamente relacionadas, que uma chama a outra, uma leva á percepção da outra, embora de um modo differente. Por quanto, a \emph{especie intelligivel} é effeito de uma \emph{paixão}; e esta paixão é produzida pela acção do phantasma, que é elevado, pela intelligencia \emph{agente}, a operar na propria intelligencia \emph{possivel} e a imprimir nella a dita \emph{especie}. Por isso a percepção da \emph{especie intelligivel} leva á percepção da dita \emph{paixão}, e esta á percepção da sua causa, que é o \emph{phantasma}. --- Além d'isso, a intelligencia, para entender, precisa sempre do apoio do phantasma. Na verdade, esta faculdade, como dissemos com S. Thomaz, mesmo depois de ter abstrahido as \emph{especies intelligiveis}, não pode actualmente entender por meio d'ellas, isto é, não pode considerar a essencia \emph{abstracta}, sem se virar de novo para o phantasma, em que vê resplandecer o seu intelligivel, que é essa mesma essencia. E, embora a intelligencia não perceba \emph{directamente} senão a essencia \emph{abstracta}, comtudo, como a essencia está contida no \emph{phantasma}, é impossivel que ella veja a essencia, sem ver, ao mesmo tempo e de algum modo, o phantasma.\par Esta doutrina é de S. Thomaz. Diz o S. Doutor: « Dicendum est quod anima coniuncta corpori per \emph{intellectum} cognoscit \emph{singulare} --- \emph{non quidem directe}, \emph{sed per quandam reflexionem}, inquantum scilicet ex hoc quod apprehendit suum intelligibile, \emph{revertitur} ad considerandum suum \emph{actum}, et \emph{speciem intelligibilem} quae est \emph{principium} suae operationis, et eius speciei \emph{originem}, et \emph{sic} venit in considerationem \emph{phantasmatum}, et \emph{singularium}, quorum sunt phantasmata. Sed reflexio haec compleri non potest nisi per adiunctionem virtutis \emph{cogitativae} et \emph{imaginativae} » (\emph{De anima}, a. 20 ad 1).\par Por isso o conhecimento do \emph{singular} é um conhecimento \emph{complexo}, constituido de um termo \emph{commum} e de um \emph{prenome demonstrativo}; por ex. \emph{este homem}. O termo \emph{commum}, emquanto tal, exprime um conceito \emph{universal} --- a \emph{essencia}, e é percebido \emph{directamente} pela intelligencia; o \emph{prenome demonstrativo} significa o \emph{sujeito}, no qual se determina e individualisa a essencia, e é percebido pela mesma intelligencia, mas de um modo \emph{indirecto}, isto é, pela reflexão sobre o phantasma. --- Na percepção do \emph{singular}, a intelligencia acceita e repete a lição, que lhe dão os sentidos, embora considere o \emph{singular} --- não como um simples facto (o que é proprio dos sentidos), mas como uma \emph{coisa}, que concretisa a forma \emph{abstracta}, contemplada na idéa \emph{universal}; visto que a intelligencia, percebendo o \emph{singular}, percebe-o depois de ter percebido o \emph{universal}, que representa a \emph{essencia}, e então é muito natural\pgn{330} que a luz d'essa idéa \emph{universal} se reflicta sobre o \emph{singular}, e que este se mostre á intelligencia como a realisação concreta d'aquella essencia.\par Concluimos --- que o homem conhece o \emph{singular corporeo} por dois modos: \emph{directamente} por meio dos sentidos, e \emph{indirectamente} por meio da intelligencia, emquanto reflecte sobre os sentidos, --- e que a mesma intelligencia, percebendo o \emph{universal} de um modo \emph{directo}, e o \emph{singular} de um modo \emph{indirecto}, pode, quando julga, unil-os, reconhecendo no \emph{singular} uma realisação concreta do \emph{universal} (\emph{Cf.} \emph{De anima}, lib. III, l. 8).}.

%%>>> t201

\pg{325} \pg{326} \tese{201}{A nossa intelligencia conhece a existencia e a essencia da alma humana por um modo indirecto.}

a) Conhece a \emph{existencia} e a \emph{essencia} da alma. Por quanto o pensamento, --- por ser um phenomeno, demonstra que \emph{existe} \pg{327}  a sua causa, que é a alma, --- e, por ser um phenomeno \emph{espiritual}, demonstra que a alma é \emph{essencialmente espiritual}.

b) Por um modo \emph{indirecto}. Por quanto uma coisa é conhecida por modo \emph{indirecto}, quando não se vê em si mesma, \pg{329}  mas numa outra coisa. Ora a \emph{existencia} e a \emph{essencia} da alma humana não são conhecidas em si mesmas, mas na \emph{existencia} \pg{330}  e \emph{natureza} das operações \emph{intellectuaes}. Logo a nossa intelligencia conhece a \emph{existencia} e a \emph{essencia} da \emph{alma humana} por um modo \emph{indirecto} \footnote{O meio, pois, por que e em que conhecemos a \emph{existencia} e a \emph{essencia} da alma, é um e o mesmo --- o \emph{pensamento}. \emph{A existencia} do pensamento demonstra a \emph{existencia} da alma; assim como a \emph{natureza} d'aquelle manifesta a \emph{natureza} d'esta. Para a percepção da \emph{existencia} da alma, basta a sua \emph{presença} e \emph{operação}; para a percepção da sua \emph{natureza} é necessario o \emph{raciocinio}.\par Sejamos mais desenvolvidos.\par a) A nossa intelligencia percebe a \emph{existencia} da alma humana. --- Esta percepção --- não é \emph{directa}; porque o objecto \emph{proprio} e \emph{proporcionado} da inteligencia humana é a \emph{essencia contida} no \emph{phantasma} ou \emph{imagem sensivel}, e não ha phantasma ou imagem sensivel que encerre a nossa alma, porque esta não póde impressionar os nossos sentidos, --- mas é \emph{indirecta}; pois conhecemol-a emquanto é principio das \emph{operações intellectuaes}.\par Esta percepção da \emph{existencia} da alma não é \emph{discursiva}, emquanto da existencia das operações, posta como premissa, se deduza, como conclusão, a \emph{existencia} da alma; porque haveria \emph{petição de principio}, pois que, quem diz que \emph{pensa}, já afirma que nelle existe uma \emph{causa}, de que deriva e em que existe o \emph{pensamento}. Logo a percepção, que a intelligencia tem da \emph{existencia} da alma, é \emph{intuitiva} ou \emph{immediata}. --- Para esta percepção, basta a essencia da propria alma, intimamente \emph{presente} á intelligencia; porque d'ella derivam as operações, em que e por que é percebida. Diz S. Thomaz: « Quantum ad actualem cognitionem, qua aliquis considerat se in actu animam habere, sic dico quod anima cognoscitur per \emph{actus suos}. In hoc enim aliquis percipit animam habere et vivere et esse, quod percipit se sentire et intelligere et alia hujusmodi vitae opera exercere: unde dicit philosophus (\emph{Eth}. l. 9, c. 9): \emph{Sentimus autem quoniam sentimus}, \emph{et intelligimus quoniam intelligimus}; \emph{et quia hoc sentimus}, \emph{intelligimus quoniam sumus}. Nullus autem percipit se intelligere, nisi ex hoc quod intelligit; quia prius est intelligere aliquid, quam intelligere se intelligere; et ideo pervenit anima ad actualiter percipiendum se esse, per illud quod intelligit, vel sentit » (\emph{De ver}., q. 10, a. 8; Cf. \emph{De mente}, a. 8). --- Por isso, quando se diz que a alma se conhece \emph{habitualmente} a si mesma pela \emph{propria essencia}, não se entende que ella\pgn{331}tenha uma perenne percepção de si mesma, mas deve entender-se que a essencia da alma faz o officio de \emph{habito}, quando ella --- a alma --- quer chegar ao conhecimento de si mesma; embora de facto não chegue a esse conhecimento senão depois de ter percebido alguma coisa, distincta d'ella.\par b) A \emph{presença} da alma, se basta para a percepção da sua \emph{existencia}, não pôde bastar para a percepção da sua \emph{natureza}. Para isso é necessario fazer primeiramente uma subtilissima analyse das operações \emph{intellectuaes}, e em seguida, tendo reconhecido que estas operações são \emph{espirituaes}, devemos, pela applicação do \emph{axioma de causalidade}, concluir que o principio, de que ellas brotam, não pode deixar de ser \emph{espiritual}. Diz S. Thomaz: « Est alius intellectus, scilicet humanus, qui nec est suum intelligere (sicut divinus), nec sui intelligere est objectum primum eius essentia (sicut angelicus), sed aliquid extrinsecum, scilicet \emph{natura materialis rei}. Et ideo id quod \emph{primo} cognoscitur ab intellectu humano est \emph{huiusmodi objectum}, et \emph{secundario} cognoscitur \emph{ipse actus}, quo cognoscitur objectum; et per \emph{actum} cognoscitur \emph{ipse intellectus}, cujus est perfectio ipsum intelligere » (\emph{Sum}. \emph{Theol}., p. I, q. 87, a. 3). E accrescenta: « Anima humana intelligit seipsam per suum intelligere, quod est actus proprius ejus, perfecte demonstrans \emph{virtutem} ejus et naturam » (\emph{Ib}., q. 88, a. 2 ad 3).\par A. \emph{Comte}, seguido pelos \emph{positivistas}, diz que a intelligencia não póde nem deve occupar-se do conhecimento da \emph{alma humana}: porque esse conhecimento pertence ao \emph{estado metaphysico}, que já passou e que é \emph{antiscientifico}. --- Com os \emph{positivistas} concordam os \emph{agnosticos}.\par Esta objecção contém muitas falsidades. --- E' falso que a intelligencia não possa conhecer a nossa alma; pois da consciencia das nossas operações nos elevamos ao conhecimento do principio, de que ellas derivam e que é a alma. --- E' falso que a intelligencia não deva occupar-se do conhecimento da alma; porque a intelligencia procura as causas dos phenomenos, ora a causa dos phenomenos, que se realizam em nós, é a alma. --- E' falso que o \emph{estado metaphysico} já passou; pois o espirito humano tem e conserva sempre uma tendencia natural para subir dos phenomenos ás causas. --- E' falso que o mesmo estado \emph{metaphysico} seja \emph{antiscientifico}; pois, como diz Cl. Bernard, « \emph{na metaphysica está a propria essencia da nossa intelligencia}, \emph{não podemos fallar senão metaphysicamente} » (\emph{Phenomenos da vida}).}.

%%>>> t202

\pg{328} \tese{202}{A nossa intelligencia conhece Deus por um modo indirecto.}

a) Conhece \textsc{Deus}. E' verdade tão evidente, que não precisa de ser demonstrada. --- O conhecimento, que a nossa intelligencia tem de \textsc{Deus}, é \emph{inadequado}, mas é \emph{verdadeiro} conhecimento; aliás não poderiamos fallar, como fallamos, da \emph{existencia} de \textsc{Deus}, da sua \emph{natureza}, \emph{attributos}.

b) Conhece por um modo \emph{indirecto}. A nossa intelligencia conhece \textsc{Deus} por um modo \emph{indirecto}; porque não conhecemos \textsc{Deus} em si mesmo, mas nas coisas creadas, emquanto dos entes contingentes, relativos, finitos, subimos ao conhecimento do Ente necessario, absoluto, infinito \footnote{A nossa intelligencia conhece \textsc{DEus} por um modo \emph{indirecto}.\par Pelo facto de a intelligencia depender \emph{objectivamente} das faculdades sensitivas, é claro que, \emph{directamente}, só podemos conhecer as coisas, que nos são apresentadas pelos seres sensiveis e materiaes.\par Ora a primeira coisa, que os seres sensiveis e materiaes apresentam á nossa intelligencia, é a \emph{essencia} d'elles. Da percepção das \emph{essencias} das coisas materiaes derivam os juizos \emph{universaes} de evidencia \emph{immediata}, que são os \emph{principios da razão}. --- Reflectindo sobre as nossas operações, percebemos a nossa existencia e a dos outros individuos, distinctos de nós.\par Este é o começo do desenvolvimento \emph{objectivo} da percepção \emph{intellectual}\pgn{332}\textsc{Deus}, pois, não é um objecto \emph{directamente} conhecido; mas elevamo-nos ao conhecimento d'Elle pelo conhecimento das coisas creadas.\par O primeiro conceito, que a intelligencia forma, pela consideração das creaturas, ácerca de \textsc{Deus}, é o da sua \emph{existencia}. Por quanto, considerando as creaturas, e vendo-as mutaveis e contingentes, subimos, em virtude \emph{do principio de causalidade}, ao conhecimento de um Ente immutavel e necessario, \emph{Causa primeira} de todos os seres finitos (\emph{Sum}. \emph{Th}. p. I, q. 2).\par Este conhecimento de \textsc{Deus}, como \emph{Causa primeira} de todas as coisas finitas, contém todos os outros conceitos, que podemos formar do Ente infinito. Porquanto, devendo toda a bondade do effeito conter-se, de um modo mais perfeito e elevado, na sua causa, é certo que toda a bondade das creaturas deve existir necessariamente em \textsc{Deus}.\par D'hai os tres processos por que nos elevamos ao conhecimento das perfeições divinas --- o processo de \emph{causalidade}, de \emph{exclusão} e de \emph{eminencia}.\par Pelo processo de \emph{causalidade}, conhecemos que \textsc{Deus} é causa não só \emph{efficiente}, mas tambem \emph{final} e \emph{exemplar}; pois repugna que a \emph{Causa primeira} procure fóra de sua essencia infinita o fim e o modelo das suas obras. --- D'ahi os \emph{attributos} divinos, relativos a este triplice genero de \emph{causalidade}, como os de \emph{conservador}, \emph{ordenador}, \emph{providente}, etc.\par Pelo processo de \emph{exclusão}, afastamos de \textsc{Deus} todas as imperfeições, proprias das creaturas, e todas as qualidades, que no seu conceito indicam imperfeição; como são a \emph{limitação}, a \emph{dependencia}, a \emph{mutabilidade}, a \emph{corruptibilidade}, a \emph{extensão}, a \emph{composição}, etc.\par Pelo processo de \emph{eminencia}, attribuimos a \textsc{Deus} todas as perfeições, que no seu conceito não indicam defeito, declarando que em \textsc{Deus} se encontram n'um grau infinitamente mais nobre e elevado, do que se encontram nas creaturas. Assim, quando dizemos que \textsc{Deus} é \emph{bom}, é \emph{bello}, é \emph{amavel}, etc., indicamos que a \emph{bondade}, a \emph{belleza}, a \emph{amabilidade}, que nas creaturas se encontram misturadas com imperfeições, em \textsc{Deus} se encontram num grau infinito e sem sombra de imperfeição (\emph{Sum}. \emph{Th}., p. I, q. 12).\par Este conhecimento de \textsc{Deus}, que adquirimos pela consideração das creaturas, é \emph{inadequado}, pois as creaturas finitas não podem reflectir toda a infinidade da perfeição do Creador, mas não é \emph{falso} e é sufficiente para excitar no nosso espirito sentimentos de admiração, de estima e de amor.\par Os \emph{positivistas} e os \emph{agnosticos} negam á nossa intelligencia o direito de se elevar, pelo conhecimento das creaturas, ao conhecimento do Creador.\pgn{333}--- Não insistimos na refutação d'esta theoria, que se basêa na negação do \emph{principio de causalidade}, --- negação que é contraria ao \emph{senso commum}.\par *\par Depois de termos fallado no desenvolvimento \emph{objectivo} do conhecimento intellectual, fallemos agora, brevemente, no seu desenvolvimento \emph{subjectivo}, isto é, nas diversas \emph{funcções}, que a intelligencia produz.\par As principaes \emph{funcções}, que a intelligencia exerce e a que as outras se reduzem, são \emph{tres}, a saber: a \emph{simples apprehensão}, o \emph{juizo} e o \emph{raciocinio}.\par a) A \emph{simples apprehensão}, que se chama tambem \emph{idea} ou \emph{conceito}, é a operação pela qual a intelligencia percebe uma coisa, sem comtudo lhe attribuir ou negar alguma propriedade ou attributo.\par A \emph{simples apprehensão} tem muitas e diversas manifestações. As principaes são a \emph{attenção}, a \emph{abstracção}, a \emph{generalisação}, a \emph{comparação} e a \emph{reflexão}.\par \emph{a}) \emph{A attenção} é a concentração da intelligencia sobre um objecto, para melhor o conhecer. --- Esta operação, embora dependente da vontade, pertence á intelligencia, porque é esta faculdade que conhece, e não aquella.\par \emph{b}) A \emph{abstracção} dá-se, quando se considera \emph{isoladamente} uma propriedade ou elemento, separando-o \emph{mentalmente} ou do sujeito a que pertence, ou das outras propriedades de que \emph{na realidade} é inseparavel. --- A \emph{abstracção} é a propria \emph{attenção}, concentrada, não sobre um \emph{objecto}, mas sobre um dos modos do objecto, com a exclusão dos outros, ou sobre a substancia, com a exclusão dos modos. --- A \emph{abstracção} admitte differentes graus. Por quanto, ou podemos prescindir da materia \emph{individual}, e só considerar a \emph{materia commum}, como faz a \emph{physica}; ou prescindimos tambem da materia \emph{commum}, e só consideramos a \emph{quantidade}, como na \emph{mathematica}; ou prescindimos da materia \emph{commum} e da \emph{quantidade}, e consideramos unicamente a \emph{essencia} das coisas, como na \emph{metaphysica}.\par \emph{c}) A \emph{generalisação} consiste em formar uma noção ou idea que, embora seja una, comtudo representa \emph{muitas} coisas. --- A \emph{generalisação} é um effeito da \emph{abstracção}, ainda que muitos escriptores cheguem a identificar uma com outra. --- A \emph{generalisação} admitte os mesmos graus, que a \emph{abstracção}.\par \emph{d}) \emph{A comparação} é uma operação da intelligencia, pela qual approximamos dois objectos ou duas ideas, para conhecermos as suas relações. Suppõe varios actos de \emph{attenção}, pelos quaes nos representamos dois objectos ou duas ideas. --- A \emph{comparação}, pois, não é a \emph{percepção} das relações: porque muitas vezes não chega a descobril-as.\par \emph{e}) A \emph{reflexão} é a operação, pela qual tornamos a pensar num objecto ou num acto já conhecido. --- Por isso a \emph{reflexão} pode ser \emph{ontologica}, ou \emph{psychologica}, conforme o seu termo é o \emph{objecto} do conhecimento \emph{directo}, ou é \emph{acto} do sujeito pensante.\par Estes actos são modificações da mesma simples \emph{apprehensão}; pois que entre elles existe apenas uma differença \emph{accidental}.\par \emph{b}) O \emph{juizo} é uma operação, pela qual affirmamos que dois objectos percebidos convêm, ou não, entre si. E' uma apprehensão \emph{complexa}. --- Na \emph{Logica} tratamos da \emph{natureza}, dos \emph{elementos} e das \emph{especies} do \emph{juizo}.\par \emph{Descartes}, \emph{Mallebranche}, seguidos por alguns escriptores modernos, \emph{sustentam}\pgn{334} que o \emph{juizo} não é operação da intelligencia, mais sim da vontade. --- Esta opinião é falsa. Por quanto, a materia do \emph{juizo} é a verdade, e a verdade não é objecto da vontade, mas da intelligencia. --- Nem se diga que o \emph{juizo} consiste no \emph{assenso}, e o \emph{assenso} é proprio da vontade. Com effeito, se é certo que o \emph{juizo} consiste no \emph{assenso}, é falso que o assenso seja proprio da \emph{vontade}. O \emph{assenso} é proprio da intelligencia, porque o acto d'esta faculdade não é precedido por outro acto espiritual. Da vontade é proprio o \emph{consenso}; porque, sendo o acto da vontade precedido necessariamente pelo da intelligencia, aquella faculdade \emph{consente} (\emph{sentit cum}) com esta.\par \emph{Reid} e \emph{Cousin} afirmam que ha juizos \emph{primitivos} ou \emph{instinctivos}, que se formam sem uma previa comparação de ideas. --- Esta opinião tambem é falsa e contraria á natureza do homem. Pelo \emph{juizo} attribuimos ou negamos a um certo sujeito um certo attributo. Ora não podemos attribuir ou negar a um sujeito um attributo, sem que anteriormente sejam percebidos e comparados entre si ambos os termos.\par \emph{Rosmini} pretende que alguns \emph{juizos} sejam ou possam ser, ao mesmo tempo, \emph{instinctivos} e \emph{comparativos}. --- E' evidente o erro d'esse escriptor; pois que o instincto exclue a comparação, e vice-versa.\par c) O \emph{raciocinio} é um acto da intelligencia, pelo qual, depois de termos comparado dois termos com um terceiro, concluimos que elles convêm, ou não, entre si. --- Não nos demoramos na exposição da doutrina ácerca do \emph{raciocinio}; pois que já o fizemos na \emph{Logica}.\par *\par A frequente repetição das operações da \emph{intelligencia} produz os \emph{habitos intellectuaes}. --- Os principaes habitos \emph{intellectuaes} são \emph{cinco}, a saber: \emph{intelligencia}, \emph{sabedoria}, \emph{sciencia}, \emph{prudencia} e \emph{arte}.\par \emph{A intelligencia}, que é uma coisa distincta da faculdade \emph{homonyma}, é o habito de conhecermos immediatamente os \emph{principios da razão}. --- A \emph{sabedoria} é o habito de conhecermos as coisas pelas suas \emph{causas supremas}. --- A \emph{sciencia} é o habito de conhecermos facilmente as consequencias nos seus principios, ou mais brevemente: é o conhecimento das coisas pelas suas causas \emph{proximas}. --- A \emph{prudencia} é o habito de conhecermos o que devemos fazer. --- A \emph{arte} é o habito de conhecermos o modo por que uma coisa deve ser feita.\par A \emph{intelligencia}, a \emph{sabedoria} e a \emph{sciencia} dizem-se habitos \emph{especulativos}; a \emph{prudencia} e a \emph{arte} chamam-se habitos \emph{practicos}.\par Os \emph{evolucionistas} admittem que todos os \emph{habitos}, \emph{intellectuaes} e \emph{moraes}, se transmittem pela \emph{hereditariedade}.\par Esta opinião é contraria á experiencia e á razão. --- E' contraria á experiencia. Porquanto a experiencia attesta que, pela \emph{hereditariedade}, são transmittidas, com a natureza, as faculdades organicas, o temperamento, e certas \emph{predisposições} intellectuaes ou moraes, mas não os habitos intellectuaes ou moraes, como a \emph{sciencia}, a \emph{virtude}, etc. --- E' contraria á razão. Com effeito, para que os \emph{habitos} podessem transmittir-se, seria necessario que a alma do filho fosse parte ou effeito da alma do pae, ou que os habitos fossem outros tantos \emph{instinctos}. Ora ambas estas hypotheses são absurdas; pois que nem a alma do filho é parte ou effeito da alma do pae, nem os habitos são \emph{instinctos}. Logo os \emph{habitos} não se transmittem por \emph{hereditariedade}.}.

%%>>> art-28

\section*{ARTIGO IV.}
\addcontentsline{toc}{section}{ARTIGO IV.}

\pg{331} \pg{332} \textbf{Modo por que se exprime o conhecimento}

\textbf{intellectual.}

%%>>> t203

\tese{203}{Linguagem.}
\emph{Linguagem} é a manifestação das idéas por meio de palavras --- \emph{Palavra} é um som \emph{articulado}, que o homem emitte pelos orgãos vocaes, com o fim de manifestar uma idéa determinada.

%%>>> t204

\pg{333} \tese{204}{A relação entre a linguagem e as idéas é convencional.}

a) Se a relação entre a linguagem e as idéas não fosse \emph{convencional}, ou \emph{contingente}, mas \emph{natural}, ou \emph{necessaria}, as palavras \pg{334}  não poderiam estar sujeitas a variações, mas deveriam conservar sempre a mesma significação; porque o signal \pg{335}  \emph{natural} é uniforme e constante, como a natureza. Ora não é assim; porque as palavras estão sujeitas a variações, na mesma lingua, como é por todos sabido. Logo a relação entre a linguagem e as idéas é \emph{convencional}.
